% Options for packages loaded elsewhere
\PassOptionsToPackage{unicode}{hyperref}
\PassOptionsToPackage{hyphens}{url}
%
\documentclass[
  10pt,
]{article}
\usepackage{lmodern}
\usepackage{amssymb,amsmath}
\usepackage{ifxetex,ifluatex}
\ifnum 0\ifxetex 1\fi\ifluatex 1\fi=0 % if pdftex
  \usepackage[T1]{fontenc}
  \usepackage[utf8]{inputenc}
  \usepackage{textcomp} % provide euro and other symbols
\else % if luatex or xetex
  \usepackage{unicode-math}
  \defaultfontfeatures{Scale=MatchLowercase}
  \defaultfontfeatures[\rmfamily]{Ligatures=TeX,Scale=1}
\fi
% Use upquote if available, for straight quotes in verbatim environments
\IfFileExists{upquote.sty}{\usepackage{upquote}}{}
\IfFileExists{microtype.sty}{% use microtype if available
  \usepackage[]{microtype}
  \UseMicrotypeSet[protrusion]{basicmath} % disable protrusion for tt fonts
}{}
\makeatletter
\@ifundefined{KOMAClassName}{% if non-KOMA class
  \IfFileExists{parskip.sty}{%
    \usepackage{parskip}
  }{% else
    \setlength{\parindent}{0pt}
    \setlength{\parskip}{6pt plus 2pt minus 1pt}}
}{% if KOMA class
  \KOMAoptions{parskip=half}}
\makeatother
\usepackage{xcolor}
\IfFileExists{xurl.sty}{\usepackage{xurl}}{} % add URL line breaks if available
\IfFileExists{bookmark.sty}{\usepackage{bookmark}}{\usepackage{hyperref}}
\hypersetup{
  pdftitle={Polarized closure of the integral Einstein obstruction},
  hidelinks,
  pdfcreator={LaTeX via pandoc}}
\urlstyle{same} % disable monospaced font for URLs
\usepackage[a4paper,margin=18mm]{geometry}
\setlength{\emergencystretch}{3em} % prevent overfull lines
\providecommand{\tightlist}{%
  \setlength{\itemsep}{0pt}\setlength{\parskip}{0pt}}
\setcounter{secnumdepth}{-\maxdimen} % remove section numbering

\usepackage{esint}
\providecommand{\Ric}{\operatorname{Ric}}
\usepackage{fvextra}
\usepackage{array}
\fvset{breaklines=true,breakanywhere=true}
\usepackage{xeCJK}
\setmainfont{DejaVu Serif}
\setmonofont{DejaVu Sans Mono}
\setCJKmainfont{Noto Serif CJK SC}
\setCJKmonofont{Noto Sans Mono CJK SC}
\setmathfont{Asana Math}
\setmathfont{Latin Modern Math}[range=\lrcorner]
\setlength{\emergencystretch}{3em}
\allowdisplaybreaks
\hypersetup{hypertexnames=false}
\pdfstringdefDisableCommands{\def\nolinkurl#1{#1}}

\title{Polarized closure of the integral Einstein obstruction}
\usepackage{etoolbox}
\makeatletter
\providecommand{\subtitle}[1]{% add subtitle to \maketitle
  \apptocmd{\@title}{\par {\large #1 \par}}{}{}
}
\makeatother
\subtitle{Gauge input, uniform constants, and all Einstein signs}
\author{}
\date{22 September 2026}

\begin{document}
\maketitle

\hypertarget{ux7b2cux4e03ux5377ux539fux59cbux79efux5206ux673aux5236ux53c2ux6570ux7ea6ux5b9aux4e0eux771fux5b9eux89c4ux8303ux4f9bux7ed9}{%
\section{第七卷：原始积分机制、参数约定与真实规范供给}\label{ux7b2cux4e03ux5377ux539fux59cbux79efux5206ux673aux5236ux53c2ux6570ux7ea6ux5b9aux4e0eux771fux5b9eux89c4ux8303ux4f9bux7ed9}}

Research snapshot: 0b1490791be560203b86a7d9834dcdc3cbc83b41.

十一份既有技术前置全文，未改写为摘要，不计为新增。经典输入与项目推导分别归属；旧的二次小h状态已被第五卷线性结果替代。此卷不意味着第一主线所有实现分支都重新打包。

Read \nolinkurl{00_READ_FIRST.md} for superseded statements. Each source below is
complete and byte-preserved. Original local links remain relative to
their source; locate other texts using \nolinkurl{10_source_map_and_history.md}.
Formula numbers restart in each source.

\begin{center}\rule{0.5\linewidth}{0.5pt}\end{center}

\hypertarget{source-1-line1_integral_rigidity_paper.md}{%
\subsection{Source 1:
\nolinkurl{line1_integral_rigidity_paper.md}}\label{source-1-line1_integral_rigidity_paper.md}}

Original path:
\nolinkurl{mathematics/current/line1_integral_rigidity_paper.md}

SHA-256:
\nolinkurl{ca664aa1f07d2907c8bd0b346d4e0b824cb858335fa2a8d15c9cce33ba7e3d54}

BEGIN COMPLETE SOURCE

\hypertarget{integral-einstein-rigidity-from-polarized-monopole-classes}{%
\section{Integral Einstein rigidity from polarized monopole
classes}\label{integral-einstein-rigidity-from-polarized-monopole-classes}}

\hypertarget{abstract}{%
\subsection{Abstract}\label{abstract}}

We prove a scale-invariant, integral obstruction to Einstein metrics on
closed oriented four-manifolds with suitable polarized Seiberg--Witten
classes. The analytic statement uses actual solutions for every positive
weight, but requires no Sobolev, injectivity-radius, diameter, volume,
pointwise-curvature or curvature-derivative bound. Its proof retains the
nonnegative errors in the monopole inequalities, bounds the spinor by a
truncated weight, constructs a globally smooth polynomial Weyl model,
and tests the contracted Bianchi identity against the divergence of that
model. A bundle Hessian estimate proved below closes this self-test
without a geometric compactness argument.

The resulting universal constant gives an integral Einstein gap for the
connected sums occurring in LeBrun's Einstein obstruction. The class may
depend on the metric, including its chamber when the positive Betti
number is one; the same selected class must work for every weight in
that metric's conformal class. Positive and negative Einstein constants
and the separate Ricci-flat case are treated explicitly. The
positive-budget application strengthens the numerical obstruction in
LeBrun's Theorem 3.3 and Corollary 3.4. Negative topological budgets are
handled directly by Chern--Gauss--Bonnet; the zero-budget, zero-surgery
corner is not an Einstein nonexistence statement.

This manuscript makes the inherited integral-rigidity argument
self-contained analytically. It does not assert priority for either that
argument or the classical gauge-theoretic inputs.

\hypertarget{conventions-and-statements}{%
\subsection{1. Conventions and
statements}\label{conventions-and-statements}}

Let \(M\) be a closed, connected, oriented smooth four-manifold, and let

\[
 V=\operatorname{Vol}_g(M),\qquad T=2\chi(M)+3\tau(M),\qquad
 P=\int_M|\operatorname{Ric}^{\circ}|^2.
\]

Integrals and norms use \(g\), unless stated otherwise. Ricci is a
covariant two-tensor; \(W^+\) is the trace-free symmetric curvature
endomorphism of \(\Lambda^+\), with its operator Hilbert--Schmidt norm.
The Laplacian is nonnegative. Write \(x_-:=\max\{-x,0\}\) and

\[
 4\pi^2T=\int_M\left(\frac{s^2}{24}+2|W^+|^2
                          -\frac{|\operatorname{Ric}^{\circ}|^2}{2}\right).
 \tag{1.1}
\]

For a Spin\(^c\) structure with determinant line \(L\), put

\[
 a=c_1(L)\in H^2(M;\mathbb R),\qquad
 a^+=\operatorname{pr}_{H_g^+}a,\qquad p=(a^+)^2.
\]

All cohomological squares mean intersection squares. The quadratic
spinor map is normalized by

\[
 |\sigma(\Phi)|^2=|\Phi|^4/8.
\]

\textbf{Definition 1.1 (the conformal monopole hypothesis).} A pair
\((g,a)\) satisfies this hypothesis if \(p>0\) and one fixed Spin\(^c\)
structure with first Chern class \(a\) has an unperturbed
Seiberg--Witten solution for every metric conformal to \(g\).
Equivalently for our purposes, for every smooth \(f>0\) there is a
solution of

\[
 D_A\Phi=0,\qquad -iF_A^+=f\sigma(\Phi).                 \tag{1.2}
\]

The implication from conformal solvability to (1.2) is proved in §2. The
class is not allowed to change as \(f\) changes. A different class may
be selected for a different initial metric.

Define the scale-invariant integral Einstein error by

\[
 \mathcal E(g):=\inf_{\lambda\ne0}
  \frac{\int_M|\operatorname{Ric}_g-\lambda g|^2}
       {\lambda^2V}.                                      \tag{1.3}
\]

\textbf{Theorem 1.2 (integral rigidity).} There is an absolute constant

\[
                       \varepsilon_0>0                   \tag{1.4}
\]

with the following property. Suppose that for a metric \(g\) there is a
class \(a\) satisfying Definition 1.1 and

\[
                     (a^+)^2\ge A>0,\qquad a^2\ge T.
 \tag{1.5}
\]

Then

\[
 \boxed{\displaystyle
 \mathcal E(g)+\left(\frac TA-\frac23\right)_+
                         \ge\varepsilon_0.}              \tag{1.6}
\]

In particular, the conclusion holds for every metric if such a class can
be selected metricwise, with the same number \(A\). No regularity of the
class selection is needed. The constant is existential and universal;
this paper does not assign it a numerical value.

The main analytic ingredient is a calibration gap. Let

\[
 \begin{aligned}
 \mathcal O&=\operatorname{Orb}\operatorname{diag}
                    (-1/6,1/12,1/12)\subset S^2_0(\mathbb R^3),\\
 F&=\operatorname{dist}_{\mathrm{HS}}(W^+,\mathcal O)^2,\qquad
 E=\operatorname{Ric}+g,\qquad e=s+4,\\
 \alpha&=\sqrt{2/3}|W^+|,\qquad q=-s+3\alpha,\\
 \mathcal C&=\int_M(|E|^2+F),\qquad
 \Delta_+=\int_Mq_+^2-72\pi^2p,\\
 \kappa&=(\mathcal C+\Delta_+)/V,\qquad
 B_a=(-a^2)_+/V.
 \end{aligned}                                              \tag{1.7}
\]

Here \(\mathcal O\) is interpreted in each oriented orthonormal frame of
\(\Lambda^+\); the distance is frame independent.

\textbf{Theorem 1.3 (analytic calibration gap).} Under Definition 1.1,

\[
 \Delta_+\ge0,\qquad (1+B_a)\kappa\ge c_{\mathrm{cal}}>0, \tag{1.8}
\]

where \(c_{\mathrm{cal}}\) is absolute.

We prove this theorem first. Throughout the analytic proof, \(C\)
denotes an absolute positive constant whose value can change from line
to line. Explicit intermediate constants are retained where they clarify
the topological transfer, not to advertise an optimized final constant.

\hypertarget{actual-weighted-solutions-and-the-retained-monopole-identities}{%
\subsection{2. Actual weighted solutions and the retained monopole
identities}\label{actual-weighted-solutions-and-the-retained-monopole-identities}}

\hypertarget{why-the-weighted-equation-has-solutions}{%
\subsubsection{2.1. Why the weighted equation has
solutions}\label{why-the-weighted-equation-has-solutions}}

Let \(\widehat g=e^{2u}g=f^{-2}g\), so \(u=-\log f\). Under the
fiberwise unitary spinor identification \(\beta\), the conformal
formulas are

\[
 \begin{aligned}
 D_A^{\widehat g}(e^{-3u/2}\beta\Phi)
        &=e^{-5u/2}\beta D_A^g\Phi,\\
 \sigma_{\widehat g}(\beta v)&=e^{2u}\sigma_g(v)
                   \quad\hbox{as covariant two-forms}.
 \end{aligned}                                              \tag{2.1}
\]

For an ordinary solution \((A,\widehat\Phi)\) for \(\widehat g\), set

\[
                 \Phi=f^{-3/2}\beta^{-1}\widehat\Phi.
\]

Then \(D_A^g\Phi=0\), and

\[
 \sigma_g(\Phi)=f^{-1}\sigma_{\widehat g}(\widehat\Phi).
\]

The Hodge star on two-forms is conformally invariant in dimension four,
so the ordinary curvature equation becomes (1.2). In particular
\(H_g^+\) does not change during this operation. This is the conformal
weighted monopole construction of
\href{https://arxiv.org/pdf/math/0611450}{LeBrun, Lemma 3.7}; the
displayed formulas specify the normalization used here.

Set

\[
             \rho=|\Phi|^2,\qquad \psi=2\sqrt2\sigma(\Phi).
 \tag{2.2}
\]

Thus \(\psi\) is a smooth self-dual two-form, \(|\psi|=\rho\), and
\(\rho\) itself is smooth, including at spinor zeros. The harmonic
projection of the determinant curvature gives

\[
 \int_M f^2\rho^2=8\|F_A^+\|_2^2\ge32\pi^2p.           \tag{2.3}
\]

Only the squared norm is used, so the equivalent opposite sign
convention for the Chern representative does not change this formula.

\hypertarget{nonnegative-errors}{%
\subsubsection{2.2. Nonnegative errors}\label{nonnegative-errors}}

The spinor Weitzenböck equation and the self-dual Hodge formula are

\[
 4\nabla_A^*\nabla_A\Phi+(s+f\rho)\Phi=0,                \tag{2.4}
\]

\[
 \|d\psi\|_2^2+\|d^*\psi\|_2^2
   =\int_M\left(|\nabla\psi|^2+\frac{s}{3}\rho^2
                                  -2\langle W^+\psi,\psi\rangle\right).
 \tag{2.5}
\]

The differential of the Hopf quadratic map in (2.2) satisfies

\[
 |d\psi_z(v)|^2
       =4|z|^2|v|^2-4\bigl(\operatorname{Im}\langle z,v\rangle\bigr)^2
       \le4|z|^2|v|^2.                                   \tag{2.6}
\]

The same identity applies covariantly because the map is equivariant and
parallel. Pairing (2.4) with \(\rho\Phi\) and integrating yields

\[
 2\int_M|d\rho|^2+4\int_M\rho|\nabla_A\Phi|^2
                  +\int_Ms\rho^2+\int_M f\rho^3=0.       \tag{2.7}
\]

Let \(w\) be the smallest eigenvalue of \(W^+\). For a trace-free
symmetric three-by-three matrix, \(w\ge-\alpha\). Define

\[
 \begin{aligned}
 D_{\mathrm{amp}}&=2\int_M|d\rho|^2,\\
 D_{\mathrm{quad}}&=\int_M(4\rho|\nabla_A\Phi|^2-|\nabla\psi|^2),\\
 D_H&=\|d\psi\|_2^2+\|d^*\psi\|_2^2=2\|d\psi\|_2^2,\\
 D_{\mathrm{eig}}&=2\int_M
           (\langle W^+\psi,\psi\rangle+\alpha\rho^2),\\
 D_0&=D_{\mathrm{amp}}+D_{\mathrm{quad}}+D_H+D_{\mathrm{eig}}.
 \end{aligned}                                              \tag{2.8}
\]

Each term is nonnegative. Equations (2.5)--(2.7) give the exact identity

\[
 D_0=\int_M\left(\frac23q\rho^2-f\rho^3\right).
 \tag{2.9}
\]

For \(v=3\rho/2\), this is

\[
                  \int_M(qv^2-fv^3)=\frac{27}{8}D_0.
 \tag{2.10}
\]

\hypertarget{signed-spectral-and-scalar-estimates}{%
\subsubsection{2.3. Signed spectral and scalar
estimates}\label{signed-spectral-and-scalar-estimates}}

Put

\[
                  b=2s/3+2w,\qquad h=b_-.
\]

Using \(w\), rather than \(-\alpha\), in (2.5)--(2.7) gives

\[
                       \int_M f\rho^3\le\int_M h\rho^2.
 \tag{2.11}
\]

Choose positive smooth functions \(h_j\ge h\) converging uniformly to
\(h\), and use the actual solutions with \(f_j=\sqrt{h_j}\). Set

\[
 C_j=\int_M f_j^2\rho^2,\quad Z_j=\int_M f_j\rho^3,
 \quad Q_j=\int_M f_j^4.
\]

Equation (2.11) and Hölder imply

\[
 32\pi^2p\le C_j\le Q_j^{1/3}Z_j^{2/3}
                            \le Q_j^{1/3}C_j^{2/3}.
\]

Consequently \(32\pi^2p\le Q_j\), and uniform approximation proves

\[
 \int_M(2s/3+2w)_-^2\ge32\pi^2p,\qquad
 \int_M q_+^2\ge72\pi^2p.                               \tag{2.12}
\]

The second follows from \(b_-\le(2/3)q_+\). This argument passes to a
limit of inequalities, not to a limiting spinor or connection.

For \(f=1\), pairing (2.4) with \(\Phi\) gives

\[
 \int_M\rho^2\le\int_Ms_-\rho
        \le\|s_-\|_2\|\rho\|_2.
\]

Together with (2.3), this proves

\[
                        \int_Ms_-^2\ge32\pi^2p.          \tag{2.13}
\]

In particular \(p>0\) excludes a scalar-flat metric. These are the
classical monopole curvature estimates in the normalization needed here;
the retained terms in (2.8) will supply their quantitative use.

\hypertarget{a-bounded-spinor-with-small-retained-errors}{%
\subsection{3. A bounded spinor with small retained
errors}\label{a-bounded-spinor-with-small-retained-errors}}

The weights used in §2.3 need not bound the spinor uniformly. For the
calibration argument we instead set

\[
       Q=\max\{q,3\},\qquad
       f_0=\max\{\sqrt Q,s_-/4\}.                        \tag{3.1}
\]

Choose smooth majorants \(f_j\ge f_0\) with \(f_j\to f_0\) uniformly,
and choose an actual solution for each \(f_j\). At a positive maximum of
\(\rho\), (2.4) implies \(s+f_j\rho\le0\). Hence

\[
                           0\le\rho\le4,\qquad 0\le v\le6.
 \tag{3.2}
\]

All errors \(o_j(1)\) below tend to zero for this fixed smooth metric
and fixed class. They require only uniform approximation of \(f_0\), not
bounds on derivatives of \(f_j\). No assertion of convergence of the
solutions is made.

The two truncations cost only scalar error:

\[
 0\le Q^2-q_+^2\le9e^2,\qquad
 0\le f_0^2-Q\le\frac9{16}e^2.                          \tag{3.3}
\]

Indeed, the first difference vanishes unless \(q<3\), which implies
\(s>-3\) and \(e>1\). For the second, a nonzero difference implies
\(s_->4\sqrt3>6\), and then \(s_-\le3|s+4|\).

Define ideal moments using \(Q\), not the full truncated weight:

\[
 C_*:=\int_M Qv^2,\qquad Z_*:=\int_M\sqrt Qv^3,
 \qquad E_Q:=\int_M Q^2,\qquad P_*:=C_*-Z_*.
 \tag{3.4}
\]

Equation (2.10), \(Q\ge q\), and \(f_j\ge\sqrt Q\) give

\[
 P_*\ge\frac{27}{8}D_0\ge0.                            \tag{3.5}
\]

The actual curvature moment is

\[
 C_{\mathrm{act}}:=\int_M f_j^2v^2
                      =18\|F_A^+\|_2^2\ge72\pi^2p.
\]

By (3.2)--(3.3),

\[
 0\le C_{\mathrm{act}}-C_*
                     \le\frac{81}{4}\int_M e^2+o_j(1).
\]

Hölder gives \(C_*\le E_Q^{1/3}Z_*^{2/3}\). Since \(Z_*\le C_*\), it
follows that \(C_*\le E_Q\). Thus, with

\[
                        S_*:=\Delta_++\frac{117}{4}\int_Me^2,
\]

we have

\[
                     0\le E_Q-C_*\le S_*+o_j(1).         \tag{3.6}
\]

For completeness, two elementary moment inequalities recover all the
errors needed below. Write

\[
 \theta=(Z_*/E_Q)^{1/3},\qquad
 H_*=E_Q^{1/3}Z_*^{2/3}-C_*.
\]

Here \(0<\theta\le1\), \(H_*\ge0\), and direct algebra gives

\[
 E_Q-C_*=H_*+\frac{1+\theta}{\theta^2}(P_*+H_*)\ge2P_*.
 \tag{3.7}
\]

For the probability measure \(d\nu=Q^2d\mu/E_Q\), let \(r=v/\sqrt Q\),
\(c=\int r^2d\nu\), \(z=\int r^3d\nu\), and \(m=\int rd\nu\). We have
\(z\le c\le z^{2/3}\), hence \(c\le1\). Also \(c^2\le mz\le mc\), so
\(m\ge c\). Therefore

\[
 \int_M Q(v-\sqrt Q)^2=E_Q(c-2m+1)\le E_Q-C_*.
 \tag{3.8}
\]

Equations (3.5)--(3.8) prove

\[
 \begin{aligned}
 D_0&\le\frac4{27}S_*+o_j(1),\\
 \int_M|d\rho|^2,\ \int_M|d\psi|^2
     &\le\frac2{27}S_*+o_j(1),\\
 \int_MQ(v-\sqrt Q)^2&\le S_*+o_j(1).
 \end{aligned}                                              \tag{3.9}
\]

The point of the ideal moments is that these conclusions do not require
any estimate of \(\int f_j^4\), which could be unnecessarily large.

\hypertarget{a-globally-smooth-approximate-weyl-tensor}{%
\subsection{4. A globally smooth approximate Weyl
tensor}\label{a-globally-smooth-approximate-weyl-tensor}}

The norm of an element of \(\mathcal O\) is \(1/(2\sqrt6)\).
Consequently

\[
 \left\|q-\frac92\right\|_2^2
       \le2\|e\|_2^2+12\int_MF\le12\mathcal C.           \tag{4.1}
\]

Clipping at \(3\) moves \(q\) toward \(9/2\), so the same bound holds
with \(Q\). Since \(Q\ge3\), (3.8) implies

\[
 \|\rho-\sqrt2\|_2^2
  \le\frac8{27}S_*+\frac2{27}\left\|Q-\frac92\right\|_2^2+o_j(1).
 \tag{4.2}
\]

Here the first term follows by writing
\(\rho-\sqrt2=(2/3)(v-\sqrt Q)+(2/3)(\sqrt Q-\sqrt{9/2})\) and using
\(Q\ge3\).

To estimate direction, choose a measurable rank-one orthogonal
projection \(\Pi\) onto a lowest eigenspace of \(W^+\). Then

\[
                  A_0=I/12-\Pi/4
\]

is a nearest point of \(\mathcal O\). On the good set \(F\le1/256\),
comparison of ordered eigenvalues with \(A_0\) gives

\[
                   \lambda_2(W^+)-\lambda_1(W^+)\ge1/8.
\]

Putting \(y=(I-\Pi)\psi\), it follows from the eigenvalue error in (2.8)
that

\[
 \int_{\{F\le1/256\}}|y|^2\le\frac{16}{27}S_*+o_j(1).
\]

The complementary set has volume at most \(256\int F\), while
\(|\psi|^2\le16\). Thus

\[
                 \int_M|y|^2\le\frac{16}{27}S_*+4096\int_MF+o_j(1).
 \tag{4.3}
\]

No derivative of \(\Pi\) will be taken. Instead define the polynomial
endomorphism

\[
       A_\psi=\frac{\rho^2}{24}I-\frac18\psi\otimes\psi,
             \qquad R=W^+-A_\psi.                       \tag{4.4}
\]

It is globally smooth, trace-free, and defined even at spinor zeros. Its
eigenvalues are \(-\rho^2/12,\rho^2/24,\rho^2/24\), and
\(|A_\psi|^2=\rho^4/96\). The rank-one identity

\[
 |\psi\otimes\psi-\rho^2\Pi|^2=2\rho^2|y|^2
\]

gives the useful comparison

\[
 |A_\psi-A_0|^2
       \le\frac{\rho^2}{16}|y|^2+\frac{(\rho^2-2)^2}{48}.
 \tag{4.5}
\]

For example, (4.5) follows by inserting the intermediate matrix
\(\rho^2 A_0/2\) and applying \(|X+Y|^2\le2|X|^2+2|Y|^2\). Since
\(\rho\le4\) and \((4+\sqrt2)^2<36\), (4.1)--(4.5) imply

\[
 \begin{aligned}
 \|R\|_2^2
 &\le\frac{44}{27}S_*+8194\int_MF
                     +\frac19\left\|Q-\frac92\right\|_2^2+o_j(1)\\
 &\le2\Delta_++8500\mathcal C+o_j(1).
 \end{aligned}                                              \tag{4.6}
\]

In the final line we used \(\int e^2\le4\int|E|^2\). In particular, each
of

\[
 \|e\|_2^2,\quad\|\rho-\sqrt2\|_2^2,\quad\|R\|_2^2,
            \quad\|d\rho\|_2^2,\quad\|d\psi\|_2^2
\]

is bounded by \(C\kappa V+o_j(1)\).

\hypertarget{a-hessian-estimate-that-needs-no-sobolev-constant}{%
\subsection{5. A Hessian estimate that needs no Sobolev
constant}\label{a-hessian-estimate-that-needs-no-sobolev-constant}}

\textbf{Lemma 5.1.} Let \(u\) be a smooth section of a metric or
Hermitian vector bundle with compatible connection over a closed
Riemannian four-manifold. Let \(\mathscr F\) be its curvature, with
component norm \(\sum_{i,j}|\mathscr F_{ij}|_{\mathrm{HS}}^2\), and set
\(K=\|u\|_\infty\). Then

\[
 \begin{aligned}
 \|\nabla u\|_4^2&\le4K\|\nabla^2u\|_2,\\
 \|\nabla^2u\|_2^2
 &\le2\|\nabla^*\nabla u\|_2^2
       +32K^2\|\operatorname{Ric}\|_2^2
       +33K^2\|\mathscr F\|_2^2.
 \end{aligned}                                              \tag{5.1}
\]

\textbf{Proof.} Take real parts of Hermitian inner products when
necessary. Write \(v_i=\nabla_i u\), \(H_{ij}=\nabla_i v_j\), and
\(L=\sum_iH_{ii}=-\nabla^*\nabla u\). Integrating the divergence of
\(X_i=|v|^2\langle u,v_i\rangle\) gives

\[
 \int|v|^4
 \le K\int |v|^2(2|H|+|L|)
 \le K\|v\|_4^2(2\|H\|_2+\|L\|_2).
\]

Because \(|L|\le2|H|\) in dimension four, this proves the first
inequality, including the case \(\|v\|_4=0\).

For the second, antisymmetry gives

\[
 H_{ij}-H_{ji}=\mathscr F_{ij}u,\qquad
 \|H\|_2^2=\int\langle H_{ij},H_{ji}\rangle
                         +\frac12\sum_{i,j}\|\mathscr F_{ij}u\|_2^2.
 \tag{5.2}
\]

The commutator on the bundle-valued one-form \(v\) yields

\[
 \sum_i\nabla_i\nabla_jv_i
     =\nabla_jL+\operatorname{Ric}_{jk}v_k
                                      +\sum_i\mathscr F_{ij}v_i.
\]

Integrating by parts in the first term of (5.2) and using the
skew-adjointness of \(\mathscr F_{ij}\), we obtain the exact identity

\[
 \begin{aligned}
 \|H\|_2^2={}&\|L\|_2^2
  -\int\operatorname{Ric}_{ij}\langle v_i,v_j\rangle
  +\sum_{i,j}\int\langle\mathscr F_{ij}v_j,v_i\rangle
  +\frac12\sum_{i,j}\|\mathscr F_{ij}u\|_2^2.
 \end{aligned}                                              \tag{5.3}
\]

In particular no derivative of curvature appears. Hölder and the first
part bound the middle terms in absolute value by

\[
   4K(\|\operatorname{Ric}\|_2+\|\mathscr F\|_2)\|H\|_2.
\]

Use Young's inequality to absorb half of \(\|H\|_2^2\), and then
\((x+y)^2\le2(x^2+y^2)\). This gives exactly the second inequality in
(5.1). No Sobolev inequality or curvature smallness has been used.
\(\square\)

\hypertarget{applying-the-lemma-to-the-bounded-spinor}{%
\subsubsection{5.2. Applying the lemma to the bounded
spinor}\label{applying-the-lemma-to-the-bounded-spinor}}

The calibration budget gives

\[
 \begin{aligned}
 \int|W^+|^2&\le2\int F+V/12,\\
 \int|\operatorname{Ric}|^2&\le2\int|E|^2+8V,\\
 \int s^2&\le8\int|E|^2+32V.
 \end{aligned}                                              \tag{5.4}
\]

Equations (3.3) and (4.1) also imply

\[
                         \int f_j^2\le C(V+\mathcal C)+o_j(1).
 \tag{5.5}
\]

Indeed \(Q\le9/2+|Q-9/2|\); Cauchy--Schwarz and
\(\sqrt{V\mathcal C}\le(V+\mathcal C)/2\) suffice.

For the determinant connection, Chern--Weil and (3.2) give

\[
 \|F_A\|_2^2
   =2\|F_A^+\|_2^2-4\pi^2a^2
   \le4\int f_j^2+4\pi^2(-a^2)_+.
 \tag{5.6}
\]

The positive-spinor connection curvature \(\mathscr F_A\) is a fixed
linear representation of the Riemannian and determinant curvatures. In
the block decomposition of the Riemannian curvature operator, its
nontrivial spin action uses the \(\Lambda^+\) output: it sees \(W^+\),
\(s\), and the trace-free Ricci block, but not \(W^-\). Consequently
(5.4)--(5.6) imply

\[
             \int|\mathscr F_A|^2
               \le C\bigl(V+\mathcal C+(-a^2)_+\bigr)+o_j(1).
 \tag{5.7}
\]

From (2.4) and \(\rho\le4\),

\[
 \|\nabla_A^*\nabla_A\Phi\|_2^2
       \le\frac12\int s^2+8\int f_j^2
       \le C(V+\mathcal C)+o_j(1).                        \tag{5.8}
\]

Apply Lemma 5.1 with \(\|\Phi\|_\infty\le2\). Since \(A_\psi\) is a
parallel polynomial of degree four in \(\Phi\),

\[
 |\nabla^2A_\psi|
      \le C\bigl(|\Phi|^3|\nabla_A^2\Phi|
                               +|\Phi|^2|\nabla_A\Phi|^2\bigr).
\]

We conclude that

\[
 \int\bigl(|\nabla_A^2\Phi|^2+|\nabla_A\Phi|^4
                                      +|\nabla^2A_\psi|^2\bigr)
       \le C\bigl(V+\mathcal C+(-a^2)_+\bigr)+o_j(1).
 \tag{5.9}
\]

This proves the required second-derivative budget without
differentiating the weight and without an estimate of the anti-self-dual
Weyl curvature.

\hypertarget{a-weak-bianchi-self-test-and-its-incompatible-lower-bound}{%
\subsection{6. A weak Bianchi self-test and its incompatible lower
bound}\label{a-weak-bianchi-self-test-and-its-incompatible-lower-bound}}

\hypertarget{tensor-normalization-and-the-bianchi-upper-bound}{%
\subsubsection{6.1. Tensor normalization and the Bianchi upper
bound}\label{tensor-normalization-and-the-bianchi-upper-bound}}

For an orthonormal basis \(\epsilon_a\) of \(\Lambda^+\), put
\(\omega_a=\sqrt2\epsilon_a\). Associate to a symmetric endomorphism
\(A\) the four-tensor

\[
 W(A)_{ijkl}=\frac12\sum_{a,b}A_{ab}(\omega_a)_{ij}(\omega_b)_{kl}.
 \tag{6.1}
\]

Let

\[
                    J_\psi=\operatorname{div}W(A_\psi),
\]

using the full three-index norm \(\sum_{j,k,l}|J_{jkl}|^2\). Changing
the overall sign of divergence changes none of the following squared
norms.

The contracted Bianchi identity expresses \(\operatorname{div}W^+\) as a
parallel algebraic contraction of \(\nabla\operatorname{Ric}^\circ\) and
\(\nabla s\). This follows by contracting the Riemannian Bianchi
identity, subtracting the Schouten part, and projecting the final
two-form pair onto \(\Lambda^+\); that projection is parallel. Since
\(\nabla\operatorname{Ric}=\nabla E\) and
\(\nabla s=\nabla\operatorname{tr}E\), integration by parts and
\(W(A_\psi)=W^+-W(R)\) show, for every smooth test tensor \(Z\), that

\[
 \left|\int\langle J_\psi,Z\rangle\right|
       \le C(\|E\|_2+\|R\|_2)\|\nabla Z\|_2.
 \tag{6.2}
\]

Set \(Z=J_\psi\). It is smooth, and
\(\|\nabla J_\psi\|_2\le C\|\nabla^2A_\psi\|_2\). Equations (4.6) and
(5.9) therefore give

\[
 \frac1V\int|J_\psi|^2
            \le C\sqrt\kappa\sqrt{1+\kappa+B_a}+o_j(1).
 \tag{6.3}
\]

This is a weak identity tested against a smooth tensor already supplied
by the spinor. It asserts no pointwise bound on
\(\nabla\operatorname{Ric}\).

\hypertarget{the-exact-divergence-identity}{%
\subsubsection{6.2. The exact divergence
identity}\label{the-exact-divergence-identity}}

Define

\[
 B_\psi=\frac{s}{3}\rho^2-2\langle W^+\psi,\psi\rangle,
 \qquad C_\psi=\langle\iota_{\nabla\rho}\psi,*d\psi\rangle.
\]

We now derive the identity used to lower-bound the left side of (6.3).
Pairing the self-dual Hodge--Weitzenböck equation with \(\rho^2\psi\)
and integrating by parts gives

\[
 \int\rho^2\bigl(|\nabla\psi|^2+B_\psi-2|d\psi|^2\bigr)
          =4\int\rho C_\psi-2\int\rho^2|d\rho|^2.
 \tag{6.4}
\]

To see the mixed terms explicitly, use \(d^*\psi=-*d\psi\),
\(d^*(\rho^2\psi)=\rho^2d^*\psi-2\rho\iota_{\nabla\rho}\psi\), and
\(d(\rho^2\psi)=\rho^2d\psi+2\rho d\rho\wedge\psi\). The two
exterior-derivative mixed terms add to \(4\rho C_\psi\); the
rough-Laplacian mixed term is \(2\rho^2|d\rho|^2\).

The required pointwise first-jet identities are

\[
 |\nabla A_\psi|^2
       =\frac{\rho^2}{32}|\nabla\psi|^2
                                +\frac{\rho^2}{96}|d\rho|^2,
 \tag{6.5}
\]

\[
 |\nabla A_\psi|^2-|J_\psi|^2
    =\frac{\rho^2}{64}(|\nabla\psi|^2-2|d\psi|^2)
      +\frac{\rho^2}{32}|d\rho|^2-\frac\rho{16}C_\psi.
 \tag{6.6}
\]

Here is a direct component verification that fixes all normalization
factors. At a point with \(\rho>0\), choose a parallel orthonormal frame
so that

\[
 \begin{aligned}
 \epsilon_1&=(dx^0\wedge dx^1+dx^2\wedge dx^3)/\sqrt2,\\
 \epsilon_2&=(dx^0\wedge dx^2+dx^3\wedge dx^1)/\sqrt2,\\
 \epsilon_3&=(dx^0\wedge dx^3+dx^1\wedge dx^2)/\sqrt2,\\
 \psi&=\rho\epsilon_1,\qquad
 \nabla_m\psi=c_m\epsilon_1+a_m\epsilon_2+b_m\epsilon_3.
 \end{aligned}
\]

Then \(c_m=\partial_m\rho\), and differentiating (4.4) gives

\[
 \nabla_m A_\psi=
 \begin{pmatrix}
  -\rho c_m/6&-\rho a_m/8&-\rho b_m/8\\
  -\rho a_m/8&\rho c_m/12&0\\
  -\rho b_m/8&0&\rho c_m/12
 \end{pmatrix}.                                         \tag{6.7}
\]

The squared matrix norm gives (6.5). For (6.6), insert (6.7) in
\(J_{jkl}=\sum_{m,a,b}(\nabla_mA)_{ab}(\epsilon_a)_{mj}(\epsilon_b)_{kl}\),
use \(\sum_{k,l}(\epsilon_b)_{kl}(\epsilon_c)_{kl}=2\delta_{bc}\), and
compute
\(d\psi_{ijk}=\nabla_i\psi_{jk}+\nabla_j\psi_{ki} +\nabla_k\psi_{ij}\)
in the displayed basis. Explicitly,

\[
 *d\psi=\frac1{\sqrt2}
 \begin{pmatrix}
  -a_2-b_3-c_1\\
   a_3-b_2+c_0\\
   a_0+b_1-c_3\\
  -a_1+b_0+c_2
 \end{pmatrix},
 \qquad
 \iota_{\nabla\rho}\psi
       =\frac{\rho}{\sqrt2}(-c_1,c_0,-c_3,c_2).
\]

These expressions and (6.7) give (6.6) by summing squares and collecting
the mixed terms in \(C_\psi\). Thus (6.5)--(6.6) are algebraic
identities in the twelve arbitrary numbers \(c_m,a_m,b_m\), with no
field equation imposed. At a spinor zero, both \(\psi\) and its first
derivative vanish, and \(d\rho=0\); the identities extend there as well.

Combining (6.4)--(6.6) eliminates \(\int\rho^2|\nabla\psi|^2\) and gives

\[
 \begin{aligned}
 \int|J_\psi|^2={}&-\frac1{64}\int\rho^2 B_\psi
       +\frac1{16}\int\rho^2|d\psi|^2
       -\frac5{96}\int\rho^2|d\rho|^2
       +\frac18\int\rho C_\psi.
 \end{aligned}                                              \tag{6.8}
\]

\hypertarget{the-positive-lower-bound-and-contradiction}{%
\subsubsection{6.3. The positive lower bound and
contradiction}\label{the-positive-lower-bound-and-contradiction}}

Because \(A_\psi\psi=-\rho^2\psi/12\), the curvature term in (6.8) is

\[
 -\frac{\rho^2B_\psi}{64}
        =-\frac{s\rho^4}{192}-\frac{\rho^6}{384}
                                +\frac{\rho^2\langle R\psi,\psi\rangle}{32}.
 \tag{6.9}
\]

At \(s=-4\), \(\rho=\sqrt2\), and \(R=0\), it equals \(1/16\). On the
interval \(0\le\rho\le4\), the error in (6.9) is bounded by

\[
                  C\bigl(|e|+|\rho-\sqrt2|+|R|\bigr).
\]

Also \(|C_\psi|\le\rho|d\rho||d\psi|\). Discard the nonnegative
\(|d\psi|^2\) term in (6.8), apply (3.9), (4.2), and (4.6), and divide
by \(V\). Cauchy--Schwarz gives

\[
 \frac1V\int|J_\psi|^2
                   \ge\frac1{16}-C(\sqrt\kappa+\kappa)+o_j(1).
 \tag{6.10}
\]

For each fixed metric let the smooth-weight approximation errors tend to
zero in (6.3) and (6.10). The two numerical inequalities remain valid
without choosing any limit of the spinors. For a sufficiently small
absolute \(\kappa_0\le1\), the lower bound is at least \(1/32\) when
\(\kappa\le\kappa_0\). The upper bound then forces

\[
                         \kappa(1+\kappa+B_a)\ge c_1>0.
\]

As \(1+\kappa+B_a\le2(1+B_a)\) in this range, (1.8) follows. For
\(\kappa>\kappa_0\), it follows after decreasing the same universal
constant. This proves Theorem 1.3. \(\square\)

\hypertarget{from-the-calibration-gap-to-topological-integral-rigidity}{%
\subsection{7. From the calibration gap to topological integral
rigidity}\label{from-the-calibration-gap-to-topological-integral-rigidity}}

We now prove Theorem 1.2, first at the negative normalization. Assume
(1.5), and put

\[
 r=\frac1V\int|\operatorname{Ric}+g|^2,\qquad
 d=\left(\frac TA-\frac23\right)_+.
\]

Continue to use \(\alpha=\sqrt{2/3}|W^+|\) and
\(w=\lambda_{\min}(W^+)\), and introduce

\[
 t=s+24\alpha=s+8\sqrt6|W^+|,\qquad
 b=2s/3+2w,\qquad
 G=(2s/3-2\alpha)^2-b_-^2.                               \tag{7.1}
\]

Since \(b\ge2s/3-2\alpha\), we have \(G\ge0\), including where the
scalar curvature is positive. Direct expansion gives

\[
 t^2+18G=216\left(s^2/24+2|W^+|^2\right)-18b_-^2.
\]

Apply (1.1) and (2.12) to obtain

\[
                    \int(t^2+18G)\le108P+864\pi^2Ad.
 \tag{7.2}
\]

Here \(p\ge A\) and \(T-2p/3\le T-2A/3\le Ad\). Since
\(\int|E|^2=P+\frac14\int e^2\), (7.2) implies

\[
             \int(|E|^2+e^2+t^2+G)\le109rV+864\pi^2Ad.
 \tag{7.3}
\]

We record the pointwise orbit comparison, including its scalar-positive
branch. Set \(z=w+\alpha\); trace-freeness gives \(0\le z\le\alpha\),
and minimizing over the orbit gives exactly

\[
 F=\frac32(\alpha-1/6)^2+z/2,\qquad
                         \alpha=(t-e+4)/24.              \tag{7.4}
\]

If \(s\le-3\), then \(b<0\) and \(G=z(-8s/3+8\alpha-4z)\ge8z\). Thus

\[
                       F\le(t^2+e^2)/192+G/16.
\]

If \(s>-3\), then \(e>1\). Use \(z\le\alpha\), \(4-e\le3e^2\), and
\(t\le(t^2+1)/2\) in (7.4) to get

\[
                       F\le(t^2+5e^2)/64.
\]

In both cases,

\[
                             F\le e^2+t^2+G.             \tag{7.5}
\]

The second completed square is

\[
 27\left(s^2/24+2|W^+|^2\right)-q^2=t^2/8.
\]

It and (1.1) yield

\[
 \Delta_+\le\frac{27}{2}rV+108\pi^2Ad.
\]

Combining this with (7.3)--(7.5),

\[
                     \kappa\le\frac{245}{2}r
                                      +972\pi^2\frac AVd.
 \tag{7.6}
\]

The scalar estimate (2.13) gives

\[
 32\pi^2A\le\int s^2\le32V+8rV,
              \qquad \frac AV\le\frac{4+r}{4\pi^2}.
 \tag{7.7}
\]

For \(r\le1\), therefore,

\[
                              \kappa\le1215(r+d).        \tag{7.8}
\]

When \(r\le1/4\), reverse triangle inequality in \(L^2\) and (1.1) also
imply

\[
 4\pi^2T
     \ge\left[\frac{(4-2\sqrt r)^2}{24}-\frac r2\right]V
     \ge V/4>0.                                          \tag{7.9}
\]

The square condition \(a^2\ge T\) then gives \(B_a=0\). Apply Theorem
1.3 and (7.8). If \(r>1/4\), no argument is needed. In either case

\[
       r+d\ge\varepsilon_0,
       \qquad \varepsilon_0:=\min\{1/4,c_{\mathrm{cal}}/1215\}.
 \tag{7.10}
\]

For reference, a weaker square hypothesis has a correspondingly
nonuniversal version. If \(a^2\ge-C_AA\), \(C_A\ge0\), the same proof
gives

\[
 r+d\ge\min\left\{1,
          \frac{c_{\mathrm{cal}}}{1215(1+5C_A/(4\pi^2))}\right\}.
 \tag{7.11}
\]

If instead \(a^2\ge T-C_AA\), replace the first entry of that minimum by
\(1/4\), using (7.9). An arbitrary monopole class must not silently be
assumed to have \(a^2\ge T\).

\hypertarget{positive-scale-arbitrary-scale-and-the-ricci-flat-case}{%
\subsection{8. Positive scale, arbitrary scale, and the Ricci-flat
case}\label{positive-scale-arbitrary-scale-and-the-ricci-flat-case}}

\hypertarget{the-positive-normalization}{%
\subsubsection{8.1. The positive
normalization}\label{the-positive-normalization}}

Put \(r_+=V^{-1}\int|\operatorname{Ric}-g|^2\). Since
\(s_-^2\le(s-4)^2\le4|\operatorname{Ric}-g|^2\), (2.13) gives

\[
                            A/V\le r_+/(8\pi^2).
\]

If \(r_+\le1/4\), the analogue of (7.9) gives \(4\pi^2T/V\ge1/4\). If
also \(d\le1/4\), then \(T/A\le2/3+d\), so

\[
  \frac14\le\frac{4\pi^2T}{V}
             \le\frac{r_+}{2}\left(\frac23+d\right)
             \le\frac{11}{96},
\]

a contradiction. In particular \(r_++d>1/4\).

\hypertarget{scale-optimization}{%
\subsubsection{8.2. Scale optimization}\label{scale-optimization}}

For \(\lambda\ne0\), apply the appropriate normalized result to
\(|\lambda|g\). The positive harmonic subspace and the conformal
solvability hypothesis are unchanged. The normalized error becomes

\[
 \frac{\int|\operatorname{Ric}_g-\lambda g|^2}{\lambda^2V}.
\]

Taking the infimum proves (1.6). This completes Theorem 1.2. \(\square\)

For interpretation, write

\[
 R_2=\int|\operatorname{Ric}|^2,\qquad S=\int s,\qquad
 L_s=\int(s-S/V)^2.
\]

When \(R_2>0\), direct minimization in \(1/\lambda\) gives

\[
              \mathcal E(g)=4-\frac{S^2}{VR_2}
                             =\frac{4P+L_s}{R_2}.        \tag{8.1}
\]

This formula includes \(S=0\), when the infimum is \(4\) and need not be
attained. If \(R_2=0\), definition (1.3) gives \(\mathcal E(g)=4\), not
zero. Thus a lower bound for \(\mathcal E\) alone does not exclude a
Ricci-flat metric. Ricci-flat metrics in Theorem 1.2 are excluded
separately by (2.13), since \(p\ge A>0\).

For an Einstein metric with nonzero Einstein constant, \(\mathcal E=0\).
A pointwise operator error also dominates the integral one: for
\(\sigma=\pm1\),

\[
 \delta_\sigma(g):=\|\operatorname{Ric}^{\sharp}-\sigma I\|_{L^\infty,\mathrm{op}}
      \quad\Longrightarrow\quad\mathcal E(g)\le4\delta_\sigma(g)^2.
 \tag{8.2}
\]

There are exactly four Ricci eigenvalues, which explains this factor.

\hypertarget{metricwise-polarized-classes-and-chamber-bookkeeping}{%
\subsection{9. Metricwise polarized classes and chamber
bookkeeping}\label{metricwise-polarized-classes-and-chamber-bookkeeping}}

The analytic theorem needs the quantifiers

\[
 \boxed{\begin{gathered}
   \text{for every }g\text{ there exists a Spin}^c
       \text{ structure }\mathfrak s_g\\
   \text{such that for every }f>0,\text{ equation (1.2)}\\
   \text{has a solution for the same }\mathfrak s_g.
 \end{gathered}}
 \tag{9.1}
\]

It does not require one universal choice of signs for all metrics. This
distinction is essential when \(b_+=1\).

\hypertarget{the-connected-sum-input-with-insertions}{%
\subsubsection{9.1. The connected-sum input with
insertions}\label{the-connected-sum-input-with-insertions}}

Let \(X\) have \(b_+(X)>0\), and suppose a Spin\(^c\) structure
\(\mathfrak s\), with real first Chern class \(c\), has a nonzero
Seiberg--Witten invariant evaluated on a homogeneous insertion \(z\).
The invariant is defined on

\[
 \mathbb A(X)=\mathbb Z[U]\otimes\Lambda^*H_1(X;\mathbb Z)/\mathrm{torsion},
                 \qquad\deg U=2,\quad\deg H_1=1.
\]

Its expected dimension is \((c^2-T_X)/4\), where
\(T_X=2\chi(X)+3\tau(X)\). Assume

\[
                    T_X>0,\qquad A:=c^2\ge T_X.          \tag{9.2}
\]

The second inequality follows automatically from the nonnegative degree
of a nonzero ordinary invariant, but is displayed because it is used in
the proof. For \(b_+=1\), the nonzero invariant here must be in the
small-perturbation chamber at the relevant polarization, not merely in
some unspecified chamber.

Let

\[
 M=X\#k\,\overline{\mathbb{CP}}^{2}\#\ell(S^1\times S^3),
 \qquad D=k+4\ell,\qquad T_M=T_X-D.                      \tag{9.3}
\]

For a prescribed positive polarization \(H=H_g^+\), choose the signs of
the exceptional generators \(E_j\) so that

\[
                         c_H^+\mathbin\cdot E_j\le0,
\]

and set

\[
                             a_H=c-\sum_{j=1}^kE_j.
 \tag{9.4}
\]

Then

\[
 \begin{aligned}
 (a_H^+)^2
   &=(c_H^+)^2-2\sum_jc_H^+\mathbin\cdot E_j
                              +\left|\sum_j(E_j)_H^+\right|^2
      \ge(c_H^+)^2\ge c^2=A,\\
 a_H^2&=A-k\ge T_X-k-4\ell=T_M.
 \end{aligned}                                              \tag{9.5}
\]

The ordinary blow-up formula with exceptional coefficient \(\pm1\)
preserves expected dimension and preserves the original insertion \(z\).
The \(S^1\times S^3\) operation raises expected dimension by one and
adds the degree-one insertion of its new circle. After \(\ell\)
operations the nonzero invariant is therefore evaluated on

\[
                           z\,\mu(C_1)\cdots\mu(C_\ell).
 \tag{9.6}
\]

These statements include positive-dimensional base invariants. They are
the special cases of the blow-up and surgery formulas in
\href{https://arxiv.org/pdf/math/0005268}{Ozsváth--Szabó, Theorem 3.2,
Proposition 2.2 and Remark 2.3}. The last remark supplies the
chamber-compatible version when \(b_+=1\). Using only a zero-dimensional
numerical invariant and omitting the insertions in (9.6) would not
justify the general statement.

\hypertarget{what-the-chamber-condition-means}{%
\subsubsection{9.2. What the chamber condition
means}\label{what-the-chamber-condition-means}}

For \(b_+(X)=1\), choose a unit forward generator \(\omega_H\) of the
positive line on \(M\) so that \(c\cdot\omega_H>0\). The \(X\) component
\(\omega_X\) has positive square, lies in the corresponding forward
cone, and satisfies \(c\cdot\omega_X>0\), because \(c^2>0\). The chosen
exceptional signs give \(E_j\cdot\omega_H\le0\), hence

\[
                          a_H\cdot\omega_H>0.             \tag{9.7}
\]

Thus the small-perturbation discriminants for the base and the selected
blow-up structure have the same sign in corresponding chamber
conventions. Adding \(S^1\times S^3\) leaves the second-cohomology
polarization unchanged. The gluing formulas consequently give a nonzero
invariant in the chamber containing zero perturbation for this selected
class, provided that was the specified nonzero base chamber.

For any metric with this polarization, take generic perturbations
tending to zero in that chamber. Nonzero invariant gives solutions, and
the usual fixed-metric Seiberg--Witten compactness gives an unperturbed
solution. A reducible limit would have \(F_A^+=0\), contradicting
\((a_H^+)^2\ge A>0\). Every metric conformal to \(g\) has the same
polarization and the same nonzero discriminant sign. Thus the
\emph{same} class (9.4) satisfies Definition 1.1 for every weight, as
required by (9.1).

For \(b_+>1\), there is no chamber ambiguity in this argument. For
\(b_+=1\), the statement must not be replaced by the assertion that
every exceptional sign choice is an ordinary all-metric monopole class.

\hypertarget{the-minimal-general-type-case}{%
\subsubsection{9.3. The minimal general-type
case}\label{the-minimal-general-type-case}}

If \(X\) is a minimal compact complex surface of general type with its
complex orientation, then

\[
                          A=c_1(X)^2=T_X>0.
\]

The canonical Spin\(^c\) structure has a nonzero Seiberg--Witten
invariant. For \(b_+>1\), the standard Kähler/symplectic nonvanishing
statement supplies the input directly. For \(b_+=1\), the canonical
class \(K\) has \(K^2>0\) and positive pairing with a Kähler form. It is
therefore in the forward time-like cone and pairs positively with every
forward period ray. For \(c=-K\), the small-perturbation chamber agrees
with the canonical nonzero chamber, up to the simultaneous choice of
sign convention for the chamber discriminant. This is the
light-cone/chamber argument underlying
\href{https://arxiv.org/pdf/math/0003068}{LeBrun's polarized
construction, Proposition 3.2}. The chamber convention and canonical
nonvanishing can also be read in
\href{https://arxiv.org/pdf/math/0012048}{Li--Liu, §2.1 and Theorem
2.2}.

The positivity of \(K\) on the symplectic or Kähler form matters:
\(b_+=1\) and \(K^2>0\) alone do not identify the nonzero
small-perturbation chamber for an arbitrary symplectic manifold.

\hypertarget{the-einstein-obstruction-and-its-precise-relation-to-lebrun}{%
\subsection{10. The Einstein obstruction and its precise relation to
LeBrun}\label{the-einstein-obstruction-and-its-precise-relation-to-lebrun}}

\textbf{Theorem 10.1 (connected-sum integral obstruction).} Under the
gauge-theoretic and chamber assumptions of §9.1, every metric on (9.3)
satisfies

\[
 \mathcal E(g)+\left(\frac{T_M}{A}-\frac23\right)_+
                         \ge\varepsilon_0.               \tag{10.1}
\]

It also satisfies the slightly weaker but base-topological statement

\[
 \boxed{\displaystyle
 \mathcal E(g)+\left(\frac13-\frac{D}{T_X}\right)_+
                         \ge\varepsilon_0.}              \tag{10.2}
\]

For minimal general type, \(A=T_X=c_1(X)^2\), so the two statements
coincide.

\textbf{Proof.} Equation (9.5) and the conformal solvability just proved
give all the hypotheses of Theorem 1.2. If \(T_M>0\), then \(A\ge T_X\)
implies \(T_M/A\le T_M/T_X\), yielding (10.2). If \(T_M\le0\), both
positive-part terms are zero. This case split avoids reversing an
inequality by dividing a negative numerator by the larger denominator.
\(\square\)

In particular, put \(\gamma=\varepsilon_0/2\). Then

\[
             D\ge(1/3-\gamma)T_X
 \quad\Longrightarrow\quad
             \mathcal E(g)\ge\varepsilon_0/2
             \quad\hbox{for every metric }g.             \tag{10.3}
\]

There is no Einstein metric in this band: nonzero Einstein constants
contradict (10.3), and zero Einstein constant contradicts (2.13). At
either normalized sign, (8.2) also gives

\[
                    \delta_\sigma(g)\ge\sqrt{\varepsilon_0/8}.
 \tag{10.4}
\]

An Einstein metric would in fact require \(D/T_X\le1/3-\varepsilon_0\).
The equality endpoint of that necessary condition is not excluded by
(10.2) alone; the closed band (10.3) avoids that endpoint issue.

\hypertarget{exact-source-scope}{%
\subsubsection{10.2. Exact source scope}\label{exact-source-scope}}

The version of \href{https://arxiv.org/pdf/math/0003068}{LeBrun's paper
verified here} is arXiv:math/0003068v1, March 2000. Its Proposition 3.2
constructs the metricwise polarized classes on the connected sums in
(9.3). Its Theorem 3.3 states the classical threshold

\[
                              k+4\ell\ge T_X/3,
\]

and Corollary 3.4 specializes the result to blow-ups of minimal complex
surfaces of general type. We use the ordinary nontrivial Seiberg--Witten
input, including its insertion and chamber meaning, not merely the
topology of a manifold with the same Euler characteristic and signature.
The proof above handles the positive-budget case of the general theorem,
not only its general-type corollary.

The numerical band below the classical threshold is nonvacuous in
families with unbounded budget. For example, let
\(X_m=\Sigma_2\times\Sigma_m\), \(m\ge2\). Then \(A=T_{X_m}=8(m-1)\).
With \(\ell=0\) and \(k=\lceil A/3\rceil-1\),

\[
                       0<1/3-k/A\le1/A\longrightarrow0.
\]

For sufficiently large \(m\), these ratios lie in (10.3). Since the
universal analytic constant has not been made numerical, this argument
does not specify the first qualifying \(m\).

\hypertarget{nonpositive-budgets-and-the-exceptional-zero-corner}{%
\subsubsection{10.3. Nonpositive budgets and the exceptional zero
corner}\label{nonpositive-budgets-and-the-exceptional-zero-corner}}

There is a direct all-metric formula, independent of monopoles. For
\(R_2>0\), (1.1) and (8.1) give

\[
 \mathcal E(g)
   =1+\frac{12\int|W^+|^2-24\pi^2T+L_s}{R_2}.           \tag{10.5}
\]

Thus \(T\le0\) implies \(\mathcal E(g)\ge1\). For a Ricci-flat metric
its value is \(4\) by definition, so the same lower bound holds. This
lower bound is not a Ricci-flat exclusion.

For an Einstein metric, (1.1) reduces to

\[
               4\pi^2T=\int(s^2/24+2|W^+|^2)\ge0.
\]

Consequently, when \(T_X<0\), or when \(T_X=0\) and \(D>0\), the
connected sum has \(T_M<0\) and admits no Einstein metric. These are
classical Chern--Gauss--Bonnet consequences and require no
positive-budget monopole argument.

When \(T_X=D=0\), however, equality permits Ricci-flat, anti-self-dual
metrics in the chosen orientation. K3 and the four-torus provide the
standard zero-budget examples. The unqualified statement of Theorem 3.3
in the verified preprint includes this corner, but its appeal to the
Hitchin--Thorpe inequality does not exclude these equality cases. We
therefore do \textbf{not} claim an unqualified subsumption of that
literal zero-corner wording, nor infer an exception-free statement for a
separately unverified published version. The conclusions established
here are: the positive-budget integral strengthening with the stated
gauge input; the classical negative-\(T_M\) exclusion; and an explicit
reservation of the Ricci-flat zero corner.

\hypertarget{logical-dependencies-and-limitations}{%
\subsection{11. Logical dependencies and
limitations}\label{logical-dependencies-and-limitations}}

The analytic proof has the following complete chain:

\[
 \begin{gathered}
 \text{one actual conformal Spin}^c\text{ structure}
 \ \Longrightarrow\ \text{solutions for smooth truncated weights}\\
 \Longrightarrow\ \rho\le4\text{ and small retained first-order defects}
 \ \Longrightarrow\ W^+\approx A_\psi\text{ with }A_\psi\text{ smooth}\\
 \Longrightarrow\ \text{bundle Hessian budget and weak Bianchi upper bound}\\
 \Longrightarrow\ \text{incompatibility with the exact positive divergence term}.
 \end{gathered}
\]

The proof neither differentiates an eigenprojection nor normalizes
\(\psi\) across its zero set. Smooth approximation of weights is used
only for numerical inequalities at a fixed metric; there is no claimed
uniform regularity or limiting solution for a nonsmooth weight. The
Hessian estimate uses the positive-spinor connection, which is why no
\(L^2\) budget for \(W^-\) enters. The topological theorem additionally
needs the class-square bound and the actual class/chamber construction;
the scalar curvature inequality by itself is not a substitute for those
inputs.

This argument proves an obstruction and an integral gap. It does not
prove convergence of near-Einstein metrics, classify all equality
metrics of other curvature functionals, produce an Einstein metric, or
settle arbitrary pointwise Ricci-pinching problems without the
monopole/topological hypotheses displayed above.

\hypertarget{references-and-provenance}{%
\subsection{References and provenance}\label{references-and-provenance}}

\begin{itemize}
\tightlist
\item
  Claude LeBrun, \emph{Ricci Curvature, Minimal Volumes, and
  Seiberg--Witten Theory},
  \href{https://arxiv.org/pdf/math/0003068}{arXiv:math/0003068v1},
  especially Proposition 3.2, Theorem 3.3 and Corollary 3.4. The source
  version and the zero-budget caveat are specified in §10.
\item
  Peter Ozsváth and Zoltán Szabó, \emph{Higher Type Adjunction
  Inequalities in Seiberg--Witten Theory},
  \href{https://arxiv.org/pdf/math/0005268}{arXiv:math/0005268},
  especially Proposition 2.2, Remark 2.3 and Theorem 3.2. These supply
  insertion-aware surgery and blow-up formulas.
\item
  Tian-Jun Li and Ai-Ko Liu, \emph{Uniqueness of Symplectic Canonical
  Class, Surface Cone and Symplectic Cone of 4-Manifolds with
  \(b^+=1\)},
  \href{https://arxiv.org/pdf/math/0012048}{arXiv:math/0012048}, §2.1
  and Theorem 2.2, for chamber conventions and canonical nonvanishing.
\item
  Claude LeBrun, \emph{Four-Manifolds, Curvature Bounds, and Convex
  Geometry}, \href{https://arxiv.org/pdf/math/0611450}{arXiv preprint},
  especially Lemma 3.7, for the conformal weighted-equation input.
\end{itemize}

This text consolidates the local weighted-monopole,
integral-calibration, all-sign, chamber-selection and LeBrun-comparison
arguments. Those earlier working records are provenance, not proof
dependencies of the analytic argument presented here. Numerical and
symbolic checks are useful for detecting transcription errors; none is
used in place of a mathematical identity or a gauge-theoretic existence
theorem.

END COMPLETE SOURCE

\begin{center}\rule{0.5\linewidth}{0.5pt}\end{center}

\hypertarget{source-2-round21_analytic.md}{%
\subsection{Source 2:
round21\_analytic.md}\label{source-2-round21_analytic.md}}

Original path: \nolinkurl{mathematics/current/round21_analytic.md}

SHA-256:
\nolinkurl{9e39e932ae9f109d9ab3f959a243046e243beb20900c3f4e95be8dc424f5eb22}

BEGIN COMPLETE SOURCE

\hypertarget{integral-rigidity-for-bounded-self-dual-forms}{%
\section{Integral rigidity for bounded self-dual
forms}\label{integral-rigidity-for-bounded-self-dual-forms}}

This note isolates the analytic mechanism behind the integral monopole
gap. It proves an abstract theorem whose input is a bounded self-dual
two-form with a Hodge-Laplacian budget, rather than a monopole spinor.
This makes the production of a calibration and its differential budget
separate inputs. It is an analytic criterion; no new family of smooth
topologies is asserted to satisfy its hypotheses.

All manifolds are smooth, closed, oriented and four-dimensional.
Curvature is normalized relative to \(\operatorname{Ric}=-g\). Two-forms
have the exterior-algebra norm; \(W^+\) has the Hilbert--Schmidt
operator norm. Write \(\fint=\operatorname{Vol}^{-1}\int\), and let
\(\Delta_H\) denote the nonnegative Hodge Laplacian.

\hypertarget{the-bounded-two-form-criterion}{%
\subsection{1. The bounded two-form
criterion}\label{the-bounded-two-form-criterion}}

For a smooth self-dual two-form \(\psi\), set

\[
 q=|\psi|^2,\qquad
 A_\psi={q\over24}I-{1\over8}\psi\otimes\psi,
 \qquad R=W^+-A_\psi.
\]

Define the defect and source budget

\[
 d=\fint_M\left(|\operatorname{Ric}+g|^2+|R|^2+(q-2)^2\right),
 \qquad B=\fint_M|\Delta_H\psi|^2.
\]

\textbf{Theorem.} For each finite \(K>0\), there is \(c(K)>0\) such that
\(\|\psi\|_\infty\le K\) implies

\[
                    \boxed{(1+B)d\ge c(K).}              \tag{1}
\]

\textbf{Small-source corollary.} For a possibly smaller \(c_0(K)>0\),

\[
 \boxed{\quad
 \fint_M\left(
 |\operatorname{Ric}+g|^2+|R|^2+(q-2)^2+|\Delta_H\psi|^2
 \right)\ge c_0(K).
 \quad}                                                     \tag{1a}
\]

No spin structure, monopole class, curvature bound, scalar sign, Sobolev
constant, injectivity radius or diameter bound is required. In
particular, for bounded \emph{harmonic} self-dual forms \(B=0\), giving
a gap for the three zeroth-order defects alone. No small
first-derivative defect is assumed, even in the nonharmonic case. The
boundedness of the form and its Hodge-Laplacian control are genuine
hypotheses, not consequences of a cohomology class alone.

The theorem excludes a specific incompatible calibration: Ricci
curvature near \(-g\), Weyl curvature near \(A_\psi\), and \(|\psi|^2\)
near two. For a product of curvature-minus-one surfaces the simple Weyl
eigenvalue is \(-2/3\), not the target \(-1/6\). Thus the theorem does
not exclude the realizable hyperbolic-product Weyl model and does not
assert a universal Ricci gap without the other summands.

\hypertarget{re-derivation-of-the-bundle-hessian-estimate}{%
\subsection{2. Re-derivation of the bundle Hessian
estimate}\label{re-derivation-of-the-bundle-hessian-estimate}}

Let \(u\) be a smooth section of a metric bundle, with compatible
connection \(\nabla\). Set \(H_{ij}=\nabla_i\nabla_j u\), including the
cotangent Levi-Civita correction, \(L=\operatorname{tr}H\), and
\(F_{ij}=[\nabla_i,\nabla_j]\). Use the component convention
\(|F|^2=\sum_{ij}|F_{ij}|_{\rm HS}^2\).

First, integration of the divergence of
\(|\nabla u|^2\langle u,\nabla u\rangle^\sharp\) gives

\[
 \|\nabla u\|_4^2\le(2+\sqrt4)\|u\|_\infty\|H\|_2.
                                                               \tag{2}
\]

This is a covariant integration-by-parts inequality, not a Sobolev
embedding. To check the curvature terms, first interchange the two
Hessian slots:

\[
 \int|H|^2-\sum_{ij}\int\langle H_{ij},H_{ji}\rangle
 =\frac12\sum_{ij}\int|F_{ij}u|^2.
\]

Integrate the second expression once and then commute derivatives on the
bundle-valued one-form \(\nabla u\). The cotangent curvature gives
Ricci; the bundle curvature acts on \(\nabla u\). Integrating the trace
term once more yields

\[
\begin{aligned}
 \int|H|^2={}&\int|L|^2
 -\int\operatorname{Ric}_{ij}\langle\nabla_i u,\nabla_j u\rangle\\
 &+\sum_{ij}\int\langle F_{ij}\nabla_j u,\nabla_i u\rangle
 +\frac12\sum_{ij}\int|F_{ij}u|^2.                         
\end{aligned}\tag{3}
\]

There is no differentiated-curvature term. Swapping the Hessian slots
\emph{before} commuting is important. The scalar specialization fixes
the Ricci sign. With \(K_u=\|u\|_\infty\), Holder, (2), and Young give

\[
 \|\nabla^2u\|_2^2\le
 2\|L\|_2^2+32K_u^2\|\operatorname{Ric}\|_2^2
                         +33K_u^2\|F\|_2^2.              \tag{4}
\]

This recovers the inherited estimate, including its coefficient and the
absence of a hidden Sobolev or curvature-derivative hypothesis.

\hypertarget{proof-of-the-two-form-criterion}{%
\subsection{3. Proof of the two-form
criterion}\label{proof-of-the-two-form-criterion}}

Assume \(d\le1\); otherwise the conclusion is immediate after reducing
\(c(K)\). Put \(\mathcal B=1+B+d\). All constants below depend only on
\(K\) and the dimension. Put \(E=\operatorname{Ric}+g\) and
\(\eta=\Delta_H\psi\).

\hypertarget{the-self-dual-connection-supplies-the-second-order-budget}{%
\subsubsection{3.1 The self-dual connection supplies the second-order
budget}\label{the-self-dual-connection-supplies-the-second-order-budget}}

The curvature of \(\Lambda^+\) is the representation of the Riemann
operator projected to its self-dual output. Hence

\[
 |F^{\Lambda^+}|^2\le C\left(|W^+|^2+s^2+
                               |\operatorname{Ric}^{\circ}|^2\right).
                                                               \tag{5}
\]

The \(W^-\) block does not act on \(\Lambda^+\). Since
\(|A_\psi|^2=q^2/96\) and \(q\le K^2\), the \(E,R\) budgets give

\[
 \fint\left(|\operatorname{Ric}|^2+|F^{\Lambda^+}|^2\right)
 \le C_K(1+d).
\]

The self-dual Weitzenbock equation is

\[
 \nabla^*\nabla\psi=\eta+2W^+\psi-(s/3)\psi.
\]

Apply (2)--(4) to \(u=\psi\). Differentiating the quadratic polynomial
\(A_\psi\) and the scalar \(q\) then gives

\[
 \fint\left(|\nabla^2\psi|^2+|\nabla\psi|^4+
              |\nabla^2A_\psi|^2+|\Delta q|^2\right)
 \le C_K\mathcal B.
                                                               \tag{6}
\]

Here \(\Delta=d^*d\) on functions, and
\(\Delta q=2\langle\nabla^*\nabla\psi,\psi\rangle -2|\nabla\psi|^2\).
All estimates are extensive before division by volume; no lower volume
bound enters.

\hypertarget{amplitude-control-and-a-coercive-divergence-identity}{%
\subsubsection{3.2 Amplitude control and a coercive divergence
identity}\label{amplitude-control-and-a-coercive-divergence-identity}}

The amplitude itself need not have a smooth square root at its zeros.
Use the globally smooth quantity \(q\), rather than divide by
\(|\psi|\). Integration by parts and (6) give

\[
 \fint|dq|^2=\fint(q-2)\Delta q
                  \le C_K\sqrt{d\mathcal B}.             \tag{7}
\]

There is a useful pointwise identity, valid for arbitrary smooth
self-dual forms without first-order equations. Let
\(J=\operatorname{div}W(A_\psi)\). The full contraction in the checker
gives

\[
 |J|^2={q|\nabla\psi|^2\over64}+{q|d\psi|^2\over32}
       -{|dq|^2\over192}
       +{\langle\iota_{\nabla q}\psi,*d\psi\rangle\over32}.
                                                               \tag{8}
\]

The last inner product is at most \(\sqrt q\,|dq|\,|d\psi|\) in absolute
value. Young's inequality therefore gives the pointwise coercivity

\[
 \boxed{\quad
 q\bigl(|\nabla\psi|^2+|d\psi|^2\bigr)
        \le64|J|^2+{4\over3}|dq|^2.
 \quad}                                                     \tag{8a}
\]

This is polynomial at zeros, and introduces no division by \(q\). It
will control the derivative terms for the two-sided selection theorem in
Section 6. Small amplitude error and a second-order budget now supply
the weighted derivative information previously obtained from monopoles.

\hypertarget{weak-bianchi-gives-a-small-self-test}{%
\subsubsection{3.3 Weak Bianchi gives a small
self-test}\label{weak-bianchi-gives-a-small-self-test}}

Let \(W(A)\) be the algebraic Weyl four-tensor represented by \(A\), and
\(J=\operatorname{div}W(A_\psi)\). The projections to \(\Lambda^+\) are
parallel. The contracted Bianchi identity, followed by integration by
parts, gives for every smooth test tensor \(T\)

\[
 \left|\int\langle J,T\rangle\right|
 \le C\left(\|E\|_2+\|R\|_2\right)\|\nabla T\|_2.      \tag{9}
\]

This is a weak derivative estimate, not an estimate on \(\nabla E\). The
tensor \(J\) is smooth, so \(T=J\) is an admissible test. Moreover
\(|\nabla J|\le C|\nabla^2A_\psi|\). Therefore (6) implies

\[
                 \fint|J|^2\le C_K\sqrt{d\mathcal B}.    \tag{10}
\]

\hypertarget{the-same-self-test-has-a-positive-lower-bound}{%
\subsubsection{3.4 The same self-test has a positive lower
bound}\label{the-same-self-test-has-a-positive-lower-bound}}

The inherited weighted Hodge and divergence identities admit the
following entirely polynomial version:

\[
\begin{aligned}
 \int|J|^2=\int\Big(&-{s q^2\over192}-{q^3\over384}
       +{q\langle R\psi,\psi\rangle\over32}
       +{q|d\psi|^2\over16}-{5|dq|^2\over384}\\
       &+{\langle\iota_{\nabla q}\psi,*d\psi\rangle\over16}
                           \Big).                         
\end{aligned}\tag{11}
\]

It follows by testing the self-dual Hodge equation against \(q\psi\) and
contracting the derivative of \(A_\psi\); the pointwise contraction is
independently checked in the accompanying program. Because (11) is
polynomial in \(\psi,\nabla\psi\), it extends across all zeros of
\(\psi\); no spinor-square representation is used.

At \(s=-4,q=2,R=0\), its algebraic part equals \(1/16\). On
\(0\le q\le K^2\), polynomial Lipschitz estimates give an averaged error
bounded by \(C_K\fint(|s+4|+|q-2|+|R|)\le C_K\sqrt d\). The last term is
bounded below by \(-\{q|d\psi|^2+|dq|^2\}/32\). Absorb its first
contribution into the positive \(q|d\psi|^2/16\) already present in
(11). The remaining derivative error is at worst \(-17|dq|^2/384\),
controlled by (7). In particular this argument does not assume \(d\psi\)
small. Consequently

\[
           \fint|J|^2\ge {1\over16}-C_K\sqrt{d\mathcal B}.\tag{12}
\]

Combining (10) and (12) yields \(1/16\le C'_K\sqrt{d(1+B+d)}\). For
\(d\le1\), \(1+B+d\le2(1+B)\), proving (1) after choosing a sufficiently
small \(c(K)>0\). For (1a), if its left side is \(z\le1\), then
\(d\le z\), \(B\le z\), and \((1+B)d\le2z\). Thus
\(c_0(K)=\min\{1,c(K)/2\}\) suffices.

\hypertarget{what-the-monopole-proof-actually-requires}{%
\subsection{4. What the monopole proof actually
requires}\label{what-the-monopole-proof-actually-requires}}

For the inherited integral gap, one fixes a metric \(g\) and a single
spin-c structure with \(p=(a_g^+)^2>0\). The analytic construction only
requires ordinary unperturbed Seiberg--Witten solutions for every metric
in \([g]\). Indeed, solving on \(\widehat g=f^{-2}g\) gives on \(g\)

\[
 D_A\Phi=0,\qquad -iF_A^+=f\sigma(\Phi),
\]

for every positive smooth weight \(f\). The self-dual polarization is
conformally invariant. Every remaining Hessian, Hodge and Bianchi step
takes place on this one metric. Neither \(b_+>1\) nor persistence over
metrics outside \([g]\) is used in those analytic steps.

This is already the scope of \href{round13_chamber.md}{the Round13
chamber result}, not a new Round21 extension. In \(b_+=1\) applications,
chamber invariance and passage from perturbed invariants to unperturbed
solutions still have to be established. The Hessian identity cannot
supply them.

For the monopole argument specifically, the determinant connection
energy is controlled by

\[
 \|F_A\|_2^2=2\|F_A^+\|_2^2-4\pi^2a^2.
\]

That is the origin of the negative-square cost. The theorem in Section 1
avoids this connection by acting directly on \(\Lambda^+\), but replaces
the monopole PDE with a bounded-form and Hodge-budget hypothesis. It
therefore does \textbf{not} remove the square cost from the existing
monopole theorem: the harmonic projection need not preserve a maximum
bound, and a small \(d\psi\) does not bound \(\Delta_H\psi\) by itself.

More precisely, let \(\kappa\) and \(B_a=(-a^2)_+/V\) have the meanings
in \href{round4_rigidity.md}{Round4, Theorem 1}. The inherited
bounded-spinor, Hessian and fourth-power-gradient estimates imply

\[
 B=\fint|\Delta_H\psi|^2\le C(1+\kappa+B_a)+o(1).
\]

Indeed, the quadratic map from \(\Phi\) to \(\psi\) gives
\(|\nabla^2\psi|\le C(|\Phi||\nabla_A^2\Phi|+|\nabla_A\Phi|^2)\); the
two-form Weitzenbock formula adds only the already controlled \(W^+,s\)
terms. The Weyl approximation and amplitude estimates give
simultaneously

\[
                         d\le C\kappa+o(1).
\]

The bounded-form theorem therefore returns
\(\kappa(1+\kappa+B_a)\ge c\), and hence the inherited
\((1+B_a)\kappa\ge c'\) after splitting \(\kappa\le1\) and \(\kappa>1\).
This is a structural recovery of the earlier calibration gap, not a
second independent topological theorem. It displays precisely where a
different construction of a bounded form could replace the monopole
input, and which quantitative budget that construction must supply.

\hypertarget{research-use-and-verification-scope}{%
\subsection{5. Research use and verification
scope}\label{research-use-and-verification-scope}}

The new analytic formulation offers a distinct input route: a family of
bounded, almost-harmonic self-dual forms with the indicated calibration
would be impossible even without monopole solvability. For a harmonic
form, only its maximum bound and the three zeroth-order calibration
defects remain. Producing these hypotheses on a new topological family
is an open step; a formal cohomology vector does not suffice.

The checker verifies the Hessian curvature coefficients on a
noncommuting constant connection, the full twelve-variable divergence
contraction, and the polynomial conversion used in (11). The general
integration-by-parts and PDE arguments are the proofs above, not
consequences of running the checker. The inherited statements have been
checked at these analytic interfaces; this note does not certify every
earlier topological input.

\hypertarget{quantitative-selection-of-the-kuxe4hler-calibration}{%
\subsection{6. Quantitative selection of the Kähler
calibration}\label{quantitative-selection-of-the-kuxe4hler-calibration}}

The forbidden Weyl model is one member of a family, not an isolated
numerical accident. For a constant real parameter \(t\), put

\[
 A_{\psi,t}=t\left({q\over3}I-\psi\otimes\psi\right),
 \quad R_t=W^+-A_{\psi,t},\quad
 d_t=\fint\left(|E|^2+|R_t|^2+(q-2)^2\right).
                                                               \tag{13}
\]

Thus \(t=1/8\) is the model in Theorem (1). At \(q=2\), the Weyl
eigenvalues prescribed by \(A_{\psi,t}\) are \((-4t/3,2t/3,2t/3)\). The
negative Kähler--Einstein value is \(t=1/2\).

\textbf{Branch-selection theorem.} Fix \(K,T<\infty\). For
\(\|\psi\|_\infty\le K\) and \(|t|\le T\),

\[
 \boxed{\quad |t|\,|1-2t|\le C(K,T)
                  \sqrt{d_t(1+B+d_t)}.\quad}             \tag{14}
\]

A quantitative spectral-selection consequence is

\[
 \boxed{\quad
 \operatorname{dist}(t,\{0,1/2\})^2
       \le C(K,T)(1+B)d_t\qquad(d_t\le1).
 \quad}                                                   \tag{14a}
\]

Indeed, \(\operatorname{dist}(t,\{0,1/2\})\le2|t|\,|1-2t|\) for every
real \(t\), since the larger of the distances to zero and one-half is at
least one-quarter. Also \(1+B+d_t\le2(1+B)\). The two branches are
vanishing Weyl amplitude and the negative Kähler--Einstein amplitude,
not an arbitrary list of forbidden constants.

Consequently \((1+B)d_t\) has a uniform positive lower bound when \(t\)
is restricted to a compact subset of \(\mathbb R\setminus\{0,1/2\}\). On
any parameter set \(t_0\le|t|\le T\), with \(t_0>0\), one further has

\[
 |2t-1|+\fint\left(|\nabla\psi|^2+|d\psi|^2\right)
 \le C(K,t_0,T)\sqrt{d_t(1+B+d_t)}\qquad(d_t\le1).       \tag{15}
\]

Thus a bounded-source-budget calibration with vanishing defect and
nonzero Weyl amplitude selects the Kähler value and an almost-parallel
form. This is an integral conclusion; it does not supply coordinates,
pointwise metric closeness or compactness.

\textbf{Proof.} Let \(J_t=\operatorname{div}W(A_{\psi,t})\). Since
\(A_{\psi,t}=8tA_\psi\), one has \(J_t=8tJ\), with \(J\) here denoting
the divergence of the original polynomial on the same form and metric.
The second-order estimates remain valid uniformly for bounded \(t\),
with \(\mathcal B_t=1+B+d_t\). In the weak Bianchi test the derivative
of \(J_t\) has one factor \(|t|\), so

\[
                  \fint|J_t|^2\le C_{K,T}\,|t|
                                      \sqrt{d_t\mathcal B_t}. \tag{16}
\]

The weighted Hodge contraction becomes exactly

\[
\begin{aligned}
 \int|J_t|^2=\int\Big(&-{t^2s q^2\over3}
       -{4t^3q^3\over3}+2t^2q\langle R_t\psi,\psi\rangle
       +4t^2q|d\psi|^2-{5t^2\over6}|dq|^2\\
       &+4t^2\langle\iota_{\nabla q}\psi,*d\psi\rangle
                    \Big).                               
\end{aligned}\tag{17}
\]

Its algebraic part at \(s=-4,q=2,R_t=0\) is \(P(t)=(16/3)t^2(1-2t)\).
The pointwise coercivity (8a) rescales to

\[
 t^2q\bigl(|\nabla\psi|^2+|d\psi|^2\bigr)
          \le |J_t|^2+{4\over3}t^2|dq|^2.                \tag{18}
\]

In (17), the absolute derivative contribution is at most
\(6|J_t|^2+(65/6)t^2|dq|^2\): the positive closedness term costs at most
\(4|J_t|^2+(16/3)t^2|dq|^2\), and Young's inequality bounds the cross
term by \(2|J_t|^2+(14/3)t^2|dq|^2\). This uses (18), not an assumption
that \(d\psi\) is small. Together with the zeroth-order algebraic error
and the amplitude estimate (7),

\[
 |P(t)|\le 7\fint|J_t|^2+
                  C_{K,T}t^2\sqrt{d_t\mathcal B_t}.       \tag{18a}
\]

This is valid on both sides of \(t=1/2\), including negative \(t\).
Combine (16) and (18a), divide by \(|t|>0\), and use \(|t|\le T\),
obtaining (14). The case \(t=0\) is automatic. For \(d_t>1\), (14)
follows after enlarging the constant since its left side is bounded and
its right square root is at least one.

For (15), (14) first bounds \(|2t-1|\) by its right side divided by
\(t_0\). On \(\{q\ge1\}\), (18), (16) and (7) give an averaged bound
\(C_{K,t_0,T}\sqrt{d_t\mathcal B_t}\) for \(|\nabla\psi|^2+|d\psi|^2\).
On \(\{q<1\}\), the relative volume is at most \(\fint(q-2)^2\le d_t\).
The fourth-power gradient estimate (6), and
\(|d\psi|^2\le3|\nabla\psi|^2\), bound its contribution by
\(C_{K,T}\sqrt{d_t\mathcal B_t}\). This proves (15) without assuming
closedness.

For comparison, the unweighted Hodge identity is

\[
 \fint|\nabla\psi|^2
 =2\fint|d\psi|^2-
  \fint\left({s q\over3}+{4tq^2\over3}
                               -2\langle R_t\psi,\psi\rangle\right).
                                                               \tag{19}
\]

The parenthesized expression differs in average from \((8/3)(2t-1)\) by
at most \(C_{K,T}\sqrt{d_t}\). Nonnegativity of the left side gives the
one-sided estimate

\[
 (2t-1)_+\le C_{K,T}\left(\sqrt{d_t}+\fint|d\psi|^2\right).
                                                               \tag{20}
\]

It becomes source-budget-free when \(d\psi=0\), but not for arbitrary
forms under zeroth-order control alone.

At zero defect and \(t\ne0\), (14) forces \(t=1/2\), and (15) gives
\(\nabla\psi=0\). The norm-squared-two self-dual form is then the
parallel Kähler form of an orthogonal complex structure, and
\(\operatorname{Ric}=-g\) makes the metric negative Kähler--Einstein.
This endpoint is genuinely realized; parallelness of the Kähler form
does not imply that the metric is a global product. At \(t=0\), the
divergence polynomial loses sensitivity; (14) makes no assertion that an
exact calibrated example exists there.

This extension identifies the realizable equality model of the analytic
mechanism. It does not change which monopole classes are available, and
it does not by itself extend the topological range of the earlier
integral gap.

\hypertarget{the-same-analytic-criterion-for-either-einstein-sign}{%
\subsection{7. The same analytic criterion for either Einstein
sign}\label{the-same-analytic-criterion-for-either-einstein-sign}}

The spectral-selection argument itself is not restricted to the negative
Einstein sign. For \(\sigma\in\{-1,+1\}\), put

\[
E_\sigma=\operatorname{Ric}-\sigma g,\qquad
d_{\sigma,t}=\fint\left(
|E_\sigma|^2+|W^+-A_{\psi,t}|^2+(q-2)^2\right).
\]

With the same bounded-form and bounded-parameter hypotheses,

\[
\boxed{\quad
|t|\,|\sigma+2t|
\le C(K,T)\sqrt{d_{\sigma,t}(1+B+d_{\sigma,t})}.
\quad}                                                     \tag{21}
\]

In particular, for \(d_{\sigma,t}\le1\),

\[
\operatorname{dist}(t,\{0,-\sigma/2\})^2
\le C(K,T)(1+B)d_{\sigma,t}.                               \tag{22}
\]

To verify this, replace the scalar target \(-4\) by \(4\sigma\) in (17).
Its calibrated algebraic value becomes \[
P_\sigma(t)=-\frac{16}{3}t^2(\sigma+2t).
\] The curvature budget uses \(|\operatorname{Ric}|\le|E_\sigma|+2\),
the amplitude estimate is unchanged, and weak Bianchi differentiates
neither constant multiple of the metric. Equations (16), (18) and (18a)
therefore hold with this \(P_\sigma\) and \(d_{\sigma,t}\); division by
\(|t|\) proves (21), including its trivial \(t=0\) case. The distance
estimate follows as before. For \(|t|\ge t_0>0\), the same
good/bad-amplitude argument controls
\(\fint(|\nabla\psi|^2+|d\psi|^2)\).

At zero defect and \(t\ne0\), this forces \(t=-\sigma/2\) and a parallel
norm-squared-two self-dual form. The metric is then Kähler--Einstein
with \(\operatorname{Ric}=\sigma g\). The positive Kähler--Einstein
amplitude is \(-1/2\), not \(+1/2\). This extension concerns the
analytic criterion only: it does not supply positive-side monopole
classes or a Yau26 positive-branch obstruction.

END COMPLETE SOURCE

\begin{center}\rule{0.5\linewidth}{0.5pt}\end{center}

\hypertarget{source-3-line1_functional_pencil.md}{%
\subsection{Source 3:
\nolinkurl{line1_functional_pencil.md}}\label{source-3-line1_functional_pencil.md}}

Original path: \nolinkurl{mathematics/current/line1_functional_pencil.md}

SHA-256:
\nolinkurl{13f0324dab4999ffba998cd60e8a1d8c152b97aae6f40218aca599230c4201ae}

BEGIN COMPLETE SOURCE

\hypertarget{quantitative-non-saturation-of-scalarweyl-functionals-near-negative-einstein-metrics}{%
\section{Quantitative non-saturation of scalar--Weyl functionals near
negative Einstein
metrics}\label{quantitative-non-saturation-of-scalarweyl-functionals-near-negative-einstein-metrics}}

23 September 2026

\hypertarget{main-theorem}{%
\subsection{1. Main theorem}\label{main-theorem}}

Let \((M,g)\) be a closed connected oriented smooth four-manifold. Fix a
Spin\(^c\) structure, with real first Chern class \(a\), whose
unperturbed Seiberg--Witten equations have a solution for every metric
conformal to \(g\). The structure can depend on \(g\); it is fixed while
the conformal weight varies. Write \[
V=\operatorname{Vol}_g(M),\qquad p=(a_g^+)^2>0,\qquad
r=\fint_M|\operatorname{Ric}+g|^2,\qquad
b_a=\frac{(-a^2)_+}{V}.
\] Ricci uses the tensor norm, two-forms the exterior-algebra norm, and
\(W^+\) the Hilbert--Schmidt norm as an endomorphism of \(\Lambda^+\).
Set \(w=\sqrt6|W^+|\). For a constant \(\beta>0\), define the
dimensionless, nonnegative functional deficit \[
\mathfrak D_\beta(g,a)
=\frac{\beta+1}{\beta}\int_M(s^2+\beta w^2)
-72\pi^2p.
\tag{1}
\] Here dimensionless means invariant under constant metric rescaling;
\(\mathfrak D_\beta/V\), like \(r\), is measured after fixing the
Einstein target to \(-g\).

\textbf{Theorem 1 (functional non-saturation).} For every finite \(L>2\)
there is \(c_L>0\) such that, simultaneously for \(2\le\beta\le L\), \[
\boxed{
(1+b_a)\left(r+\frac{\mathfrak D_\beta(g,a)}V\right)
\ge c_L(\beta-2)^2.
}
\tag{2}
\] There is no scalar-curvature band, Sobolev constant,
injectivity-radius, diameter, volume, or full-curvature hypothesis.
Actual conformal monopole solvability and the displayed square cost are
the inputs.

If the selected class additionally satisfies
\(a^2\ge T:=2\chi(M)+3\tau(M)\), there is \(c'_L>0\) with \[
\boxed{
r+\frac{\mathfrak D_\beta(g,a)}V\ge c'_L(\beta-2)^2.
}
\tag{3}
\] In particular, for an exact metric \(\operatorname{Ric}=-g\) and such
a class, \[
\int(s^2+6\beta|W^+|^2)
\ge\frac{72\pi^2\beta}{\beta+1}p
+c'_L\frac{\beta}{\beta+1}(\beta-2)^2V.
\tag{4}
\]

The endpoint \(\beta=2\) must remain: a negative Kähler--Einstein
surface with its canonical Spin\(^c\) structure has \(r=0\) and
\(\mathfrak D_2=0\). Thus (2) is a compatibility theorem, not a
prohibition of the Kähler--Einstein branch.

The proof supplies all auxiliary forms and their joint analytic budget
from (1). It extends the former calibration at \(\beta=8\) to a
parameter family of curvature functionals. It does not, by itself,
identify a new smooth topology with no Einstein metric.

\hypertarget{the-exact-deficit-decomposition}{%
\subsection{2. The exact deficit
decomposition}\label{the-exact-deficit-decomposition}}

Put \[
q=-s+w,\quad q_+=\max(q,0),\quad q_-=\max(-q,0),\quad
\Delta_+=\int q_+^2-72\pi^2p.
\] The signed monopole estimate gives \(\Delta_+\ge0\). The relevant
classical inequality follows from the rescaled equations; the positive
part is needed when scalar curvature changes sign. Completing a square
gives the exact identity \[
\boxed{
\mathfrak D_\beta
=\Delta_++\int q_-^2+\frac1\beta\int(s+\beta w)^2.
}
\tag{5}
\] Indeed
\((\beta+1)(s^2+\beta w^2)/\beta =(-s+w)^2+(s+\beta w)^2/\beta\).

Write \(e=s+4\), \(z=s+\beta w\), and \[
q_\beta=4\frac{\beta+1}{\beta},\qquad
\kappa=r+\mathfrak D_\beta/V.
\] Then \[
w-4/\beta=(z-e)/\beta,\qquad
q-q_\beta=z/\beta-(1+1/\beta)e,\qquad
\int e^2\le4Vr.
\tag{6}
\] Consequently, uniformly for \(2\le\beta\le L\), \[
\int\bigl[(q-q_\beta)^2+(w-4/\beta)^2\bigr]
\le C_L V\kappa,
\quad
\int(|\operatorname{Ric}|^2+s^2+|W^+|^2)
\le C_L V(1+\kappa).
\tag{7}
\] These bounds control the Weyl norm, not its eigenspaces. The latter
come from the actual monopole below.

\hypertarget{actual-bounded-spinors-with-no-scalar-band}{%
\subsection{3. Actual bounded spinors with no scalar
band}\label{actual-bounded-spinors-with-no-scalar-band}}

Set \[
Q=\max(q,3),\qquad
f_0=\max\{\sqrt Q,(-s)_+/4\}.
\] Choose smooth positive majorants \(f_j\to f_0\) uniformly on the
fixed compact manifold. Conformal covariance, applied to
\(\widehat g=f_j^{-2}g\), supplies solutions \[
D_A\Phi=0,\quad -iF_A^+=f_j\sigma(\Phi),\quad
\psi=2\sqrt2\sigma(\Phi),\quad
\rho=|\Phi|^2=|\psi|,\quad u=3\rho/2.
\tag{8}
\] We use \(|\sigma(\Phi)|^2=|\Phi|^4/8\). The spinor equation is
\(4\nabla_A^*\nabla_A\Phi+(s+f_j\rho)\Phi=0\). At a maximum of \(\rho\),
\(f_j\rho\le-s\), so \[
0\le\rho\le4,\qquad 0\le u\le6.
\tag{9}
\] Every approximation error below tends to zero for the fixed \(g,a\).
No uniform smoothing rate along a sequence and no bounds on derivatives
of \(f_j\) are required.

The cutoff costs only scalar error: \[
0\le Q^2-q_+^2\le9e^2,\qquad
0\le f_0^2-Q\le\frac9{16}e^2.
\tag{10}
\] For the second inequality, a nonzero difference forces
\((-s)_+>4\sqrt3>6\), whence \((-s)_+\le3|e|\). For the first, \(q<3\)
implies \(s>-3\), hence \(e>1\).

Define nonnegative defects \[
\begin{aligned}
D_0={}&2\int|d\rho|^2
+\int(4\rho|\nabla_A\Phi|^2-|\nabla\psi|^2)\\
&+2\int|d\psi|^2+2\int E_\lambda,\qquad
E_\lambda=\langle W^+\psi,\psi\rangle+w\rho^2/3\ge0.
\end{aligned}
\tag{11}
\] Nonnegativity uses the differential inequality for the parallel
spinor quadratic map and \(\lambda_{\min}(W^+)\ge-w/3\). The spinor and
self-dual-form Weitzenböck identities, subtracted after testing the
spinor against \(\rho\Phi\), give \[
\frac{27}{8}D_0=\int(qu^2-f_j u^3).
\tag{12}
\]

For completeness, the moment argument that retains the truncation cost
is as follows. Set \[
C=\int Qu^2,\quad Z=\int\sqrt Q\,u^3,\quad E_Q=\int Q^2.
\] Then \(C-Z\ge27D_0/8\ge0\), while the curvature equation gives \[
C_{\rm act}:=\int f_j^2u^2=18\|F_A^+\|_2^2\ge72\pi^2p,
\quad 0\le C_{\rm act}-C\le\frac{81}{4}\int e^2+o(1).
\] Hölder gives \(C\le E_Q^{1/3}Z^{2/3}\), hence \[
0\le E_Q-C\le S+o(1),\qquad
S:=\Delta_++\frac{117}{4}\int e^2\le117V\kappa.
\] If \(\theta=(Z/E_Q)^{1/3}\), the nonnegative Hölder slack implies
\(E_Q-C\ge2(C-Z)\). The associated amplitude estimate is
\(\int Q(u-\sqrt Q)^2\le E_Q-C\): with probability measure
\(Q^2d\mu/E_Q\), let \(h=u/\sqrt Q\), \(c=\mathbb E h^2\),
\(z_3=\mathbb E h^3\), and \(m=\mathbb E h\). Then
\(z_3\le c\le z_3^{2/3}\), so \(c\le1\); Cauchy--Schwarz gives
\(c^2\le m z_3\le mc\), and hence \(m\ge c\). Thus
\(\mathbb E(h-1)^2\le1-c\). We have proved \[
D_0\le\frac{4S}{27}+o(1),\qquad
\int Q(u-\sqrt Q)^2\le S+o(1).
\tag{13}
\]

\hypertarget{produce-the-full-calibrated-form-including-its-weyl-direction}{%
\subsection{4. Produce the full calibrated form, including its Weyl
direction}\label{produce-the-full-calibrated-form-including-its-weyl-direction}}

Because \(Q\ge3\), \(u\le6\), and \((Q-q)^2\le e^2\), (13) yields \[
\int(q-u^2)^2
\le2(1+2\sqrt3)^2S+2\int e^2+o(1)
\le C V\kappa+o(1).
\tag{14}
\] The middle inequality uses
\(|Q-u^2|\le(1+2\sqrt3)\sqrt Q\,|\sqrt Q-u|\). Together with (7), this
controls \(\int(u^2-q_\beta)^2\).

Now choose the normalization, rather than assuming it: \[
c_\beta^2=\frac{9\beta}{8(\beta+1)},\qquad
\omega=c_\beta\psi,\qquad t_\beta=1/\beta,\qquad
\alpha_\beta=\frac9{8(\beta+1)}.
\] It follows that \[
\|\omega\|_\infty\le3\sqrt2,\qquad
\int(|\omega|^2-2)^2\le C_L V\kappa+o(1).
\tag{15}
\] Define the smooth polynomial model \[
A_{\omega,t}=t\left(\frac{|\omega|^2}{3}I-\omega\otimes\omega\right).
\] The exact identity for arbitrary \(\psi\), including its zeros, is \[
\left|W^+-\alpha_\beta
       \left(\frac{\rho^2}{3}I-\psi\otimes\psi\right)\right|^2
=\frac16(w-2\alpha_\beta\rho^2)^2
+2\alpha_\beta E_\lambda.
\tag{16}
\] This identity uses only trace-freeness of \(W^+\), not an eigenline
selection. The tensor on its left is exactly \(W^+-A_{\omega,1/\beta}\).
Also \(2\alpha_\beta\rho^2=u^2/(\beta+1)\), and \[
w-\frac{u^2}{\beta+1}
=\left(w-\frac4\beta\right)
-\frac{u^2-q_\beta}{\beta+1}.
\] Equations (7), (13)--(16) give, for the same smooth solution, \[
d_\beta:=\fint\left(
 |\operatorname{Ric}+g|^2+|W^+-A_{\omega,1/\beta}|^2
+(|\omega|^2-2)^2\right)
\le C_L\kappa+o(1).
\tag{17}
\] In particular the small Weyl-direction error has been proved from the
actual gauge equations. It is not a hidden geometric hypothesis.

\hypertarget{supply-the-second-order-budget}{%
\subsection{5. Supply the second-order
budget}\label{supply-the-second-order-budget}}

For a bounded section \(v\) of a metric bundle with compatible
connection curvature \(\mathscr F\), covariant integration by parts on a
closed four-manifold gives \[
\|\nabla v\|_4^2\le4\|v\|_\infty\|\nabla^2v\|_2,
\] \[
\|\nabla^2v\|_2^2
\le2\|\nabla^*\nabla v\|_2^2
+32\|v\|_\infty^2\|\operatorname{Ric}\|_2^2
+33\|v\|_\infty^2\|\mathscr F\|_2^2.
\tag{18}
\] The curvature convention is
\(|\mathscr F|^2=\sum_{ij} |\mathscr F_{ij}|_{\rm HS}^2\). These are the
inherited Hessian identities, not a Sobolev embedding. One first
interchanges Hessian slots and then commutes on \(\nabla v\); no
derivative of curvature is introduced.

Equations (7), (10) give \(\int f_j^2\le C_L V(1+\kappa)+o(1)\). The
determinant Chern--Weil identity and (9) imply \[
\|F_A\|_2^2=2\|F_A^+\|_2^2-4\pi^2a^2
\le4\int f_j^2+4\pi^2(-a^2)_+.
\] The curvature of the positive-spinor connection is bounded in norm by
the \(W^+,s,\operatorname{Ric}^{\circ},F_A\) blocks; \(W^-\) acts
trivially on this representation. Its squared integral is therefore at
most \(C_LV(1+\kappa+b_a)+o(1)\). The spinor equation and \(|\Phi|\le2\)
likewise bound \(\int|\nabla_A^*\nabla_A\Phi|^2\) by
\(C_LV(1+\kappa)+o(1)\).

Apply (18) to \(\Phi\), and twice differentiate the parallel quadratic
map \(\Phi\mapsto\omega\). The resulting terms are bounded by
\(C(|\Phi||\nabla_A^2\Phi|+|\nabla_A\Phi|^2)\). Finally use
\(\Delta_H\omega=\nabla^*\nabla\omega+(s/3)\omega-2W^+\omega\). This
proves \[
B_\omega:=\fint|\Delta_H\omega|^2
\le C_L(1+\kappa+b_a)+o(1),
\] \[
\boxed{(1+B_\omega)d_\beta
\le C_L\kappa(1+\kappa+b_a)+o(1).}
\tag{19}
\] All the forms are smooth actual solutions. None of these estimates
differentiates the smoothing weight.

\hypertarget{spectral-compatibility-and-completion-of-the-proof}{%
\subsection{6. Spectral compatibility and completion of the
proof}\label{spectral-compatibility-and-completion-of-the-proof}}

The bounded self-dual-form compatibility theorem states that, for
\(\|\omega\|_\infty\le K\) and bounded constant \(t\), \[
|t|\,|1-2t|
\le C\sqrt{d_t(1+B_\omega+d_t)},
\tag{20}
\] where \(d_t\) is as in (17). Its proof is recorded in full in
\href{round21_analytic.md}{the analytic criterion}, Sections 2, 3 and 6.
The mechanism relevant to its use here is worth making explicit. The
same Hessian estimate bounds second derivatives of \(A_{\omega,t}\). The
weak contracted Bianchi identity, tested against
\(J_t=\operatorname{div}W(A_{\omega,t})\), gives \[
\fint|J_t|^2\le C|t|\sqrt{d_t(1+B_\omega+d_t)}.
\] If \(v=|\omega|^2\), the polynomial pointwise contraction gives \[
t^2v(|\nabla\omega|^2+|d\omega|^2)
\le |J_t|^2+\tfrac43t^2|dv|^2,
\quad
\fint|dv|^2\le C\sqrt{d_t(1+B_\omega+d_t)}.
\] Testing the self-dual Hodge equation against \(v\omega\) gives an
algebraic term \(16t^2(1-2t)/3\); all its derivative terms are
controlled by the preceding two displays. The zeroth-order errors cost
at most \(Ct^2\sqrt{d_t}\). Thus \[
\left|\frac{16}{3}t^2(1-2t)\right|
\le7\fint|J_t|^2+C t^2\sqrt{d_t(1+B_\omega+d_t)}.
\] Division by \(|t|\), with the case \(t=0\) separate, proves (20). The
argument neither differentiates the Ricci error nor assumes
\(d\omega=0\).

Substitute the produced form, with \(t=1/\beta\), into (20). Using
(17)--(19), then letting its smoothing allowance tend to zero, \[
\frac{(\beta-2)^2}{\beta^4}
\le C_L\kappa(1+\kappa+b_a).
\tag{21}
\] Absorb \(\beta^4\le L^4\) into the constant. When \(\kappa\le1\), the
right side is bounded by \(2C_L(1+b_a)\kappa\). When \(\kappa>1\), (2)
is immediate after choosing \(c_L\le(L-2)^{-2}\). This proves (2).

For (3), Chern--Weil gives \[
4\pi^2T=\int\left(2|W^+|^2+s^2/24
-|\operatorname{Ric}^{\circ}|^2/2\right)
\ge-\tfrac12 Vr.
\] Hence \(a^2\ge T\) implies \(b_a\le r/(8\pi^2)\). For \(\kappa\le1\)
this is bounded by \(1/(8\pi^2)\); for \(\kappa>1\) use the same trivial
split. Equation (4) is (3) with \(r=0\), after multiplying by
\(\beta V/(\beta+1)\).

\hypertarget{the-allowed-endpoint-without-a-scalar-curvature-band}{%
\subsection{7. The allowed endpoint without a scalar-curvature
band}\label{the-allowed-endpoint-without-a-scalar-curvature-band}}

The same construction has a structural consequence even when the
parameter approaches the allowed endpoint. It is not obtained by
declaring that endpoint forbidden.

Let \(H\) be the harmonic self-dual representative of \(a_g^+\), with
convention \(\operatorname{proj}_{\mathcal H^+}(iF_A^+)=2\pi H\). Set \[
\zeta_\beta=-\frac{3\pi\beta}{\beta+1}H.
\] \textbf{Theorem 2 (actual harmonic alignment).} Under the hypotheses
of Theorem 1, if \(\kappa\le1\), the same smooth forms can be chosen
with arbitrarily small allowance \(\eta>0\) so that \[
\fint|\omega-\zeta_\beta|^2\le C_L\kappa+\eta,
\quad
\fint|\nabla\omega|^2
\le C_L\sqrt{\kappa(1+\kappa+b_a)}+\eta.
\tag{22}
\] The amplitude, polynomial Weyl fit and Hodge-Laplacian bounds remain
(15), (17) and (19). In particular, at \(\beta=2\), the harmonic target
is \(-2\pi H\), and \(\kappa=r+144\pi^2D_{\rm LB}/V\). This conclusion
needs neither a scalar band nor a scalar spectral gap.

\textbf{Proof.} Retain the actual curvature moment from Section 3. Since
\(C\le E_Q\), its Hodge remainder satisfies \[
18\|iF_A^+-2\pi H\|_2^2
=C_{\rm act}-72\pi^2p\le S+o(1).
\] The curvature equation therefore gives \[
\|f_j\psi+4\sqrt2\pi H\|_2^2\le4S/9+o(1).
\tag{23}
\] Furthermore \[
\int(f_j-\sqrt{q_\beta})^2
\le2\int(f_j-\sqrt Q)^2
+2\int(\sqrt Q-\sqrt{q_\beta})^2
\le C_LV\kappa+o(1).
\] Here \((f_0-\sqrt Q)^2\le f_0^2-Q\le9e^2/16\), and the square-root
map is uniformly Lipschitz on \([3,\infty)\). Use \(|\psi|\le4\) and
multiply (23) by \(c_\beta^2/q_\beta\). The identity
\(4\sqrt2\pi c_\beta/\sqrt{q_\beta} =3\pi\beta/(\beta+1)\) proves the
first estimate in (22).

For the second, the derivative conclusion of the bounded-form theorem on
\(1/L\le t\le1/2\) is
\(\fint|\nabla\omega|^2\le C_L\sqrt{d_t(1+B_\omega+d_t)}\) when
\(d_t\le1\). If \(\kappa\) is small, (17) ensures this. For the
remaining compact range of \(\kappa\le1\), the integrated Hodge equation
and the bound on \(B_\omega\) give
\(\fint|\nabla\omega|^2\le C_L\sqrt{1+b_a}\); enlarging the constant
recovers (22). All allowances are taken small simultaneously for the
same solution.

The harmonic form need not have a uniform maximum bound. Thus (22) does
not license replacing \(\omega\) by \(\zeta_\beta\) inside the quadratic
Weyl model. It does give a precise global spectral conclusion: if
\(\Pi_{>\lambda}\) is the spectral projection of the rough Laplacian on
self-dual forms, then for every \(\lambda>0\), \[
\fint|\Pi_{>\lambda}\zeta_\beta|^2
\le C_L\left[\kappa+
\frac{\sqrt{\kappa(1+\kappa+b_a)}}{\lambda}\right].
\tag{24}
\] Indeed, project \(\zeta_\beta-\omega\) and use
\(\|\Pi_{>\lambda}\omega\|_2^2 \le\lambda^{-1}\|\nabla\omega\|_2^2\),
then let the allowance tend to zero. This controls an actual specified
harmonic class; it does not assert global Kähler topology or exact
metric correction.

At \(\beta=2\), there is a stronger direct estimate, independent of
\(b_a\). In the notation of (11), the retained Hodge identity is \[
\int|\nabla\psi|^2
=2\int|d\psi|^2+2\int E_\lambda
-\tfrac13\int(s+2w)\rho^2.
\] Since \(\rho\le4\), \(\int(s+2w)^2\le2\mathfrak D_2\), and
\(c_2^2=3/4\), this gives simultaneously \[
\fint|\nabla\omega|^2\le G(\kappa)+\eta,
\quad G(\kappa)=13\kappa+4\sqrt{2\kappa},
\quad
\fint|d\omega|^2\le13\kappa/2+\eta.
\tag{25}
\] The preceding harmonic-alignment estimate can likewise be made
explicit: \[
\fint|\omega+2\pi H|^2\le49\kappa+\eta.
\tag{26}
\] For (26), (6) gives \(\fint(q-6)^2\le18\kappa\). The cutoff estimate
implies \(\fint(f_j-\sqrt6)^2\le9\kappa+o(1)\), and (23) then costs at
most \(4\cdot9\kappa+S/(9V)\le49\kappa\). Thus in (24) the endpoint
right side can be replaced by \(98\kappa+2G(\kappa)/\lambda\), with no
square-cost hypothesis. The
\href{line1_composed_application.md}{endpoint application} gives this
statement and its spectral consequences separately.

\hypertarget{what-is-strengthened-and-what-is-not}{%
\subsection{8. What is strengthened, and what is
not}\label{what-is-strengthened-and-what-is-not}}

At \(\beta=2\), \[
\mathfrak D_2=144\pi^2D_{\rm LB},\qquad
D_{\rm LB}=\frac1{4\pi^2}\int(s^2/24+|W^+|^2/2)-p/2.
\] The classical equality set is negative cscK, compatible with the
selected Spin\(^c\) structure, not necessarily Einstein. Equation (2)
does not turn its zero deficit into an unconditional Ricci estimate.

At \(\beta=8\), Chern--Weil gives \[
\mathfrak D_8
=108\pi^2(T-2p/3)+\frac{27}{2}\int|\operatorname{Ric}^{\circ}|^2.
\tag{27}
\] Thus the old topological calibration is one member of the family. For
other parameters, (2) controls a different scalar--Weyl functional, with
its own genuine angular calibration supplied in Section 4. This is a
functional extension of the method, not a second discovery of the old
\(1/3\) blow-up obstruction.

An exact negative Kähler--Einstein metric with canonical class has
\(s=-4,w=2,p=V/(2\pi^2)\). Direct substitution yields \[
\frac{\mathfrak D_\beta}{V}
=\frac{4(\beta-2)^2}{\beta}.
\tag{28}
\] The quadratic vanishing at the allowed endpoint is consequently the
correct order; the theorem does not claim an optimal coefficient.

For \(0<\beta\le2\), the classical scalar and scalar--Weyl bounds
already provide a stronger elementary bound. If \(x=\|s\|_2,y=\|w\|_2\),
then \(x\ge\sqrt{32\pi^2p}\) and \(x+y\ge\sqrt{72\pi^2p}\). Minimizing
\(x^2+\beta y^2\) on these constraints gives
\(\int(s^2+\beta w^2) \ge8\pi^2p(4+\beta)\). Thus this range does not
need the new compatibility argument. For \(\beta>2\), that scalar
constraint is inactive at the classical relaxed minimizer. The
additional obstruction in (2) is differential, through weak Bianchi and
the actual bounded form.

For the Yau26 negative input, \(r\le4\delta_-(g)^2\). Under
\(a^2\ge T\), (3) therefore gives the usable estimate \[
4\delta_-(g)^2+\mathfrak D_\beta/V
\ge c'_L(\beta-2)^2.
\tag{29}
\] To obtain a new purely topological application one must additionally
force \(\mathfrak D_\beta/V\) small from independently justified
topological or geometric data. No such new topology is asserted here.
Neither (29) nor the allowed-root form supply proves existence of an
Einstein metric on an arbitrary smooth manifold.

\hypertarget{sources-and-verification}{%
\subsection{Sources and verification}\label{sources-and-verification}}

\begin{itemize}
\tightlist
\item
  C. LeBrun, \emph{Ricci Curvature, Minimal Volumes, and Seiberg--Witten
  Theory}, \href{https://arxiv.org/abs/math/0003068}{Theorem 2.3 and
  Propositions 3.1, 4.1}: the classical curvature inequalities and cscK
  equality set.
\item
  C. LeBrun, \emph{Four-manifolds, curvature bounds, and convex
  geometry}, \href{https://arxiv.org/abs/math/0611450}{Lemma 3.7}:
  weighted equations.
\item
  \href{round4_rigidity.md}{Integral truncated-weight proof}, Sections
  2--4: the inherited cutoff and moment estimates, reproduced above.
\item
  \href{round21_analytic.md}{Bounded-form spectral theorem}: the
  analytic compatibility theorem, including its polynomial contractions.
\item
  \href{line1_scalar_band_remainder.md}{Scalar-band supply}: the prior
  allowed cscK branch, which is not itself a forbidden-spectrum
  application.
\end{itemize}

Exact Einstein saturation forcing Kähler--Einstein structure already
appears in LeBrun's Corollary 2.5; it is not a new classification here.
The quantitative persistence under integral Ricci error, and the actual
all-scalar form-production estimates, are the assertions at issue.

The parameterized all-scalar functional estimate (2) is the claim proved
in this note. Its finite algebra is checked separately; those checks do
not replace the gauge-solvability hypothesis or the analytic proofs. No
comprehensive priority claim over later literature is made.

END COMPLETE SOURCE

\begin{center}\rule{0.5\linewidth}{0.5pt}\end{center}

\hypertarget{source-4-line1_compactified_pencil.md}{%
\subsection{Source 4:
\nolinkurl{line1_compactified_pencil.md}}\label{source-4-line1_compactified_pencil.md}}

Original path:
\nolinkurl{mathematics/current/line1_compactified_pencil.md}

SHA-256:
\nolinkurl{0d971e1519c9d8b5edddcbd71f158c3ca08acf943d5ae4a7eb3f359af2628bed}

BEGIN COMPLETE SOURCE

\hypertarget{a-compactified-functional-pencil-and-a-two-sided-einstein-energy-window}{%
\section{A compactified functional pencil and a two-sided Einstein
energy
window}\label{a-compactified-functional-pencil-and-a-two-sided-einstein-energy-window}}

23 September 2026. Baseline: \texttt{4401d28}.

This note extends the finite-parameter non-saturation theorem to its
vanishing-Weyl endpoint, with an explicit degenerating weight. It then
obtains the upper, as well as the lower, quantitative energy envelope.
The analytic input is the existing bounded-form compatibility theorem;
the classical feasible energy interval is not a new curvature
inequality.

\hypertarget{uniformity-on-the-unbounded-parameter-interval}{%
\subsection{1. Uniformity on the unbounded parameter
interval}\label{uniformity-on-the-unbounded-parameter-interval}}

Let \((M,g)\) be closed, connected, oriented and four-dimensional. Fix
an actual Spin\(^c\) structure whose ordinary Seiberg--Witten equations
are solvable for every metric conformal to \(g\). Write \[
 p=(a_g^+)^2>0,\quad V=\operatorname{Vol}(g),\quad
 r=\fint|\operatorname{Ric}+g|^2,\quad b_a=(-a^2)_+/V,
\] \[
 D_\beta=\frac{\beta+1}{\beta}\int(s^2+6\beta|W^+|^2)
                 -72\pi^2p,\qquad \beta\ge2.
\] The conventions are tensor Ricci norm and Hilbert--Schmidt Weyl
endomorphism norm. The class is fixed while the conformal weight varies.

\textbf{Theorem 1.} There is a universal \(c>0\) such that, for every
finite \(\beta\ge2\), \[
 \boxed{(1+b_a)(r+D_\beta/V)
       \ge c\frac{(\beta-2)^2}{\beta^4}.}                 \tag{1}
\] If \(a^2\ge T=2\chi+3\tau\), the factor \(1+b_a\) can be removed,
after reducing the universal constant. Unlike the earlier coefficient
\(c_L\), the constant in (1) does not depend on a parameter cutoff. Its
explicit weight vanishes at both the Kähler endpoint and the
vanishing-Weyl endpoint.

\textbf{Proof.} Put \(w=\sqrt6|W^+|\), \(q=-s+w\), \(e_s=s+4\),
\(z=s+\beta w\), and \(\kappa=r+D_\beta/V\). The exact nonnegative
decomposition in \href{line1_functional_pencil.md}{the finite-parameter
proof}, Section 2, implies \[
 \fint e_s^2\le4r,\quad \fint z^2\le\beta D_\beta/V,
 \quad\Delta_+/V\le D_\beta/V.
\] Consequently \[
 \fint(w-4/\beta)^2
 \le2D_\beta/(\beta V)+8r/\beta^2\le3\kappa,
\] \[
 \fint(q-q_\beta)^2\le14\kappa,
 \qquad q_\beta=4+4/\beta\in[4,6].                       \tag{2}
\] In particular the normalized Ricci, scalar and \(W^+\) squared
energies are bounded by \(C(1+\kappa)\), uniformly for \(\beta\ge2\).

Here is the required tracking of the form-production constants;
uniformity is not inferred merely by sending \(L\) to infinity. The
truncated weight \(Q=\max(q,3)\), \(f_0=\max(\sqrt Q,(-s)_+/4)\) is
independent of \(\beta\). The inherited moment argument gives the same
bounds \(|\psi|\le4\), \(D_0\le4S/27+o(1)\) and the same amplitude
control, with \(S\le117V\kappa\). Its errors involve only \(e_s\),
\(\Delta_+\) and \(Q\). For the normalized form and Weyl model the
coefficients are \[
 c_\beta^2=\frac{9\beta}{8(\beta+1)}\in[3/4,9/8],\quad
 \alpha_\beta=\frac9{8(\beta+1)}\le3/8,\quad
 t=1/\beta\in(0,1/2].
\] Set \(\omega=c_\beta\psi\) and
\(A_{\omega,t}=t(|\omega|^2I/3-\omega\otimes\omega)\). The polynomial
fit identity is \[
 |W^+-A_{\omega,t}|^2
 =\tfrac16(w-2\alpha_\beta|\psi|^2)^2
                                  +2\alpha_\beta E_\lambda.
\] The amplitude estimate, (2), and
\(2\alpha_\beta|\psi|^2=u^2/(\beta+1)\) therefore yield, with one smooth
actual solution and an arbitrarily small allowance, \[
 \|\omega\|_\infty\le3\sqrt2,\quad
 d_t:=\fint\bigl(|\operatorname{Ric}+g|^2
   +|W^+-A_{\omega,t}|^2+(|\omega|^2-2)^2\bigr)
                       \le C\kappa+o(1).                 \tag{3}
\] No inverse power of \(t\) is used here. The determinant energy
identity and the bounded-spinor Hessian estimate then give \[
 B_\omega:=\fint|\Delta_H\omega|^2
                   \le C(1+\kappa+b_a)+o(1),             \tag{4}
\] again uniformly: the needed curvature energies are bounded by (2),
and the normalization of the quadratic spinor map is bounded above. This
is exactly the estimate proved in Sections 3--5 of the finite-parameter
paper, with every parameter-dependent coefficient now bounded
explicitly.

The bounded-form compatibility theorem, with the fixed parameter range
\(0\le t\le1/2\), gives \[
 t^2(1-2t)^2\le C\kappa(1+\kappa+b_a).                   \tag{5}
\] For \(\kappa\le1\) this proves (1); for \(\kappa>1\) use
\(t^2(1-2t)^2\le1/64\). Let the smoothing allowance tend to zero. If
\(a^2\ge T\), Chern--Weil gives \(b_a\le r/(8\pi^2)\), and the same
split removes the square-cost factor. This proves the theorem. No
derivative estimate for the harmonic projection is asserted uniformly as
\(t\to0\). \(\square\)

\hypertarget{a-pencil-with-an-actual-endpoint-at-zero}{%
\subsection{2. A pencil with an actual endpoint at
zero}\label{a-pencil-with-an-actual-endpoint-at-zero}}

Define on \(0\le t\le1/2\) \[
 P_t=t(1+t)\int s^2+6(1+t)\int|W^+|^2-72\pi^2tp.
                                                               \tag{6}
\] For \(t>0\), this is \(tD_{1/t}\); at zero it is exactly
\(P_0=6\int|W^+|^2\). Hence \(P_t\ge0\) on the closed interval.
Multiplying (1) by \(t\), and using \(r+tD/V\ge t(r+D/V)\), proves \[
 \boxed{(1+b_a)(r+P_t/V)\ge c t^3(1-2t)^2.}             \tag{7}
\] The same square-cost removal applies when \(a^2\ge T\). At \(t=0\)
the inequality is trivial and supplies no nonzero lower amplitude or
actual anti-self-dual Einstein existence. At \(t=1/2\), the endpoint
includes negative Kähler--Einstein metrics. The cubic decay in (7) is a
proved sufficient rate from this transfer, not an optimality claim.

For a normalized negative Kähler--Einstein metric and its canonical
class, \[
 P_t/V=4(1-2t)^2.
\] This provides an actual normalization check. It does not establish
sharpness at the other endpoint.

\hypertarget{the-missing-upper-energy-envelope}{%
\subsection{3. The missing upper energy
envelope}\label{the-missing-upper-energy-envelope}}

Assume now that \(A>0\), \[
 T=2\chi+3\tau=A(1-h),\qquad 1/4<h<1/3,                 \tag{8}
\] and for every metric one can choose an actual conformally solvable
class with \(p\ge A\) and \(a^2\ge T\). Put \[
 e=\frac{\int|\operatorname{Ric}|^2}{8\pi^2A},\quad
 F(h)=\frac{(4-\sqrt{1-3h})^2}{9},\quad
 G(h)=\frac{(4+\sqrt{1-3h})^2}{9},
\] \[
 t_*(h)=\frac{1-2\sqrt{1-3h}}{8+2\sqrt{1-3h}}\in(0,1/8),
 \qquad \Theta(h)=t_*^3(1-2t_*)^2.                       \tag{9}
\] Use the all-sign relative defect \[
 \mathcal E(g)=\inf_{\lambda\ne0}
 \frac{\int|\operatorname{Ric}-\lambda g|^2}{\lambda^2V}.
\]

\textbf{Theorem 2.} There is a universal \(c_*>0\) such that every
smooth metric under (8) satisfies \[
 \boxed{\mathcal E(g)+(G(h)-e)_+\ge c_*\Theta(h).}        \tag{10}
\] Combining with the inherited lower-envelope theorem, every Einstein
candidate must satisfy the strictly smaller energy window \[
 \boxed{F(h)+\varepsilon h^2\le e
                       \le G(h)-c_*\Theta(h).}          \tag{11}
\] Ricci-flat and positive Einstein metrics are independently excluded
by the positive actual monopole scalar estimate. Equation (11) concerns
all remaining Einstein metrics, not just minimizers or a gluing
neighborhood.

\textbf{Proof.} Write \(R=\int|\operatorname{Ric}|^2\),
\(P=\int|\operatorname{Ric}^{\circ}|^2\), \(C_A=8\pi^2A\). The exact
Chern--Weil reduction of (6) is \[
 P_t=k(t)R+\ell(t)P+12\pi^2(1+t)T-72\pi^2tp,
\] \[
 k(t)=\frac{(1+t)(8t-1)}2,\qquad
 \ell(t)=2(1+t)(1-2t).                                  \tag{12}
\] For \(t=t_*(h)\), direct substitution gives \[
 k(t)G(h)+\tfrac32(1+t)(1-h)-9t=0,
\] and therefore \[
 P_t=C_A k(t)(e-G(h))+\ell(t)P-72\pi^2t(p-A)
       \le C_A|k(t)|(G(h)-e)_++\ell(t)P.                \tag{13}
\] This is the opposite sign of the coefficient in the lower-envelope
argument. It is why the finite-parameter range beyond \(\beta=8\)
controls an upper energy wall.

First fix the target to \(\operatorname{Ric}=-g\) and
\(r=\fint|\operatorname{Ric}+g|^2\le1\). The scalar monopole bound and
the trace estimate imply \[
 C_A/V\le(2+\sqrt r)^2\le9,\qquad P/V\le r.
\] On \(0\le t\le1/8\), \(|k(t)|\le1/2\) and \(0\le\ell(t)\le2\). Using
(7), with its square-cost factor removed, in (13) gives \[
 c\Theta(h)\le3r+\tfrac92(G(h)-e)_+
              \le\tfrac92\{r+(G(h)-e)_+\}.             \tag{14}
\] For \(r>1\) the desired estimate follows by decreasing the constant.
Constant rescaling and taking the infimum over negative target constants
prove (10) for the negative-target infimum \(\mathcal E_-\).

The passage to all signs needs a different argument from the earlier
lower-envelope transfer: a very large \(e\) now makes the energy excess
zero. Here the upper Chern--Weil budget closes that possibility. Let
\(I=\int s\). If \(I\le0\), direct minimization gives
\(\mathcal E_+=4\). If \(I>0\), it gives
\(\mathcal E_+=\mathcal E=4-I^2/(RV)\). The signed scalar estimate
\(\int s_-^2\ge4C_A\) implies \[
 \mathcal E_+\ge4/e.                                   \tag{15}
\] If \(\mathcal E_+<1/2\), the chirality upper bound gives \[
 e(1-\mathcal E_+)\le
 U=3\min(2\chi+3\tau,2\chi-3\tau)/A\le3,
\] so \(e\le6\). Equation (15) would imply \(\mathcal E_+\ge2/3\), a
contradiction. Thus \(\mathcal E_+\ge1/2\). Since \(\Theta\) is bounded,
this proves (10) also for the positive target infimum, after reducing
\(c_*\); take the minimum of the two signwise bounds. \(\square\)

\hypertarget{geometry-of-the-two-endpoints-and-an-actual-application}{%
\subsection{4. Geometry of the two endpoints and an actual
application}\label{geometry-of-the-two-endpoints-and-an-actual-application}}

The classical Einstein constraints are, with nonnegative \(x,y\), \[
 x^2=4e,\quad x^2+8y^2=12(1-h),\quad x\ge2,\quad x+y\ge3.
\] When \(1/4<h<1/3\), the ellipse has \(x<3\). Solving these
inequalities gives the full classical interval \(F(h)\le e\le G(h)\).
Thus the interval itself is classical; the positive quantitative margins
in (11) are the new application of the differential non-saturation
mechanism.

At \(h\downarrow1/4\), the upper optimal parameter tends to \(t=0\), and
\(\Theta(h)\sim(8/27)(h-1/4)^3\). At \(h\uparrow1/3\), both optimal
parameters meet \(t=1/8\); the interval collapses, consistently with the
old critical gap. At the lower-envelope endpoint \(h\downarrow0\), the
optimal parameter instead tends to the allowed Kähler value \(t=1/2\).
These are different geometric degeneracies; they cannot be treated as
one compact parameter range with a fixed forbidden distance.

For an actual example, let \[
 M=(\Sigma_2\times\Sigma_3)\#5\overline{\mathbb{CP}}^{\,2}.
\] Its polarized actual classes give \(A=16\), \(h=5/16\), \(a^2=T=11\).
Here \[
 F=25/16,\qquad G=289/144,\qquad t_*=1/17,
 \qquad \Theta=225/17^5.
\] Any Einstein metric, if one exists, must consequently satisfy \[
 \frac{25}{16}+\varepsilon\frac{25}{256}
 \le\frac{R}{128\pi^2}
 \le\frac{289}{144}-c_*\frac{225}{17^5}.                 \tag{16}
\] This is not a nonexistence claim for this example. Its bare
two-chirality upper bound is only \(e\le33/16\), so the upper envelope
contains information not supplied by that bound alone. The actual
monopole selection is the same established polarized input as in the
main paper; no new gauge existence theorem is asserted here.

\hypertarget{relation-to-the-research-goal}{%
\subsection{5. Relation to the research
goal}\label{relation-to-the-research-goal}}

The first-line theorem now supplies a two-sided restriction on the
entire energy spectrum of possible Einstein outputs in the displayed
regime. It does not merely study an energy-minimizing sequence. For a
pointwise input, \(\mathcal E\le4\delta_\sigma^2\) converts (10) into
the corresponding upper-energy constraint on that input; a small
residual alone does not force the input energy into a forbidden region.

This note neither classifies the allowed spectra of all Einstein metrics
nor constructs an Einstein metric. The separate Bauer--Furuta
application in this increment is the new topological-category result.
The present contribution is the closed-parameter formulation and its
upper-wall application, using the existing analytic mechanism.

Sources: \href{https://arxiv.org/abs/math/0003068}{LeBrun's scalar--Weyl
inequalities and Chern--Weil reduction}; the local
\href{line1_functional_pencil.md}{functional-pencil proof}, Sections
2--6; \href{round21_analytic.md}{bounded-form compatibility}, Section 6;
\href{line1_energy_envelope_gap.md}{lower energy envelope}; and
\href{line1_two_chirality_topology.md}{two-chirality upper bound}.

\nolinkurl{check_line1_compactified_pencil.py} verifies the exact
coefficients, both envelope roots, endpoint rates, actual example, and
several deliberate sign errors. These checks do not substitute for the
written uniform analytic estimates or actual gauge-solvability inputs.

END COMPLETE SOURCE

\begin{center}\rule{0.5\linewidth}{0.5pt}\end{center}

\hypertarget{source-5-line1_energy_envelope_gap.md}{%
\subsection{Source 5:
\nolinkurl{line1_energy_envelope_gap.md}}\label{source-5-line1_energy_envelope_gap.md}}

Original path:
\nolinkurl{mathematics/current/line1_energy_envelope_gap.md}

SHA-256:
\nolinkurl{8cd0c429b2f403b1c31bb24333efc7397afba5742394a2951dd56aa5e24c3ad4}

BEGIN COMPLETE SOURCE

\hypertarget{a-uniform-ricci-energy-gap-above-the-classical-einstein-envelope}{%
\section{A uniform Ricci-energy gap above the classical Einstein
envelope}\label{a-uniform-ricci-energy-gap-above-the-classical-einstein-envelope}}

23 September 2026. This is an application of the parameterized
functional theorem, with actual polarized monopole classes. The
elementary energy envelope below follows from classical inequalities.
The added conclusion is a uniform positive separation from that
envelope, persisting under integral Einstein error. No new topological
nonexistence theorem is asserted.

\hypertarget{statement-on-a-realized-family}{%
\subsection{1. Statement on a realized
family}\label{statement-on-a-realized-family}}

Let \[
 M=X\#k\overline{\mathbb {CP}}^{\,2}\#m(S^1\times S^3),\qquad
 A=c_1^2(X)>0,\quad h=\frac{k+4m}{A},\quad 0<h<\frac13,
 \tag{1}
\] where \(X\) is a minimal complex surface of general type. Put \[
 R=\int_M|\operatorname{Ric}|^2,\qquad
 e=\frac{R}{8\pi^2A},\qquad
 F(h)=\frac{(4-\sqrt{1-3h})^2}{9},
\] \[
 \mathcal E(g)=\inf_{\lambda\ne0}
 \frac{\int_M|\operatorname{Ric}-\lambda g|^2}
      {\lambda^2\operatorname{Vol}(g)}.
 \tag{2}
\] Both \(e\) and \(\mathcal E\) are invariant under constant rescaling.

\textbf{Theorem.} Relative to Theorem 1 of
\href{line1_functional_pencil.md}{the functional-pencil proof}, there is
a universal \(\varepsilon>0\), independent of \(X,k,m,h,g\), such that
\[
 \boxed{\mathcal E(g)+(e-F(h))_+\ge\varepsilon h^2.}
 \tag{3}
\] Thus, for example, \(\mathcal E(g)\le\varepsilon h^2/2\) implies
\(e\ge F(h)+\varepsilon h^2/2\). Every Einstein metric on \(M\), if one
exists, necessarily satisfies \[
 \boxed{\frac{R}{8\pi^2A}\ge F(h)+\varepsilon h^2.}
 \tag{4}
\] An actual nonzero monopole projection excludes Einstein constants
zero and positive, so (4) concerns all Einstein candidates, not just a
chosen sign.

The same theorem holds whenever \(T=2\chi+3\tau=A(1-h)\) and each metric
admits an actual conformally solvable Spin\(^c\) structure with
\(p=(a_g^+)^2\ge A\) and \(a^2\ge T\). No regularity of the class
selection in \(g\) is needed. In (1), the polarized sign selection gives
\[
 p\ge A,\qquad a^2=A-k\ge A-k-4m=T.
 \tag{5}
\] This is gauge-theoretic input, not just an intersection-form
calculation; the conformal class keeps its polarization and its selected
structure. For \(b_+=1\), the small-perturbation chamber and, for
handles, the loop insertions are those in the
\href{lebrun_stability_closure.md}{polarized closure}. The source is
\href{https://arxiv.org/pdf/math/0003068}{LeBrun, Proposition 3.2 and
its proof}.

For a concrete family without chamber issues, take
\(X=\Sigma_g\times\Sigma_q\), \(g,q\ge2\). Then \(A=8(g-1)(q-1)\) and
\(b_+>1\). In particular,
\((\Sigma_2\times\Sigma_2)\#2\overline{\mathbb {CP}}^{\,2}\) has
\(A=8,h=1/4\). Formula (3) applies to every smooth metric on this actual
closed manifold. Any Einstein metric on it would satisfy \[
 \frac{R}{64\pi^2}\ge\frac{49}{36}+\frac{\varepsilon}{16}.
\] We do not assume or produce an Einstein metric on it.

\hypertarget{the-classical-envelope-is-not-the-new-result}{%
\subsection{2. The classical envelope is not the new
result}\label{the-classical-envelope-is-not-the-new-result}}

For an Einstein metric, set \[
 x^2=\frac{\int s^2}{8\pi^2A},\qquad
 y^2=\frac{6\int|W^+|^2}{8\pi^2A}.
\] The classical scalar and scalar--Weyl inequalities and Chern--Weil
give \[
 x\ge2,\qquad x+y\ge3,\qquad x^2+8y^2=12(1-h),\qquad e=x^2/4.
 \tag{6}
\] These imply \(x\ge(8-2\sqrt{1-3h})/3\), hence \(e\ge F(h)\). Equality
is feasible in this finite relaxation; it is not claimed to be realized
by a metric. The constraints are the scalar estimate,
\href{https://arxiv.org/pdf/math/0003068}{LeBrun's Theorem 2.3}, and the
exact Chern--Weil identity. This optimization is entirely classical.

The all-metrics infimum in (1) is \(\inf e=1+h\), by
\href{https://arxiv.org/pdf/math/0003068}{LeBrun's Theorem 4.3}. One has
the exact comparison \[
 F(h)-(1+h)=\frac{4h^2}{(1+\sqrt{1-3h})^2}>0.
 \tag{7}
\] The new term in (4) is therefore an additional uniform energy
separation above the already optimized classical restriction, not the
separation in (7). Classical equality rigidity already excludes exact
Einstein saturation away from the Kähler--Einstein endpoint; it does not
alone give a uniform gap for varying manifolds or almost-Einstein
metrics.

\hypertarget{exact-chernweil-reduction-and-the-proof}{%
\subsection{3. Exact Chern--Weil reduction and the
proof}\label{exact-chernweil-reduction-and-the-proof}}

Use tensor Ricci norms and Hilbert--Schmidt Weyl operator norms. Write
\[
 P=\int|\operatorname{Ric}^{\circ}|^2,\quad S_2=\int s^2,\quad
 W_2=\int|W^+|^2.
\] The identities \[
 S_2=4(R-P),\qquad W_2=2\pi^2T-R/12+P/3
\] turn the functional in the companion theorem into \[
 \mathfrak D_\beta
 =K_\beta R+L_\beta P+12\pi^2(\beta+1)T-72\pi^2p,
 \tag{8}
\] where \[
 K_\beta=\frac{(\beta+1)(8-\beta)}{2\beta},\qquad
 L_\beta=\frac{2(\beta+1)(\beta-2)}{\beta}.
\] Let \(z=\sqrt{1-3h}\) and choose \[
 \beta=\beta(h)=\frac{8-2z}{1+2z}\in(2,8).
 \tag{9}
\] Direct substitution gives \[
 K_\beta F(h)+\tfrac32(\beta+1)(1-h)-9=0.
\] Using (5) in (8), with no discarded negative error, yields \[
 \mathfrak D_\beta
 =8\pi^2A K_\beta(e-F(h))+L_\beta P-72\pi^2(p-A)
 \le8\pi^2A K_\beta(e-F(h))+L_\beta P.
 \tag{10}
\]

First fix the Einstein target to \(-g\), and let
\(V=\operatorname{Vol}(g)\), \(r=V^{-1}\int|\operatorname{Ric}+g|^2\).
The functional-pencil theorem, with \(L=8\) and \(a^2\ge T\), supplies a
universal \(c>0\) such that \[
 c(\beta-2)^2\le r+\mathfrak D_\beta/V.
\] Since \(P/V\le r\), equation (10) implies \[
 c(\beta-2)^2
 \le (1+L_\beta)r+
 K_\beta\frac{8\pi^2A}{V}(e-F(h)).
 \tag{11}
\] The scalar monopole bound and the trace inequality give \[
 32\pi^2A\le S_2,\qquad
 \sqrt{S_2/V}\le4+2\sqrt r,
 \qquad \frac{8\pi^2A}{V}\le(2+\sqrt r)^2.
 \tag{12}
\] For \(r\le1\), replace the last factor in (11) by its positive part
before using (12). On \([2,8]\), \[
 0\le K_\beta\le9/2,\qquad 0\le L_\beta\le27/2.
\] Consequently \[
 c(\beta-2)^2\le\frac{81}{2}\bigl[r+(e-F(h))_+\bigr].
 \tag{13}
\] Finally, \[
 \beta(h)-2
 =\frac{18h}{(1+z)(1+2z)}\ge3h.
\] Thus (3), with \(r\) in place of \(\mathcal E\), holds for \(r\le1\)
with coefficient \(2c/9\). For \(r>1\) it is immediate after choosing
\(\varepsilon=\min\{2c/9,1\}\), because \(h^2<1/9\). For any
\(\lambda<0\), apply this result to \((-\lambda)g\): its normalized
error is the corresponding quotient in (2), while \(e\) and \(h\) are
unchanged. Taking the infimum proves (3) first for \(\mathcal E_-\), the
infimum restricted to negative \(\lambda\).

To include positive target constants, put \(C=8\pi^2A\), \(I=\int s\),
and let \(\mathcal E_+\) be the infimum for \(\lambda>0\). The signed
scalar estimate is \[
 \int s_-^2\ge32\pi^2p\ge4C,\qquad s_-:=\max\{-s,0\}.
 \tag{14}
\] For clarity, it follows directly from an ordinary monopole solution:
with \(\rho=|\Phi|^2\), the integrated spinor identity gives
\(\int\rho^2\le\int s_-\rho\), while Hodge projection gives
\(\int\rho^2\ge32\pi^2p\); apply Cauchy--Schwarz. If \(I>0\), then \[
 \frac{I^2}{V}\le\int s_+^2\le S_2-4C\le4R-4C,
 \qquad \mathcal E_+=4-\frac{I^2}{RV}\ge\frac4e.
\] If \(I\le0\), direct minimization gives \(\mathcal E_+=4\). When
\(e\le F(h)+1\le25/9\), therefore, \(\mathcal E_+\ge36/25>1\). When
\(e>F(h)+1\), the energy excess alone exceeds one. Hence
\(\mathcal E_++(e-F(h))_+\ge1\ge\varepsilon h^2\). Together with
\(\mathcal E=\min\{\mathcal E_-,\mathcal E_+\}\), this proves (3) for
all target signs.

Only the functional-pencil theorem is the nonclassical analytic input.
There is no geometric compactness, scalar band, auxiliary-form
assumption, or additional square-cost budget in this transfer.

\hypertarget{why-the-parameter-eight-does-not-give-this-energy-separation}{%
\subsection{4. Why the parameter eight does not give this energy
separation}\label{why-the-parameter-eight-does-not-give-this-energy-separation}}

At \(\beta=8\), \(K_8=0\). Equation (8) loses \(R\) altogether: \[
 \mathfrak D_8=108\pi^2(T-2p/3)+27P/2.
 \tag{15}
\] This is the old topological gap. Varying \(\beta\) retains the
independent energy coordinate and gives the separation in (3).

There is an exact algebraic non-implication check. At \(h=1/4\), set
\(p=A\), \(a^2=T=3A/4\), \(P=0\), \(s=-4\), and prescribe the constant
Weyl norm \(w=\sqrt6|W^+|=8/7\). Choose \[
 V=2\pi^2A(49/36),\qquad R=4V.
\] These moment data satisfy Chern--Weil, the scalar inequality, every
classical bound \(\mathfrak D_\gamma\ge0\), \(\gamma>0\), and
\(e=F(1/4)=49/36\). However, \[
 \mathfrak D_{7/2}=0,\qquad \mathfrak D_8/V=162/49.
 \tag{16}
\] They violate the new parameter-\(7/2\) estimate. They do not violate
any possible universal coefficient in the global parameter-eight member
\(r+\mathfrak D_8/V\ge36c_8\): necessarily \(c_8\le1/11\). Indeed, for
an actual negative Kähler--Einstein metric \(g_0\), the rescalings
\(g=tg_0\) give \[
 r+\mathfrak D_8/V=4(1-1/t)^2+18/t^2
 \ge36/11,
\] with equality at \(t=11/2\). Yet \(162/49>36/11\). This proves
independence from the parameter-eight numerical inequality even at the
largest coefficient it could possibly have.

The data in (16) are a formal curvature-moment model, not a metric
construction. The manifold and its actual monopole-class selection in
Section 1 are realized; realizing the forbidden moments would contradict
the theorem. This argument does not claim independence from every
possible use of the old differential proof machinery.

\hypertarget{the-remaining-topological-input-is-precise}{%
\subsection{5. The remaining topological input is
precise}\label{the-remaining-topological-input-is-precise}}

Equation (3) does not exclude all energies on any new smooth topology.
For normalized Einstein metrics, the missing independent number is \[
 e=\frac{R}{8\pi^2A}=\frac{V}{2\pi^2A}.
 \tag{17}
\] It is not fixed by \(T\), \(A\), or \(p\). Chern--Weil fixes a linear
combination of scalar and Weyl energies, not this scalar/volume
coordinate. An independently proved upper bound
\(e<F(h)+\varepsilon h^2\) for every Einstein candidate would convert
(4) into a new topological obstruction; no such bound is supplied here.
An all-metrics energy \emph{infimum} cannot serve as that upper bound.

For the pointwise negative Ricci input
\(\delta_-=\|\operatorname{Ric}^{\sharp}+I\|_{\infty,\mathrm{op}}\), (3)
implies the raw consequence \[
 4\delta_-^2+(e-F(h))_+\ge\varepsilon h^2.
 \tag{18}
\] Thus a construction with controlled energy below the strengthened
envelope cannot be arbitrarily Ricci-pinched. This is a restriction on
such input metrics, not a proof that their smooth manifold has no
Einstein metric at a different energy.

\nolinkurl{check_line1_energy_envelope_gap.py} checks the exact
identities, envelope optimization, constants, scale transformation, and
formal non-implication witness. These finite checks do not certify
Seiberg--Witten solvability or the inherited differential theorem. No
priority claim over unexamined later literature is made.

END COMPLETE SOURCE

\begin{center}\rule{0.5\linewidth}{0.5pt}\end{center}

\hypertarget{source-6-line1_two_chirality_topology.md}{%
\subsection{Source 6:
\nolinkurl{line1_two_chirality_topology.md}}\label{source-6-line1_two_chirality_topology.md}}

Original path:
\nolinkurl{mathematics/current/line1_two_chirality_topology.md}

SHA-256:
\nolinkurl{91eb14bc07220c6d66e2b97af97501d324bc1dad999ba7b9d2f41368217fa5a6}

BEGIN COMPLETE SOURCE

\hypertarget{from-the-energy-envelope-to-a-two-chirality-topological-gap}{%
\section{From the energy envelope to a two-chirality topological
gap}\label{from-the-energy-envelope-to-a-two-chirality-topological-gap}}

23 September 2026. This note composes the existing functional-pencil
theorem with both Chern--Weil identities. Its new application is a
uniform integral Einstein obstruction on a realized family of mixed
blow-ups and handle additions lying away from the previously recorded
one-third critical band. All statements depend on the analytic theorem
in \href{line1_functional_pencil.md}{the functional-pencil manuscript},
through \href{line1_energy_envelope_gap.md}{the energy-envelope
theorem}. The elementary Chern--Weil calculation is classical. No
independent literature-priority certification is claimed for this
composition.

\hypertarget{an-upper-bound-that-applies-to-every-candidate}{%
\subsection{1. An upper bound that applies to every
candidate}\label{an-upper-bound-that-applies-to-every-candidate}}

Let \[
 M=X\#k\overline{\mathbb{CP}}^{\,2}\#m(S^1\times S^3),\qquad
 A=c_1^2(X)>0,\quad h=\frac{k+4m}{A}\in(0,1/3),
 \tag{1}
\] where \(X\) is a minimal complex surface of general type. Write \[
 T_\pm=2\chi(M)\pm3\tau(M),\qquad
 U=\frac{3\min(T_+,T_-)}{A}
   =\frac{3(2\chi(M)-3|\tau(M)|)}{A},
\] \[
 R=\int_M|\operatorname{Ric}|^2,\quad
 P=\int_M|\operatorname{Ric}^{\circ}|^2,\quad
 e=\frac{R}{8\pi^2A},\quad
 F(h)=\frac{(4-\sqrt{1-3h})^2}{9}.
 \tag{2}
\] The Ricci norms are tensor norms; Weyl norms are Hilbert--Schmidt
endomorphism norms on two-forms. Set \[
 \mathcal E(g)=\inf_{\lambda\ne0}
 \frac{\int|\operatorname{Ric}-\lambda g|^2}
      {\lambda^2\operatorname{Vol}(g)}.
 \tag{3}
\] The scalar monopole estimate makes \(R>0\). Direct minimization, with
\(I=\int s\) and \(V=\operatorname{Vol}(g)\), gives \[
 \mathcal E=4-\frac{I^2}{RV}\ge\frac{4P}{R}.
 \tag{4}
\] For \(I=0\), the value four is an infimum as \(|\lambda|\to\infty\);
the same formula remains valid. The inequality uses
\(I^2/V\le\int s^2=4(R-P)\), not a sign assumption on \(s\).

The two Chern--Weil formulas give, separately for either sign, \[
 \int|W^\pm|^2=2\pi^2T_\pm-\frac{R}{12}+\frac P3\ge0.
\] Consequently \[
 e\le U+\frac{4P}{8\pi^2A}\le U+e\mathcal E,
 \qquad\boxed{e(1-\mathcal E)\le U.}
 \tag{5}
\] For Einstein metrics, this is the genuine upper bound \(e\le U\). It
controls all candidates, regardless of their proximity to a preferred
construction or variational minimum. No energy infimum has been used as
an upper bound. No monopole structure in the opposite orientation is
assumed: only the nonnegativity of its Weyl energy is used.

\hypertarget{a-uniform-topological-inequality-for-all-smooth-metrics}{%
\subsection{2. A uniform topological inequality for all smooth
metrics}\label{a-uniform-topological-inequality-for-all-smooth-metrics}}

\textbf{Theorem 1.} There is a universal \(c>0\) such that every metric
on every manifold (1) satisfies \[
 \boxed{\mathcal E(g)+(U-F(h))_+\ge c h^2.}
 \tag{6}
\] One may take \(c=\varepsilon/7\), where \(0<\varepsilon\le1\) is the
constant in the energy-envelope theorem \[
 \mathcal E(g)+(e-F(h))_+\ge\varepsilon h^2.
 \tag{7}
\] The same argument works in the abstract actual-class setting of that
theorem: \(T_+=A(1-h)\), and for every metric an actual conformally
solvable Spin\(^c\) structure has \((a_g^+)^2\ge A\) and
\(a_g^2\ge T_+\).

\textbf{Proof.} Since \(T_+/A=1-h\), one has \(U\le3\). If
\(\mathcal E\ge1/2\), (6) is immediate because \(c h^2\le1/63\). If
\(\mathcal E<1/2\), (5) forces \(U>0\) and \[
 e\le\frac{U}{1-\mathcal E}
 \le U+2U\mathcal E\le U+6\mathcal E.
\] Hence \[
 (e-F(h))_+\le(U-F(h))_++6\mathcal E.
\] Substitute this into (7) and use
\(7\mathcal E+(U-F)_+\le7[\mathcal E+(U-F)_+]\). This proves (6),
including \(U\le0\) through the first case. \(\square\)

\textbf{Corollary 2.} If \[
 U\le F(h)+\frac c2 h^2,
 \tag{8}
\] then every smooth metric satisfies \[
 \mathcal E(g)\ge\frac c2 h^2.
 \tag{9}
\] There are no Einstein metrics of any sign on this manifold. Nonzero
Einstein constants would give \(\mathcal E=0\). For the zero constant,
the signed scalar monopole inequality \(\int s_-^2\ge32\pi^2A>0\)
excludes Ricci-flatness separately. For either normalized pointwise
branch, \[
 \delta_\sigma(g)=\|g^{-1}\operatorname{Ric}-\sigma I\|_{\infty,\mathrm{op}},
 \qquad\sigma\in\{-1,1\},
\] one has \(\mathcal E\le4\delta_\sigma^2\). Thus \[
 \boxed{\delta_\sigma(g)\ge h\sqrt{c/8}.}
 \tag{10}
\] This is an obstruction to both pointwise inputs on these types, not a
Yau26 counterexample.

The previously obtained family-uniform theorem used only \(T_+/A=1-h\).
Formula (6) adds the independent opposite-chirality budget \(T_-\). When
\(T_-\) is large it yields no new restriction; the mixed examples below
have a small enough \(T_-\) for it to matter.

\hypertarget{realized-mixed-ball-quotient-examples-outside-the-old-band}{%
\subsection{3. Realized mixed ball-quotient examples outside the old
band}\label{realized-mixed-ball-quotient-examples-outside-the-old-band}}

Choose any compact smooth complex-hyperbolic surface \(X_0\), with
complex orientation, and a sequence of connected finite unramified
covers \(X_j\to X_0\) of unbounded degree. Such a sequence exists: the
cocompact lattice is an infinite finitely generated linear group in
characteristic zero (use its faithful adjoint representation), hence is
residually finite. A finitely generated residually finite infinite group
has finite-index subgroups of unbounded index. Otherwise it has only
finitely many finite-index subgroups, their intersection is
finite-index, and residual finiteness makes that intersection trivial,
forcing the group to be finite. All covers remain compact complex
hyperbolic and minimal of general type.

Put \(A_j=c_1^2(X_j)\to\infty\). The Chern numbers are \[
 \chi(X_j)=A_j/3,\qquad \tau(X_j)=A_j/9.
\] Define integers \[
 k_j=\left\lfloor\frac{29 A_j}{2025}\right\rfloor,
 \qquad
 m_j=\left\lfloor\frac{397 A_j}{32400}\right\rfloor-2,
 \qquad
 M_j=X_j\#k_j\overline{\mathbb{CP}}^{\,2}
          \#m_j(S^1\times S^3).
 \tag{11}
\] For sufficiently large \(j\), both integers are positive. No
prescribed divisibility of the cover degree or \(A_j\) is needed.

Let \[
 h_0=\frac{19}{300},\quad x_0=\frac{29}{2025},\quad
 z_0=\frac{397}{32400},\quad F_0=F(h_0)=\frac{961}{900}.
\] Then \(x_0+4z_0=h_0\), and \[
 1+15x_0-12z_0=F_0.
 \tag{12}
\] Write \(k_j=x_0A_j-\alpha_j\) and \(m_j=z_0A_j-\beta_j-2\), where
\(0\le\alpha_j,\beta_j<1\). Consequently \[
 h_j=h_0-\frac{\alpha_j+4\beta_j+8}{A_j},
\] \[
 \frac{3T_-(M_j)}{A_j}
 =1+\frac{15k_j-12m_j}{A_j}
 =F_0+\frac{-15\alpha_j+12\beta_j+24}{A_j}.
 \tag{13}
\] The signature is positive, since
\(\tau(M_j)/A_j=1/9-k_j/A_j\ge1/9-x_0>0\). Thus the second expression in
(13) is exactly \(U_j\). Moreover \[
 F'(h)=\frac{4-\sqrt{1-3h}}{3\sqrt{1-3h}},
 \qquad 0<F'(h)\le\frac{31}{27}<2
 \quad(0\le h\le h_0).
\] For large \(j\), therefore, \[
 0<h_j<h_0,\qquad
 \boxed{0<U_j-F(h_j)<\frac{62}{A_j}.}
 \tag{14}
\] Indeed, the second correction in (13) lies strictly between \(9/A_j\)
and \(36/A_j\), and \(0<F_0-F(h_j)<26/A_j\).

Since \(h_j\to h_0>0\), (14) eventually implies (8). We have proved:

\textbf{Theorem 3.} The realized sequence (11), for every sufficiently
large \(j\), satisfies \[
 \boxed{\inf_{g\in\operatorname{Met}(M_j)}\mathcal E(g)
       \ge\frac{c h_0^2}{8}>0.}
 \tag{15}
\] In particular, none of these smooth manifolds carries any Einstein
metric. For every smooth metric on every sufficiently large member, \[
 \delta_\sigma(g)\ge h_0\sqrt{c/32},\qquad \sigma=\pm1.
 \tag{16}
\] The threshold index depends on the universal analytic constant and
the chosen cover sequence; no explicit numerical lower bound for that
constant is claimed. The manifolds have unbounded Euler characteristic,
so an infinite subsequence consists of pairwise distinct smooth types.

\hypertarget{actual-gauge-input-not-merely-numerical-chern-data}{%
\subsubsection{Actual gauge input, not merely numerical Chern
data}\label{actual-gauge-input-not-merely-numerical-chern-data}}

For every metric on (11), select the exceptional signs at its
polarization as in Section 2 of \href{lebrun_stability_closure.md}{the
polarized closure proof}. The actual structure has \[
 (a_g^+)^2\ge A_j,\qquad a_g^2=A_j-k_j\ge A_j-k_j-4m_j=T_+(M_j).
\] Nonvanishing after handle addition includes the \(m_j\) loop
insertions. The structure is fixed throughout the conformal class; the
unchanged Hodge star on two-forms supplies the same class for every
auxiliary positive weight. This is precisely the quantifier
\(\forall g\,\exists\mathfrak s_g\,\forall f>0\) used in the
energy-envelope theorem. The classical input is LeBrun's Proposition 3.2
and its polarized proof, not residual finiteness or the intersection
form alone. Alternatively, sufficiently high covers have \(b_+>1\),
since \(\chi(\mathcal O_{X_j})=A_j/9\to\infty\) and \(b_+=1+2p_g\);
hence these examples do not require the small-\(b_+\) chamber extension.

\hypertarget{precisely-which-prior-tests-these-examples-pass}{%
\subsection{4. Precisely which prior tests these examples
pass}\label{precisely-which-prior-tests-these-examples-pass}}

The following comparisons establish additional coverage relative to the
recorded first-line bounds and to the displayed classical tests. They do
not prove that every Einstein-obstruction paper has been exhaustively
searched.

\begin{enumerate}
\def\labelenumi{\arabic{enumi}.}
\tightlist
\item
  \textbf{The recorded uniform one-third band does not detect them.} Its
  stated constant is
  \(\varepsilon_0=\min\{1/4,c_{\rm cal}/1215\}\le1/4\), whereas
  \(1/3-h_j>1/3-h_0=27/100>1/4\). Thus setting \(\mathcal E=0\) does not
  contradict that inequality. In particular this is not another choice
  of topology inside its stated near-one-third strip.
\item
  \textbf{Both strict Hitchin--Thorpe inequalities hold.}
  \(T_+/A_j\to281/300>0\) and \(T_-/A_j\to961/2700>0\).
\item
  \textbf{The classical optimized scalar--Weyl moment bounds have
  slack.} Formula (14) gives \(U_j>F(h_j)\). Choose a formal Einstein
  energy \(e_j=(U_j+F(h_j))/2\); it has positive energies for both Weyl
  halves and satisfies the scalar and all functional-pencil
  nonnegativity bounds for large \(j\). The equality case at \(U=F\) was
  already classically forbidden; the new content is the uniform open
  strip \(0<U-F\lesssim h^2\). These moment values are a numerical
  non-implication witness, not a metric construction.
\item
  \textbf{The degree-to-complex-hyperbolic entropy obstruction has
  slack.} Collapsing the added summands gives a degree-one map
  \(M_j\to X_j\). Kotschick's Corollary 3, equation (15), requires
  \(\chi(M_j)>\frac32|\tau(M_j)|+\frac{32}{81}\chi(X_j)\) for an
  Einstein metric. Here the left difference divided by \(A_j\) tends to
  \(F_0/6=961/5400>32/243\), so this test is satisfied strictly. Entropy
  supplies a lower-volume restriction, not the missing upper bound used
  in (5).
\item
  \textbf{The elementary Seiberg--Witten Miyaoka--Yau test has slack.}
  \((\chi-3\tau)(M_j)/A_j\to59/1800>0\).
\end{enumerate}

The mechanism is now different from a bare improvement of the
\(T_+/A=2/3\) threshold: the actual class forces a quantitative lower
energy envelope while the opposite Weyl half supplies a topological
upper envelope. Their incompatibility eliminates every candidate
Einstein metric on the same smooth type. It simultaneously rules out
arbitrarily small normalized integral error on those types.

\hypertarget{sources-and-reproducible-scope}{%
\subsection{Sources and reproducible
scope}\label{sources-and-reproducible-scope}}

\begin{itemize}
\tightlist
\item
  \href{https://arxiv.org/html/math/0003068}{LeBrun, \emph{Ricci
  Curvature, Minimal Volumes, and Seiberg--Witten Theory}, Proposition
  3.2, Theorem 3.3 and Corollary 2.5}. These provide the classical
  polarized gauge input, the one-third obstruction, and the exact
  equality comparison.
\item
  \href{https://arxiv.org/html/math/0410215}{Kotschick, \emph{Entropies,
  volumes, and Einstein metrics}, Theorem 1 and Corollary 3}. The
  comparison above uses its explicitly stated complex-hyperbolic degree
  bound, not an uncomputed exact minimal entropy for the new manifolds.
\item
  Mal'cev's residual-finiteness theorem for finitely generated linear
  groups in characteristic zero supplies the finite covers; see
  \href{https://arxiv.org/abs/1306.2385}{Nica's account and proof}. The
  unbounded-index deduction is given above.
\item
  The nonclassical differential input remains the local
  \href{line1_functional_pencil.md}{functional-pencil proof}, with its
  \href{line1_energy_envelope_gap.md}{energy-envelope transfer}. This
  note does not replace that proof by the finite calculations below.
\end{itemize}

\nolinkurl{check_line1_two_chirality_topology.py} verifies the
Chern--Weil and defect algebra, all rational ratios, the floor-error
estimates, and the displayed comparison margins; it rejects incorrect
opposite- orientation handle and trace-free-Ricci coefficients. It does
not certify the analytic pencil theorem, Seiberg--Witten gluing or
literature priority.

END COMPLETE SOURCE

\begin{center}\rule{0.5\linewidth}{0.5pt}\end{center}

\hypertarget{source-7-line1_bauer_furuta_geography.md}{%
\subsection{Source 7:
\nolinkurl{line1_bauer_furuta_geography.md}}\label{source-7-line1_bauer_furuta_geography.md}}

Original path:
\nolinkurl{mathematics/current/line1_bauer_furuta_geography.md}

SHA-256:
\nolinkurl{1ef42b7981722b1ec8a4325e85e460cc5dda1805b64827def3c7d4afe07de3e9}

BEGIN COMPLETE SOURCE

\hypertarget{two-chirality-integral-rigidity-beyond-ordinary-seibergwitten-type}{%
\section{Two-chirality integral rigidity beyond ordinary Seiberg--Witten
type}\label{two-chirality-integral-rigidity-beyond-ordinary-seibergwitten-type}}

23 September 2026. This note applies the existing quantitative
functional pencil to Bauer--Furuta detected connected sums. The new
application is an actual family with even positive and negative Betti
numbers and no almost-complex structure in either orientation. The
differential estimate is inherited; the actual gauge supply, geography,
topology and application are proved together below. No claim of
exhaustive literature priority is made.

\hypertarget{the-analytic-input-and-the-distinct-gauge-category}{%
\subsection{1. The analytic input and the distinct gauge
category}\label{the-analytic-input-and-the-distinct-gauge-category}}

Use tensor norms for Ricci and the Hilbert--Schmidt norm on the Weyl
endomorphism of two-forms. For a closed connected oriented smooth
four-manifold set

\[
 T_\pm=2\chi\pm3\tau,\qquad
 \mathcal E(g)=\inf_{\lambda\ne0}
 \frac{\int|\operatorname{Ric}_g-\lambda g|^2}
      {\lambda^2\operatorname{Vol}_g(M)}.
\]

The actual-class version of the existing
\href{line1_two_chirality_topology.md}{two-chirality theorem}, Section
2, says: if there is a number \(A>0\) such that every metric admits an
actual conformally solvable Spin\(^c\) structure satisfying

\[
 (a_g^+)^2\ge A,\qquad a_g^2\ge T_+,
 \qquad h=1-T_+/A\in(0,1/3),
\]

then a universal constant \(c>0\) gives

\[
 \mathcal E(g)+(U-F(h))_+\ge ch^2,\qquad
 U=\frac{3\min(T_+,T_-)}A,
 \qquad F(h)=\frac{(4-\sqrt{1-3h})^2}{9}.
 \tag{1}
\]

The constant is independent of the number of ordinary basic classes, the
rank of cohomology, and the virtual dimension of the chosen structure.
The only nonclassical analytic premise in this note is (1), whose proof
uses the \href{line1_functional_pencil.md}{functional-pencil theorem}
and \href{line1_energy_envelope_gap.md}{energy-envelope transfer}.

The following source theorem supplies actual structures where ordinary
Seiberg--Witten invariants vanish. For two or three simply connected
symplectic four-manifolds \(X_i\) with \(b_+(X_i)\equiv3\pmod4\),
Bauer's nonvanishing theorem gives nonzero Bauer--Furuta invariants for
the connected sum of their canonical structures. Adding a manifold with
negative-definite intersection form and determinant square \(-b_2\)
preserves this nonvanishing. Independent canonical and exceptional signs
give actual monopole classes. See Ishida--LeBrun, Theorem 9, Corollary
8, Proposition 10 and Corollary 11 in
\href{https://arxiv.org/html/math/0111228}{\emph{Curvature, Connected
Sums, and Seiberg--Witten Theory}}. For two or three positive summands
no fourth-summand congruence is needed.

\hypertarget{a-quantitative-theorem-for-decomposable-symplectic-building-blocks}{%
\subsection{2. A quantitative theorem for decomposable symplectic
building
blocks}\label{a-quantitative-theorem-for-decomposable-symplectic-building-blocks}}

\textbf{Theorem 1.} Let \(r\in\{2,3\}\). Let \(X_1,\ldots,X_r\) be
simply connected symplectic four-manifolds with
\(b_+(X_i)\equiv3\pmod4\), and suppose

\[
 A=\sum_{i=1}^r c_1^2(X_i)>0.
\]

For nonnegative integers \(k,m\), put

\[
 M=\Bigl(\#_{i=1}^r X_i\Bigr)
       \# k\overline{\mathbb{CP}}^{\,2}
       \# m(S^1\times S^3),\qquad
 D=4(r-1)+k+4m,\qquad h=D/A.
 \tag{2}
\]

If \(0<h<1/3\), every smooth metric on \(M\) satisfies (1), with the
same universal constant as in the existing two-chirality theorem. In
particular, if

\[
 U\le F(h)+\frac c2h^2,
 \tag{3}
\]

then

\[
 \mathcal E(g)\ge\frac c2h^2,
 \qquad
 \|g^{-1}\operatorname{Ric}_g-\sigma I\|_{\infty,\mathrm{op}}
       \ge h\sqrt{c/8}\quad(\sigma=\pm1).
 \tag{4}
\]

Under condition (3), \(M\) has no Einstein metric of any sign.

\textbf{Proof.} Canonical classes on the symplectic summands have
mod-two Seiberg--Witten invariant one. Write \(c_i=c_1(X_i)\) and
\(\alpha=\sum_i c_i\), using the orthogonal connected-sum splitting. The
source theorem gives, for every exceptional sign choice, an actual
monopole class

\[
 a_{\epsilon}=\alpha+\sum_{j=1}^k\epsilon_jE_j,
 \qquad \epsilon_j\in\{\pm1\},\quad E_iE_j=-\delta_{ij}.
 \tag{5}
\]

Here \(\epsilon=(\epsilon_1,\ldots,\epsilon_k)\). These are first Chern
classes of specified Spin\(^c\) structures with solutions for every
metric, not merely formal characteristic vectors. The handle factors are
included in the negative-definite gluing theorem; its proof allows a
nontrivial Picard torus on that summand.

Topology and the almost-complex structure on each summand give

\[
 T_+(M)=A-4(r-1)-k-4m=A-D,
 \qquad a_{\epsilon}^2=A-k,
\]

\[
 a_{\epsilon}^2-T_+(M)=4(r-1+m)\ge0.
 \tag{6}
\]

In particular these are special classes, with actual virtual dimension
\(r-1+m\). For any metric, averaging (5) over the exceptional signs
gives

\[
 2^{-k}\sum_{\epsilon}(a_{\epsilon,g}^+)^2
   =(\alpha_g^+)^2+\sum_j(E_{j,g}^+)^2
   \ge\alpha^2=A.
 \tag{7}
\]

Choose one maximizing structure. The choice can depend on the metric,
but the resulting structure solves the unperturbed equations on all
metrics and hence on every conformal metric. The same chosen class
therefore supplies every positive auxiliary weight in the analytic
theorem. Its square is positive here, since \(k/A\le h<1/3\), so the
negative-square cost is actually zero. Equations (6)--(7) verify all
inputs to (1), proving (1) and (4).

A nonzero Einstein constant makes \(\mathcal E=0\), contradicting (4).
Ricci-flatness is not a zero-defect case of this definition; instead the
actual class and (7) give the signed scalar estimate
\(\int s_-^2\ge32\pi^2A>0\), which rules it out separately. This proves
the final assertion. \(\square\)

This extension changes the source of the gauge information. Splitting
\(M\) into two pieces of positive \(b_+\) makes its ordinary
Seiberg--Witten invariants vanish by the standard connected-sum
vanishing theorem. Actual Bauer--Furuta classes persist nevertheless.
Thus ordinary basic classes on \(M\) cannot be silently substituted for
(5).

\hypertarget{why-low-signature-building-blocks-do-not-realize-the-new-strip}{%
\subsection{3. Why low-signature building blocks do not realize the new
strip}\label{why-low-signature-building-blocks-do-not-realize-the-new-strip}}

This is a useful selection calculation, not a general obstruction to the
connected-sum method. Write

\[
 \rho=\frac{\sum_i\tau(X_i)}A,\quad x=k/A,
 \qquad h=\frac{4(r-1)+k+4m}A.
\]

The exact opposite-orientation budget is

\[
 \frac{3T_-}A=3(1-6\rho)-3h+18x.
 \tag{8}
\]

For zero-signature building blocks and \(h<1/3\), the signature of \(M\)
is \(-k\); hence \(U=3T_+/A=3(1-h)>2\), while \(F(h)<16/9\). They cannot
approach \(U=F(h)\) in this range. This explains why simply taking
several product-like inputs does not realize the desired two-chirality
strip.

For positive signature, (8) instead shows what must change: the building
blocks need a sufficiently positive ratio \(\rho\), and handle addition
must supply enough of the surgery budget. Both are realized in the
following construction.

\hypertarget{an-actual-infinite-family-with-no-almost-complex-structure}{%
\subsection{4. An actual infinite family with no almost-complex
structure}\label{an-actual-infinite-family-with-no-almost-complex-structure}}

Braungardt--Kotschick's geography theorem realizes every sufficiently
large lattice point strictly below any fixed slope less than nine by
simply connected minimal symplectic manifolds: Theorem 4 of
\href{https://arxiv.org/pdf/math/0306021}{\emph{Einstein metrics and the
number of smooth structures on a four-manifold}}. Take its parameter
\(\epsilon=1/20\). Consequently, for every sufficiently large integer
\(n\), choose such a manifold \(X_n\) with

\[
 \chi_h(X_n)=20n,\qquad c_1^2(X_n)=178n.
 \tag{9}
\]

Indeed \(178n\le(9-1/20)20n-c(1/20)=179n-c(1/20)\) for large \(n\). Here
\(\chi_h=(\chi+\tau)/4\) is the topological Todd number for a symplectic
manifold; no complex structure on \(X_n\) is assumed. The two defining
topological equations imply

\[
 \chi(X_n)=62n,\quad \tau(X_n)=18n,\quad
 b_+(X_n)=40n-1\equiv3\pmod4,\quad b_-(X_n)=22n-1.
 \tag{10}
\]

Take two copies of \(X_n\). Set

\[
 A_n=356n,\quad h_0=\frac{19}{300},\quad F_0=F(h_0)=\frac{961}{900},
 \quad u_0=\frac{105}{89},
\]

\[
 x_0=\frac{781}{180225},\qquad y_0=\frac{42533}{2883600}.
 \tag{11}
\]

These positive rational numbers satisfy

\[
 x_0+4y_0=h_0,\qquad u_0+15x_0-12y_0=F_0.
 \tag{12}
\]

Define the even nonnegative surgery counts, for sufficiently large
\(n\),

\[
 k_n=2\left\lfloor\frac{x_0A_n}{2}\right\rfloor,
 \qquad m_n=2\left\lfloor\frac{y_0A_n}{2}\right\rfloor-4,
\]

\[
 M_n=X_n\#X_n\#k_n\overline{\mathbb{CP}}^{\,2}
                    \#m_n(S^1\times S^3).
 \tag{13}
\]

\textbf{Theorem 2 (realized Bauer--Furuta strip).} There exists \(n_0\)
such that, for every \(n\ge n_0\) and every smooth metric on \(M_n\),

\[
 \boxed{\mathcal E(g)\ge\frac{c h_0^2}{8}>0,\qquad
 \delta_\sigma(g)\ge h_0\sqrt{c/32}\quad(\sigma=\pm1).}
 \tag{14}
\]

The manifolds \(M_n\) are pairwise nonhomeomorphic, have vanishing
ordinary Seiberg--Witten invariants, and admit no almost-complex
structure in either orientation. Both strict Hitchin--Thorpe
inequalities hold. Their numerical data lie on the allowed side of the
classical optimized energy envelope, but inside its new quantitative
exclusion strip for large \(n\).

\textbf{Proof.} Connected-sum topology gives

\[
 \chi(M_n)=124n-2+k_n-2m_n,\qquad \tau(M_n)=36n-k_n,
\]

\[
 T_+=A_n-4-k_n-4m_n,
 \quad T_-=140n-4+5k_n-4m_n.
 \tag{15}
\]

Write \(k_n=x_0A_n-\alpha_n\) and \(m_n=y_0A_n-\beta_n-4\), with
\(0\le\alpha_n,\beta_n<2\). Since
\(\tau(M_n)/A_n\ge9/89-x_0=196/2025>0\), the smaller chirality is
\(T_-\). Thus, exactly,

\[
 h_n=h_0-\frac{\alpha_n+4\beta_n+12}{A_n},\qquad
 U_n=F_0+\frac{-15\alpha_n+12\beta_n+36}{A_n}.
 \tag{16}
\]

For large \(n\), \(0<h_n<h_0\). The numerator of the correction in
\(U_n\) lies strictly between 6 and 60. Also

\[
 F'(h)=\frac{4-\sqrt{1-3h}}{3\sqrt{1-3h}},\qquad
 0<F'(h)\le31/27<2\quad(0\le h\le h_0).
\]

Since \(\alpha_n+4\beta_n+12<22\), it follows that

\[
 \boxed{0<U_n-F(h_n)<\frac{104}{A_n},\qquad h_n\longrightarrow h_0.}
 \tag{17}
\]

Theorem 1 applies by (10). Choose \(n_0\) so that \(h_n\ge h_0/2\) and
\(104/A_n\le c h_0^2/8\). Then (1) implies
\(\mathcal E\ge ch_0^2/4-ch_0^2/8\), proving (14).

The Betti numbers are

\[
 b_1(M_n)=m_n,\qquad b_+(M_n)=80n-2,\qquad
 b_-(M_n)=44n-2+k_n.
 \tag{18}
\]

They are all even. An almost-complex four-manifold necessarily has
\((\chi+\tau)/4\in\mathbb Z\), equivalently \(1-b_1+b_+\) even. Equation
(18) makes this number odd. Reversing orientation replaces \(b_+\) by
the also-even \(b_-\), so the same contradiction holds there. No
almost-complex structure exists in either orientation. The positive and
negative summands in (18) also allow the ordinary connected-sum
vanishing argument in both orientations. The second Betti number
\(b_2=124n-4+k_n\) is strictly increasing, which gives pairwise distinct
homeomorphism types even without fixing the orientation.

Finally \(T_+/A_n\to281/300>0\) and \(T_-/A_n\to961/2700>0\), proving
strict Hitchin--Thorpe. By (17), the classical Einstein energy interval
\([F(h_n),U_n]\) has positive length. Both Weyl energies can be positive
inside it. It is the quantitative separation, not the classical
scalar--Weyl numerical inequalities alone, that eliminates this
interval. \(\square\)

\hypertarget{the-precise-increase-in-scope}{%
\subsection{5. The precise increase in
scope}\label{the-precise-increase-in-scope}}

\begin{enumerate}
\def\labelenumi{\arabic{enumi}.}
\tightlist
\item
  These are not merely new choices of a complex general-type base in the
  earlier application. A Kähler surface has odd \(b_+=1+2p_g\), and
  blowups and \(S^1\times S^3\) additions preserve \(b_+\). Both \(b_+\)
  and \(b_-\) in (18) are even, so neither orientation of \(M_n\) can be
  diffeomorphic to any member of that former category.
\item
  Ordinary Seiberg--Witten invariants on the whole manifold vanish. The
  application genuinely uses actual classes retained by stable
  cohomotopy. It is not obtained by declaring arbitrary cohomology
  vectors to be monopoles. The virtual dimension is \(1+m_n\), not zero.
\item
  The previously recorded universal one-third estimate has
  \(\varepsilon_0\le1/4\), whereas \(1/3-h_n>1/3-h_0=27/100\). It does
  not itself force positive Einstein defect on this family. The
  classical connected-sum nonexistence threshold \(D\ge A/3\) likewise
  does not apply.
\item
  The new differential input is still the existing functional pencil.
  Bauer's nonvanishing and Braungardt--Kotschick geography are classical
  inputs. This note supplies a realized extension of the quantitative
  obstruction, not a new Bauer--Furuta theorem or new symplectic
  geography.
\item
  These manifolds cannot be Yau26 counterexamples: (14) forbids the
  required arbitrarily tight pointwise inputs. The result establishes a
  further category of no-Einstein smooth types on which the original
  pointwise hypotheses are uniformly impossible.
\end{enumerate}

For large \(n\), \(\pi_1(M_n)\) is the infinite free group of rank
\(m_n\). Myers' theorem already excludes strictly positive Ricci
curvature on these objects. The positive-sign pointwise corollary is
therefore not a second independent discovery; the substantive extension
is the uniform all-metrics integral gap and its negative-sign pointwise
consequence, supplied by the actual Bauer--Furuta classes.

The threshold \(n_0\) depends on the inherited analytic constant and the
geography existence threshold. No explicit numerical value is claimed.
We choose one available \(X_n\) at each lattice point. Although the
source geography gives infinitely many smooth structures there, this
note does not assert that all remain pairwise distinct after the new
connected sum.

\nolinkurl{check_line1_bauer_furuta_geography.py} checks the exact
topology, rational parameters, parity obstruction, sign of the floor
errors, and comparison with the old recorded band. Those checks do not
prove gauge nonvanishing, the geography theorem or the inherited
analytic theorem.

END COMPLETE SOURCE

\begin{center}\rule{0.5\linewidth}{0.5pt}\end{center}

\hypertarget{source-8-round14_twisted_rigidity.md}{%
\subsection{Source 8:
\nolinkurl{round14_twisted_rigidity.md}}\label{source-8-round14_twisted_rigidity.md}}

Original path:
\nolinkurl{mathematics/current/round14_twisted_rigidity.md}

SHA-256:
\nolinkurl{77b52d2c7a8b3e8a9fd6a71f0109cbebcd825e34d428b1be30dffaf1bb505bc8}

BEGIN COMPLETE SOURCE

\hypertarget{twisted-monopole-extension-and-positive-b-plus-connected-sums}{%
\section{Twisted monopole extension and positive-b-plus connected
sums}\label{twisted-monopole-extension-and-positive-b-plus-connected-sums}}

\hypertarget{a-larger-analytic-input-than-all-metric-monopoles-on-a-cover}{%
\subsection{1. A larger analytic input than all-metric monopoles on a
cover}\label{a-larger-analytic-input-than-all-metric-monopoles-on-a-cover}}

The finite-cover criterion in \nolinkurl{round14_rigidity.md} does not
exhaust the available gauge input. For its analytic proof, the cover
equations only need to be solvable for the \textbf{deck-invariant
curvature weights actually used by the proof}. Pin-minus monopoles
supply exactly this input, including on connected sums whose ordinary
Seiberg--Witten invariants vanish. The argument below checks that
restriction explicitly.

\textbf{Invariant-weight lemma.} Let \((\widetilde X,\widetilde g)\) be
a closed oriented four-manifold with a finite group of
orientation-preserving isometries \(G\). Fix one ordinary spin-c
structure with real Chern class \(\widetilde a\). Suppose that for every
smooth positive \(G\)-invariant function \(f\), the weighted equations

\[
 D_A\Phi=0,\qquad -iF_A^+=f\sigma(\Phi)                 \tag{1.1}
\]

have a solution in this structure. Then the integral calibration theorem
of round 4 holds for \((\widetilde X,\widetilde g,\widetilde a)\), with
its original universal constant and square cost. Solutions are not
required to be invariant, and no solvability on non-invariant metrics or
for non-invariant weights is asserted.

\textbf{Proof.} There are only three kinds of weight in the full proof
chain.

\begin{enumerate}
\def\labelenumi{\arabic{enumi}.}
\tightlist
\item
  The signed scalar inequality uses \(f=1\).
\item
  The spectral inequality uses \(f_\epsilon=\sqrt{q_\epsilon}\), where
  \(q_\epsilon\) is a positive smooth majorant of
  \(h=(-2s/3-2\lambda_{\min}(W^+))_+\), converging uniformly to \(h\).
\item
  The bounded-spinor estimate uses positive smooth majorants of
  \(f_0=\max\{\sqrt{\max(-s+\sqrt6|W^+|,3)},(-s)_+/4\}\).
\end{enumerate}

Both \(h\) and \(f_0\) are \(G\)-invariant. Start with any smooth
majorant converging uniformly and replace it by
\(|G|^{-1}\sum_{\gamma\in G}f\circ\gamma\). This preserves positivity,
the majorant property and uniform convergence. The square root in case 2
also remains invariant. Thus (1.1) supplies every solution invoked in
round 3, Section 6, and round 4, Sections 3--6.

All subsequent operations take place on the smooth ordinary spinor and
line bundles upstairs: Weitzenbock identities, the Hessian estimate, the
determinant-curvature identity, the polynomial Weyl tensor and its weak
Bianchi self-test. They impose no equivariance condition on the spinor
or test tensor. The measurable lowest-eigenline projector need not be
selected equivariantly either; it is used only for pointwise norm
estimates. No additional weighted solution is requested at that step.
Therefore the proof and all constants are unchanged. QED.

\hypertarget{twisted-integral-gap-theorem}{%
\subsection{2. Twisted integral-gap
theorem}\label{twisted-integral-gap-theorem}}

Let \(\pi:\widetilde X\to X\) be a nontrivial double cover of a closed
oriented four-manifold, and let \(\ell\) denote the associated real sign
local system. A genuine spin-c-minus structure has an associated
\(O(2)\) characteristic bundle and a twisted Euler class
\(a\in H^2(X;\ell)\). Assume that for the metric being tested its
Pin-minus equations have solutions throughout that conformal class. This
assumption holds, in particular, for a nonzero Pin-minus basic class
with \(b_+^\ell\ge2\).

The required lifting statement is the one in
\href{https://arxiv.org/pdf/1505.05948}{Ishida--Matsuo--Nakamura,
Section 2 and Proposition 8}: the pulled-back structure is ordinary
spin-c, with Chern class \(\widetilde a=\pi^*a\), and a lifted twisted
solution is an ordinary Seiberg--Witten solution, after matching the
constant spinor normalization. Consequently conformal covariance
supplies (1.1) for every positive deck-invariant weight: the weight
descends, the corresponding conformal metric is still a metric on \(X\),
and one lifts its solution before applying the ordinary conformal
transformation.

For the lifted metric, the harmonic projection and intersection form
obey

\[
 \widetilde a^2=2a^2,\qquad
 (\widetilde a^+)^2=2(a^+)^2,\qquad
 \widetilde V=2V,\qquad \widetilde T=2T.                 \tag{2.1}
\]

Here products of two sign-local-system classes are untwisted. Pullback
commutes with the twisted Laplacian and with Hodge star; the factor two
in the norms is integration over the covering sheets. This verifies the
normalization rather than treating the twisted Euler class as an
ordinary Chern class downstairs.

\textbf{Theorem.} Suppose that for every metric \(g\) on \(X\) one may
select such a twisted class on a fixed double cover, with the same class
working for all weights in that conformal class, and

\[
                  (a_g^+)^2\ge A>0,\qquad a_g^2\ge T_X.
                                                               \tag{2.2}
\]

Then the existing universal constant gives

\[
 \boxed{\displaystyle
 \mathcal E_X(g)+\left(\frac{T_X}{A}-\frac23\right)_+
                  \ge c_*>0\qquad\hbox{for every }g.}           \tag{2.3}
\]

If the last square hypothesis is instead \(a_g^2\ge-C_A A\), the
corresponding constant is \(c(C_A)>0\). There is no degree loss or new
geometric condition.

\textbf{Proof.} Apply the invariant-weight lemma upstairs. The Hessian
proof uses the actual Hermitian positive-spinor bundle and its line
connection there, so there is no missing global complex structure on the
downstairs real spinor bundle. Its determinant energy remains

\[
 \int_{\widetilde X}|F_A|^2
    =2\int_{\widetilde X}|F_A^+|^2-4\pi^2\widetilde a^2.
\]

The normalized negative-square cost is
\(( -\widetilde a^2)_+/\widetilde V=(-a^2)_+/V\). The lifted monopole
quadratic form is sign-equivariant; its square \(\psi\otimes\psi\), and
hence \(A_\psi=|\psi|^2I/24-\psi\otimes\psi/8\), is an ordinary
Weyl-operator model. Alternatively all testing may simply be performed
upstairs without descending this model. The topology-transfer algebra,
positive-sign scalar obstruction and optimization in the Einstein
constant are unchanged. Use (2.1) and exact invariance of \(\mathcal E\)
under pullback to obtain (2.3). QED.

The square hypothesis in (2.2) is explicit. It is \textbf{not} inferred
from the virtual dimension formula of arbitrary Pin-minus moduli spaces,
which differs from the ordinary SW dimension formula. The concrete
application below verifies it from the actual glued classes.

\hypertarget{concrete-positive-b-plus-surgeries}{%
\subsection{3. Concrete positive-b-plus
surgeries}\label{concrete-positive-b-plus-surgeries}}

Let \(M\) be a closed oriented four-manifold with \(b_+(M)\ge2\) and a
spin-c structure with odd ordinary SW invariant and

\[
                     c^2=A=T_M=2\chi(M)+3\tau(M)>0.
\]

For example, any symplectic \(M\) with positive canonical square and
\(b_+\ge2\) supplies such a class. Let \(N\) have \(b_+(N)=0\), set
\(k=b_2(N)\), \(l=b_1(N)\), and let

\[
 Z=\mathop{\#}_{j=1}^{n}(S^2\times\Sigma_{r_j})
      \#\mathop{\#}_{i=1}^{m}(S^1\times Y_i),
 \qquad r_j\ge1,\qquad n+m\ge1,                        \tag{3.1}
\]

where the \(Y_i\) are arbitrary closed oriented three-manifolds. Put
\(X=M\#N\#Z\). The choice \(N=S^4\) simply means no \(N\) summand. No
Yamabe assumption on \(N\) or \(Y_i\) is imposed: only the nonvanishing
theorem, not the paper's Yamabe-equality conclusion, is used.

The classical input is
\href{https://arxiv.org/pdf/1505.05948}{Ishida--Matsuo--Nakamura,
Proposition 9 and its proof}. It provides actual Pin-minus basic
classes, with independently selectable exceptional signs, of the form

\[
 a=c+\sum_{j=1}^{k}\epsilon_j e_j+t,
 \qquad (a_g^+)^2\ge A,
 \qquad a^2=A-k.                                       \tag{3.2}
\]

Here the first summand is untwisted, the \(e_j\) form the negative
diagonal part, and \(t\) is the torsion twisted class on \(Z\). The
nonvanishing proof, including its local-system construction, is an
external gauge-theoretic input; the Euler-characteristic computation
alone does not produce these classes.

For connected sums, \(T(U\#V)=T(U)+T(V)-4\). Since \(T_N=4-4l-k\),
\(T_{S^2\times\Sigma_r}=8-8r\), and \(T_{S^1\times Y}=0\), define

\[
       D=k+4l+\sum_{j=1}^{n}(8r_j-4)+4m.
\]

Then

\[
 T_X=A-D,\qquad a^2-T_X=D-k
          =4l+\sum_j(8r_j-4)+4m\ge0.                   \tag{3.3}
\]

The square condition in (2.2) is now verified. Thus

\[
 \boxed{\displaystyle
 \mathcal E_X(g)+\left(\frac13-\frac DA\right)_+\ge c_* .}
                                                               \tag{3.4}
\]

This permits positive-b-plus connected-sum summands, unlike the earlier
negative-definite connected-sum corollary. When at least one
\(S^2\times\Sigma_{r_j}\) occurs, ordinary SW and ordinary
stable-cohomotopy SW invariants vanish in this family, as noted in the
same paper's introduction. Their vanishing cannot be used to dismiss the
twisted input. Its double cover also has several positive-b-plus
connected-sum factors, so the proof must not pretend that it has an
ordinary nonzero SW invariant.

\hypertarget{an-explicit-critical-family-with-vanishing-ordinary-invariants}{%
\subsubsection{An explicit critical family with vanishing ordinary
invariants}\label{an-explicit-critical-family-with-vanishing-ordinary-invariants}}

For each \(j\ge1\), take

\[
 M_j=\Sigma_2\times\Sigma_{3j+1},\qquad
 Z_j=\mathop{\#}_{i=1}^{2j}(S^2\times T^2),\qquad
 X_j=M_j\#Z_j.
\]

The product canonical invariant is \(\pm1\), and

\[
 A_j=8(2-1)((3j+1)-1)=24j,\qquad
 D_j=4(2j)=8j,\qquad
 T_{X_j}=16j,\qquad \frac{T_{X_j}}{A_j}=\frac23.
\]

The new estimate therefore gives the uniform endpoint statement

\[
                  \boxed{\mathcal E_{X_j}(g)\ge c_*}
\]

for every \(j\) and every smooth metric. Ordinary SW and ordinary
stable-cohomotopy SW invariants vanish on these manifolds, while the
twisted classes in (3.2) remain available. This is the quantitative
critical situation: the classical strict exact inequality alone does not
supply a common positive gap for arbitrary metrics across this infinite
family.

\hypertarget{application-to-the-new-closed-shear-manifolds}{%
\subsubsection{Application to the new closed shear
manifolds}\label{application-to-the-new-closed-shear-manifolds}}

Take the smallest closed shear example from round 12: \(\chi(M)=4\),
\(\tau(M)=0\), \(b_+(M)=5\), hence \(A=8\). The canonical invariant is
\(\pm1\). For

\[
                         X=M\#(S^2\times T^2)
\]

one has \(D=4\), \(T_X=4\), and (3.4) becomes

\[
                          \mathcal E_X(g)\ge c_* .      \tag{3.5}
\]

The unmodified \(M\) admits the already constructed \(\delta_-\to0\)
metrics, but \textbf{no construction on this modified smooth manifold
can retain arbitrarily small integral Einstein defect}. This excludes
all pointwise near-Einstein constructions there as well. It is an
ansatz-independent conclusion, despite the ordinary invariants of the
modified manifold being zero. Ricci-flat metrics are separately excluded
by the positive twisted monopole scalar estimate.

This last example has \(D/A=1/2\), strictly beyond the classical
critical ratio \(1/3\). A positive defect bound there is also available
from classical twisted curvature estimates away from equality. It is a
useful application to the current construction, not the appropriate
example on which to base the novelty of the quantitative endpoint.

\hypertarget{position-relative-to-the-known-theorem}{%
\subsection{4. Position relative to the known
theorem}\label{position-relative-to-the-known-theorem}}

The cited 2016 paper already proves exact Einstein obstructions on this
connected-sum family; its Theorem 16 requires \(D<A/3\) for Einstein
existence. No priority is claimed for that classical obstruction or for
Pin-minus equations. The present consequence replaces a strict exact
threshold by the uniform, all-metric relative integral inequality (3.4),
using the existing self-test theorem with genuinely verified twisted
input.

The extension does not improve the unmodified closed shear ratio. Its
near-Einstein sequence bounds twisted projection budgets just as it
bounds ordinary ones; applying a twisted scalar inequality along that
sequence cannot create a larger budget for the same manifold. The
positive result is instead a new range of surgeries and quotients to
which the integral rigidity mechanism applies.

The finite checks in \nolinkurl{round14_rigidity_check.py} include the
factors of two, the connected-sum loss and the explicit shear
application. Weighted solvability and the invariant-weight lemma are
proved above, not tested by those finite calculations.

END COMPLETE SOURCE

\begin{center}\rule{0.5\linewidth}{0.5pt}\end{center}

\hypertarget{source-9-lebrun_stability_closure.md}{%
\subsection{Source 9:
\nolinkurl{lebrun_stability_closure.md}}\label{source-9-lebrun_stability_closure.md}}

Original path:
\nolinkurl{mathematics/current/lebrun_stability_closure.md}

SHA-256:
\nolinkurl{19c436bcd1edf762899c498478ee9be11956ac38c80d9a674f365375c014a9d3}

BEGIN COMPLETE SOURCE

\hypertarget{statement-and-analytic-dependency}{%
\section{Statement and analytic
dependency}\label{statement-and-analytic-dependency}}

Let \(X\) be a closed connected oriented smooth four-manifold, let

\[
T_X=2\chi(X)+3\tau(X)>0,
\qquad M=X\#k\overline{\mathbb{CP}}^{\,2}
                 \#\ell(S^1\times S^3),\qquad D=k+4\ell.
\]

Assume that a Spin\(^c\) structure with real first Chern class \(c\) has
a nonzero ordinary Seiberg--Witten invariant, with its insertion
retained if the expected dimension is positive. When \(b_+(X)=1\),
nonvanishing is in the small-perturbation chamber, not an arbitrarily
chosen chamber. Write \(A=c^2\ge T_X>0\), and define

\[
\mathcal E(g)=\inf_{\lambda\ne0}
 \frac{\int_M|\operatorname{Ric}_g-\lambda g|_g^2\,d\mu_g}
 {\lambda^2\operatorname{Vol}_g(M)}.
\]

\textbf{Stability theorem, relative to the integral calibration lemma
below.} There is a universal constant \(\varepsilon>0\) such that every
smooth metric on \(M\) satisfies both

\[
\mathcal E(g)+\left(\frac{T_X-D}{A}-\frac23\right)_+
 \ge\varepsilon,
\qquad
\mathcal E(g)+\left(\frac13-\frac D{T_X}\right)_+
 \ge\varepsilon.                                      \tag{1}
\]

In particular, if \(X\) is any minimal complex surface of general type,
including \(b_+(X)=1\), then \(A=T_X=c_1^2(X)>0\), and

\[
\boxed{\mathcal E(g)+
 \left(\frac13-\frac{k+4\ell}{c_1^2(X)}\right)_+
 \ge\varepsilon.}                                    \tag{2}
\]

This note proves the gauge-to-obstruction implication, including its
constants and the zero Einstein constant. Its nonclassical analytic
input is stated precisely in Section 3. The full proof of that lemma
remains the integral calibration manuscript; restating it here is not an
independent verification of its Hessian and weak-Bianchi arguments.
Neither a universal supply of special monopole classes on arbitrary
four-manifolds nor positivity of a universal special-class invariant is
assumed.

\hypertarget{the-actual-polarized-class-and-every-auxiliary-weight}{%
\section{The actual polarized class and every auxiliary
weight}\label{the-actual-polarized-class-and-every-auxiliary-weight}}

\hypertarget{sign-selection-and-the-chamber}{%
\subsection{Sign selection and the
chamber}\label{sign-selection-and-the-chamber}}

Fix \(g\), and let \(H=H_g^+\) be its positive harmonic subspace. Pull
\(c\) back to \(M\). Choose exceptional generators \(E_{j,H}\),
\textbf{with signs chosen at this \(H\)}, so that

\[
E_{i,H}\cdot E_{j,H}=-\delta_{ij},\qquad
c\cdot E_{j,H}=0,\qquad c_H^+\cdot E_{j,H}\le0.
\]

Polarized Seiberg--Witten gluing supplies an actual Spin\(^c\) structure
with

\[
a_H=c-\sum_j E_{j,H}                                  \tag{3}
\]

and nonzero invariant at \(H\). Thus its unperturbed equations are
solvable for every \(H\)-adapted metric. The classical input is the
construction in LeBrun's Proposition 3.2 {[}L01{]}, not a consequence
merely of the following intersection-form calculation:

\[
\begin{aligned}
(a_H^+)^2
 &=(c_H^+)^2-2\sum_j c_H^+\cdot E_{j,H}
                  +\left|\sum_j E_{j,H}^+\right|^2
 \ge(c_H^+)^2\ge A,\\
a_H^2&=A-k\ge T_X-k-4\ell=T_M.
\end{aligned}                                           \tag{4}
\]

For \(b_+=1\), choose the positive generator \(\omega_H\) so that
\(c\cdot\omega_H>0\). This is possible because \(c^2>0\). The sign
selection implies

\[
a_H\cdot\omega_H=c\cdot\omega_H-\sum_j E_{j,H}\cdot\omega_H>0.
\]

The corresponding base projection of \(\omega_H\) has positive square.
Hence the small-perturbation chamber agrees with the one in the
chamber-compatible blow-up formula. Handle additions do not change the
intersection form. If the base invariant uses an insertion \(z\), retain
it and add the \(\ell\) degree-one loop insertions. Indeed,

\[
\frac{a_H^2-T_M}{4}=\frac{A-T_X}{4}+\ell.
\]

The surgery and insertion statements used here are classical gauge
theory {[}OS{]}; this dimension equality alone would not prove
nonvanishing.

For a minimal general-type complex surface, the canonical Spin\(^c\)
structure has \(c=-K_X\), \(K_X^2>0\), and \(K_X\cdot[\omega]>0\) for a
Kähler form. When \(b_+=1\), put \(K_X\) in the forward cone. Its
pairing with every unit forward positive generator is at least
\(\sqrt{K_X^2}>0\). The zero-perturbation discriminant for \(-K_X\) is
therefore negative, identifying the small chamber with the canonical
symplectic chamber where nonvanishing holds {[}LL{]}. This is why
\(b_+=1\) is included. A symplectic structure without this chamber
verification would not suffice.

\hypertarget{conformal-covariance-with-the-bundle-identification-explicit}{%
\subsection{Conformal covariance, with the bundle identification
explicit}\label{conformal-covariance-with-the-bundle-identification-explicit}}

Let \(\mathfrak q_g\) denote the quadratic spinor-to-self-dual-two-form
map, normalized by \(|\mathfrak q_g(\Phi)|^2=|\Phi|^4/8\). For any
smooth \(f>0\), set

\[
\widehat g=f^{-2}g=e^{2u}g,\qquad u=-\log f.
\]

The Hodge star on two-forms is unchanged, so
\(H_{\widehat g}^+=H_g^+=H\). Use the \textbf{same} Spin\(^c\) structure
from (3) to solve the unperturbed equations for \(\widehat g\):

\[
D_{\mathcal A}^{\widehat g}\widehat\Phi=0,\qquad
-iF_{\mathcal A}^+=\mathfrak q_{\widehat g}(\widehat\Phi).
\]

Let \(\beta:S_g\to S_{\widehat g}\) be the fibrewise unitary spinor
identification. The conformal identities are

\[
D_{\mathcal A}^{\widehat g}(e^{-3u/2}\beta\Phi)
 =e^{-5u/2}\beta D_{\mathcal A}^{g}\Phi,
\qquad
\mathfrak q_{\widehat g}(\beta v)=e^{2u}\mathfrak q_g(v),
\]

where the second equality compares covariant two-forms. With
\(\Phi=f^{-3/2}\beta^{-1}\widehat\Phi\), they give

\[
D_{\mathcal A}^{g}\Phi=0,\qquad
\mathfrak q_g(\Phi)=f^{-1}\mathfrak q_{\widehat g}(\widehat\Phi),
\qquad -iF_{\mathcal A}^+=f\mathfrak q_g(\Phi).           \tag{5}
\]

This is the conformal computation underlying {[}L06, Lemma 3.7{]},
applied only to the one conformal class needed here. The exact order is

\[
\forall g\ \exists\mathfrak s_g\ \forall f\in C^\infty(M,(0,\infty))
\ \exists(\mathcal A_f,\Phi_f)\text{ solving (5)}.       \tag{6}
\]

No continuity of \(g\mapsto\mathfrak s_g\) is required. The class may
change between metrics; it is held fixed while \(f\) varies. Ordinary
solvability just at \(g\) would not establish (6).

\hypertarget{the-analytic-input-and-the-uniform-constant}{%
\section{The analytic input and the uniform
constant}\label{the-analytic-input-and-the-uniform-constant}}

All Ricci norms below are tensor Hilbert--Schmidt norms. The norm on
\(W^+\) is its Hilbert--Schmidt norm as an endomorphism of
\(\Lambda^+\). Put \(V=\operatorname{Vol}(g)\), \(T=T_M\), and

\[
\begin{gathered}
\mathscr O=\operatorname{Orb}\operatorname{diag}(-1/6,1/12,1/12),
F=\operatorname{dist}(W^+,\mathscr O)^2,\qquad
q=-s+\sqrt6|W^+|,\qquad p=(a_H^+)^2,\\
\Delta=\int q_+^2-72\pi^2p,\qquad
\kappa=\frac1V\left(\int(|\operatorname{Ric}+g|^2+F)+\Delta\right),
\qquad B=\frac{(-a_H^2)_+}{V}.
\end{gathered}
\]

\textbf{Integral calibration lemma (the nonclassical input).} For an
actual structure satisfying (6) and \(p>0\), the signed spectral
monopole inequality gives \(\Delta\ge0\), and a universal
\(c_{\rm cal}>0\) satisfies

\[
                    (1+B)\kappa\ge c_{\rm cal}.         \tag{7}
\]

The proof in the integral calibration manuscript constructs bounded
weighted spinors, obtains a covariant Hessian estimate using integral
Ricci and connection curvature, forms a smooth polynomial Weyl tensor,
and tests weak Bianchi against its divergence. Its hypotheses involve
the selected structure at the metric under consideration, not an
all-metric invariant. This note uses (7) explicitly; it does not replace
that proof by an appeal to a previous round's conclusion.

We next give the complete constant transfer. Put

\[
r=V^{-1}\int|\operatorname{Ric}+g|^2,\qquad
d=(T/A-2/3)_+.
\]

The algebraic comparison detailed in the appendix yields

\[
\kappa\le\frac{245}{2}r+972\pi^2\frac AVd.              \tag{8}
\]

The scalar monopole bound gives

\[
32\pi^2A\le\int s_-^2\le\int s^2\le32V+8rV,
\qquad A/V\le(4+r)/(4\pi^2).
\]

Consequently \(r\le1\) implies \(\kappa\le1215(r+d)\). For \(r\le1/4\),
Chern--Weil and the reverse triangle inequality give

\[
\begin{aligned}
4\pi^2T/V
&=V^{-1}\int\left(s^2/24+2|W^+|^2
                    -|\operatorname{Ric}^{\circ}|^2/2\right)\\
&\ge(4-2\sqrt r)^2/24-r/2\ge1/4.
\end{aligned}                                           \tag{9}
\]

Since \(a_H^2\ge T>0\), we have \(B=0\) in this branch, even when the
family allows arbitrarily many blow-ups. Equations (7)--(9) imply

\[
r+d\ge\varepsilon,\qquad
\varepsilon=\min\{1/4,c_{\rm cal}/1215\}>0.             \tag{10}
\]

For \(r>1/4\), (10) is immediate. This proves uniformity with respect to
the manifold, the metric, the number of exceptional classes, and their
selection. Smooth approximations of the auxiliary weight are taken with
\(g,a_H\) fixed before their approximation error tends to zero; no
uniform rate of smoothing across metrics is used.

For the positive normalization, set
\(r_+=V^{-1}\int|\operatorname{Ric}-g|^2\). On the scalar-negative set,
\(s_-^2\le(s-4)^2\), whence

\[
A/V\le r_+/(8\pi^2).
\]

If \(r_+\le1/4\) and \(d\le1/4\), the analogue of (9) gives
\(4\pi^2T/V\ge1/4\), whereas \(T\le A(2/3+d)\) gives

\[
4\pi^2T/V\le(r_+/2)(2/3+d)\le11/96<1/4.
\]

Thus \(r_++d>1/4\). For any \(\lambda\ne0\), apply the appropriate
normalization to \(\widetilde g=|\lambda|g\). Constant rescaling
preserves \(H\), (4), and the weighted input. In dimension four,

\[
\frac1{V_{\widetilde g}}\int
 |\operatorname{Ric}_{\widetilde g}-\operatorname{sign}(\lambda)
             \widetilde g|_{\widetilde g}^2
=\frac{\int|\operatorname{Ric}_g-\lambda g|_g^2}{\lambda^2V_g}.
\]

Taking the infimum proves the first inequality in (1). Apply the same
argument with the smaller positive projection budget \(T_X\le A\) to
obtain the second. No comparison of ratios with a potentially negative
numerator is necessary. Equation (2) follows by \(A=T_X\).

\hypertarget{scalar-flat-exclusion-and-all-einstein-constants}{%
\section{Scalar-flat exclusion and all Einstein
constants}\label{scalar-flat-exclusion-and-all-einstein-constants}}

With \(f=1\), the spinor Weitzenböck formula, integrated on the closed
manifold, gives

\[
0=\int(4|\nabla_{\mathcal A}\Phi|^2+s|\Phi|^2+|\Phi|^4).
\]

Writing \(\rho=|\Phi|^2\), Cauchy--Schwarz and the curvature equation
yield

\[
\int s_-^2\ge\int\rho^2
=8\int|F_{\mathcal A}^+|^2
\ge32\pi^2(a_H^+)^2\ge32\pi^2A>0.                     \tag{11}
\]

The spinor cannot vanish identically because \(p>0\). In particular,
there are no scalar-flat metrics in this family, for any \(k,\ell\).

For \(\operatorname{Ric}_g=\lambda_0g\) with \(\lambda_0\ne0\), the
definition gives \(\mathcal E(g)=0\). For a Ricci-flat metric it instead
gives \(\mathcal E(g)=4\), since \(|\lambda g|_g^2=4\lambda^2\). The
Ricci-flat case is therefore excluded by (11), not by assigning it zero
defect.

Choose \(\gamma=\varepsilon/2\). From (1) and (11),

\[
\boxed{D\ge(1/3-\gamma)T_X
\quad\Longrightarrow\quad
M\text{ admits no Einstein metric of any sign}.}        \tag{12}
\]

The sharper necessary inequality for a nonzero Einstein constant is
\(D/T_X\le1/3-\varepsilon\). The closed threshold (12) uses a smaller
constant intentionally; equality at \(1/3-\varepsilon\) is not excluded
by (1) alone.

For general-type \(X\) and \(\ell=0\), (12) contains every case of
LeBrun's Corollary 3.4. In its closed band, every metric also satisfies

\[
\mathcal E(g)\ge\varepsilon/2,\qquad
\delta_\sigma(g)\ge\sqrt{\varepsilon/8},\quad
\delta_\sigma(g)=\|g^{-1}\operatorname{Ric}_g-\sigma I\|_{\infty,\rm op},
\quad \sigma=\pm1,
\]

because \(\mathcal E(g)\le4\delta_\sigma(g)^2\).

\hypertarget{theorem-3.3-the-positive-budget-and-the-zero-corner}{%
\section{Theorem 3.3: the positive budget and the zero
corner}\label{theorem-3.3-the-positive-budget-and-the-zero-corner}}

The positive-budget connected-sum setting of Theorem 3.3 is covered by
(1) and (12), including its handle term and the specified chamber. No
new gauge theorem is required beyond its polarized input.

For nonpositive budgets, use Chern--Weil directly. If \(T_M<0\), the
identity in (9) excludes Einstein metrics. There is also an elementary
integral bound that does not require any monopole input. Set

\[
R_2=\int|\operatorname{Ric}|^2,\quad S=\int s,\quad
L=\int(s-S/V)^2,\quad W_2=\int|W^+|^2.
\]

When \(R_2>0\), optimization in \(1/\lambda\) gives

\[
\mathcal E=4-\frac{S^2}{VR_2}
=1+\frac{12W_2-24\pi^2T_M+L}{R_2}.                     \tag{13}
\]

Thus \(T_M\le0\) implies \(\mathcal E\ge1\), with the case \(R_2=0\)
handled by \(\mathcal E=4\). This is a classical algebraic consequence,
not a new calibration theorem.

If \(T_X<0\), or \(T_X=0<D\), then \(T_M=T_X-D<0\), and the exact
nonexistence conclusion follows. However, \textbf{\(T_X=D=0\) is
different}: an Einstein metric is then forced to be Ricci-flat and
anti-self-dual, not impossible. K3 with its complex orientation is an
actual example. It has nonzero Seiberg--Witten invariant and admits
Ricci-flat metrics.

The inspected author-hosted PDF and arXiv v1 of {[}L01{]} are the March
2000 preprint, not the journal typesetting. Their unrestricted wording
of Theorem 3.3 must not be imported across this zero corner. The
journal's publication metadata alone does not establish how that corner
was worded in its final text. Our precise claim is the quantitative
extension for \(T_X>0\), together with the classical \(T_M<0\)
exclusion. This qualification does not narrow Corollary 3.4: there
\(T_X=c_1^2(X)>0\).

\hypertarget{what-is-strengthened}{%
\section{What is strengthened}\label{what-is-strengthened}}

The additional threshold interval contains actual members of the
general-type family. Let

\[
X_m=\Sigma_2\times\Sigma_m,\qquad A_m=8(m-1),\qquad
k_m=\lceil A_m/3\rceil-1,\qquad m\ge2.
\]

These are minimal general-type bases and

\[
0<1/3-k_m/A_m\le1/A_m\longrightarrow0.
\]

For all sufficiently large \(m\), their blow-ups satisfy (12), although
\(k_m<A_m/3\). Their Euler characteristics tend to infinity. Thus the
strengthening is not merely an enlargement of unattained real parameter
values. This does not assert novelty over every other obstruction in the
literature, nor produce an explicit first \(m\) without an explicit
calibration constant.

The correct implication is \textbf{integral stability plus scalar-flat
exclusion implies the classical exact obstruction}. The exact
obstruction alone would not imply a uniform integral gap. LeBrun's
absolute Ricci-energy infima and minimal-volume results are different
conclusions, not consequences established by this note.

\hypertarget{appendix-algebra-behind-the-constant-transfer}{%
\section{Appendix: algebra behind the constant
transfer}\label{appendix-algebra-behind-the-constant-transfer}}

Write \(e=s+4\), \(\alpha=\sqrt{2/3}|W^+|\), \(w=\lambda_{\min}(W^+)\),
\(t=s+24\alpha\), and

\[
b=2s/3+2w,\quad
G=(2s/3-2\alpha)^2-b_-^2\ge0,\quad
P=\int|\operatorname{Ric}^{\circ}|^2.
\]

The signed spectral monopole bound is \(\int b_-^2\ge32\pi^2p\). Here
\(v_+=\max(v,0)\) and \(v_-=\max(-v,0)\) for a real number \(v\). The
pointwise identity

\[
t^2+18G=216(s^2/24+2|W^+|^2)-18b_-^2
\]

therefore gives

\[
\int(t^2+18G)\le108P+864\pi^2Ad.
\]

Also \(F\le e^2+t^2+G\). One direct verification sets
\(z=w+\alpha\in[0,\alpha]\), so

\[
F=\tfrac32(\alpha-1/6)^2+z/2,\qquad
\alpha=(t-e+4)/24.
\]

For \(s\le-3\), \(b<0\) and \(G\ge8z\), giving
\(F\le(t^2+e^2)/192+G/16\). For \(s>-3\), \(e>1\), and \(z\le\alpha\),
\(4-e\le3e^2\), \(t\le(t^2+1)/2\) give \(F\le(t^2+5e^2)/64\). Hence

\[
\int(|\operatorname{Ric}+g|^2+F)
 \le109rV+864\pi^2Ad.
\]

Finally, \(27(s^2/24+2|W^+|^2)-q^2=t^2/8\), so \(q_+^2\le q^2\) and
\(p\ge A\) yield

\[
\Delta\le(27/2)rV+108\pi^2Ad.
\]

Adding proves (8). These estimates involve only the selected class's
projection bound and its square; neither a finite class set nor a
continuous selection is hidden in the constants.

\hypertarget{references-and-proof-locations}{%
\section{References and proof
locations}\label{references-and-proof-locations}}

\begin{itemize}
\tightlist
\item
  {[}L01{]} C. LeBrun, \emph{Ricci Curvature, Minimal Volumes, and
  Seiberg--Witten Theory}, Invent. Math. \textbf{145} (2001), 279--316.
  \href{https://www.math.stonybrook.edu/~claude/min.pdf}{Author-hosted
  preprint}, Definition 1, Proposition 3.2, Theorem 3.3, Corollary 3.4;
  \href{https://doi.org/10.1007/s002220100148}{publication record}.
\item
  {[}L06{]} C. LeBrun, \emph{Four-manifolds, curvature bounds, and
  convex geometry}, \href{https://arxiv.org/abs/math/0611450}{preprint},
  Lemma 3.7.
\item
  {[}LL{]} T.-J. Li and A.-K. Liu, \emph{Uniqueness of symplectic
  canonical class, surface cone and symplectic cone of 4-manifolds with
  \(b^+=1\)},
  \href{https://arxiv.org/abs/math/0012048}{arXiv:math/0012048}, Section
  2.1.
\item
  {[}OS{]} P. Ozsváth and Z. Szabó, \emph{Higher type adjunction
  inequalities in Seiberg--Witten theory},
  \href{https://arxiv.org/abs/math/0005268}{arXiv:math/0005268},
  Proposition 2.2, Remark 2.3, Theorem 3.2.
\end{itemize}

The package's analytic proof of (7) is
\href{round4_rigidity.md}{Integral calibration}, Sections 1--6; the
exact spinor and weighted Hodge identities are in
\href{round3_monopole_input.md}{Monopole input}, Sections 6--8. The
earlier \href{round13_chamber.md}{polarized proof} and
\href{round21_chambers.md}{chamber analysis} remain intact. The present
note consolidates their interfaces; it is not a new independent main
theorem or an audit of every analytic identity in those manuscripts.

END COMPLETE SOURCE

\begin{center}\rule{0.5\linewidth}{0.5pt}\end{center}

\hypertarget{source-10-20260926_boundary_compatibility.md}{%
\subsection{Source 10:
\nolinkurl{20260926_boundary_compatibility.md}}\label{source-10-20260926_boundary_compatibility.md}}

Original path:
\nolinkurl{continuation/20260926_boundary_compatibility.md}

SHA-256:
\nolinkurl{e182615cd26b7031d4275ae9c5930e244e6088cbad39ac9770c835ba7f5bca19}

BEGIN COMPLETE SOURCE

\hypertarget{endpoint-supply-joint-class-certificates-and-an-actual-small-topology-limit}{%
\section{Endpoint supply, joint-class certificates, and an actual
small-topology
limit}\label{endpoint-supply-joint-class-certificates-and-an-actual-small-topology-limit}}

26 September 2026. This is a continuation of I08, I09 and I12, not a new
independent PDE engine. It incorporates and extends the owner's two
supplied web analyses:
\href{../guidance/source/20260926_web_eight_volume_synthesis.zh.md}{eight-volume
synthesis} and
\href{../guidance/source/20260926_web_boundary_proof_program.zh.md}{boundary
proof program}. The latter supplied the joint-matrix argument and the
scalar-slack target; the former supplied the varying-topology rounding
construction. The exact coercive PDE premises remain those of the
inherited manuscripts. No global priority or external certification of
those premises is asserted.

The increments below are: a uniform near-closedness consequence of the
actual field construction; an explicit classical-budget dominance
theorem for the proposed joint certificate; and a same-scope small-h
family with an explicit upper bound from classical gluing. The zero-Weyl
weak-gap and optimal h-order questions remain open. No theorem of
unrestricted Einstein realization is claimed here.

\hypertarget{actual-near-closedness-at-the-zero-weyl-endpoint}{%
\subsection{1. Actual near-closedness at the zero-Weyl
endpoint}\label{actual-near-closedness-at-the-zero-weyl-endpoint}}

Use tensor Ricci norms and endomorphism Weyl norms. Let an actual
Spin\^{}c structure on a closed oriented four-manifold be
monopole-solvable for every metric conformal to g, with the structure
fixed during all weight changes. Put \[
 V=\operatorname{Vol}(g),\quad p=(a_g^+)^2>0,\quad
 r=\fint|\operatorname{Ric}+g|^2,\quad
 D_\beta=\frac{\beta+1}{\beta}\int(s^2+6\beta|W^+|^2)-72\pi^2p,
 \quad \kappa=r+D_\beta/V.
\] Here beta is any finite number at least two, and t=1/beta.

\textbf{Proposition 1 (the same supplied field is nearly closed).} The
bounded smooth forms used in
\href{../mathematics/current/line1_compactified_pencil.md}{the
compactified pencil} can simultaneously be chosen with arbitrarily small
allowance such that \[
 \fint|d\omega|^2\le\frac{39}{4}\kappa+o(1).                 \tag{1.1}
\] The constant is independent of beta, volume, diameter and Sobolev
constants. This is a consequence of existing estimates, not a new
regularity theorem.

\textbf{Proof.} In
\href{../mathematics/current/line1_functional_pencil.md}{the
actual-field proof}, equations (11)--(13), the nonnegative quantity D\_0
satisfies \[
 2\int|d\psi|^2\le D_0\le4S_0/27+o(1),\qquad S_0\le117V\kappa.
\] The same smooth solution supplies omega=c\_beta psi, with
\(c_\beta^2=9\beta/[8(\beta+1)]\le9/8\). Multiplying proves (1.1). The
allowance tends to zero for the fixed metric and class; no uniform rate
of smoothing over a family is used. In particular, the beta=2
almost-parallel estimate is not being applied at other beta values. QED.

The inherited bounds, on this same form, are \[
 \|\omega\|_\infty\le3\sqrt2,\quad
 \fint\bigl(|\operatorname{Ric}+g|^2+|W^+-A_{\omega,t}|^2
                    +( |\omega|^2-2)^2\bigr)\le C\kappa+o(1),
\] \[
 A_{\omega,t}=t(|\omega|^2I/3-\omega\otimes\omega),\qquad
 \fint|\Delta_H\omega|^2\le C(1+\kappa+(-a^2)_+/V)+o(1).    \tag{1.2}
\] If a\^{}2\textgreater=T\_+=2chi+3tau, Chern--Weil implies
\((-a^2)_+/V\le r/(8\pi^2)\). For kappa\textless=1 the last budget is
thus bounded by an absolute B\_0. Also \[
 \fint|W^+|^2\le C\kappa+432t^2+o(1),                    \tag{1.3}
\] because \(|A_{\omega,t}|^2=(2/3)t^2|\omega|^4\).

\hypertarget{the-remaining-endpoint-theorem-is-not-proved}{%
\subsubsection{The remaining endpoint theorem is not
proved}\label{the-remaining-endpoint-theorem-is-not-proved}}

The new report asks whether, for a smooth self-dual two-form on a closed
oriented Riemannian four-manifold and fixed K and B\_0, the assumptions
\(\|\omega\|_\infty\le K\) and \(\fint|\Delta_H\omega|^2\le B_0\) imply
a positive universal lower bound for \[
 Z=\fint\bigl(|\operatorname{Ric}+g|^2+|W^+|^2
                     +( |\omega|^2-2)^2+|d\omega|^2\bigr). \tag{1.4}
\] The exact zero system is impossible: it would give a negative
Einstein, anti-self-dual, almost-Kaehler four-manifold. Armstrong's
local obstruction is recorded in Apostolov--Draghici,
\href{https://arxiv.org/pdf/math/0302152}{Corollaries 2--3}. Exact
impossibility is not a proof of a positive infimum in this noncompact
class of geometries.

There is also a useful non-parallelism conclusion. For a self-dual form
with Z as in (1.4), the integrated Weitzenboeck identity gives \[
 \left|\fint|\nabla\omega|^2-\frac83\right|
 \le2Z+\left(\frac43+\frac{8K^2}{3}\right)\sqrt Z.         \tag{1.5}
\] To see this, use
\(\fint|\nabla\omega|^2=2\fint|d\omega|^2+ 2\fint\langle W^+\omega,\omega\rangle-(1/3)\fint s|\omega|^2\),
the trace bound \((\fint|s+4|^2)^{1/2}\le2\sqrt Z\), and the amplitude
bound K. Thus a hypothetical zero-defect sequence here has derivative
energy tending to 8/3, not to zero. The near-parallel KE-end mechanism
is not the appropriate endpoint argument.

There is a concrete analytical distinction. The inherited bounded-form
estimates control a Hessian in L2 and the first derivative in L4.
Armstrong's prolongation argument uses the additional structure
generated by the first derivative of the almost-complex form; repeating
curvature commutation there acts on its cotangent index as well. The
original Lambda+ curvature budget cannot simply be reused as a bound for
that tensor connection. Neither a third derivative estimate nor the
corresponding opposite-Weyl contractions has been supplied. See
\href{https://johnarmstrongmaths.com/publications/resources/thesis.pdf}{Armstrong's
thesis, section 3.3.3}. This identifies an unfinished route, not a
counterexample to (1.4) and not a proof that such extra hypotheses are
necessary for every possible proof.

\textbf{Conditional consequence.} If (1.4) has a lower bound
eta(K,B\_0)\textgreater0, then, under the above actual input and
a\^{}2\textgreater=T\_+, \[
 r+D_\beta/V\ge c\frac{(\beta-2)^2}{\beta^2},\qquad
 r+P_t/V\ge ct(1-2t)^2,\quad P_t=tD_{1/t}.               \tag{1.6}
\] Indeed (1.1)--(1.3) give Z\textless=C kappa+432t\^{}2+o(1). Fix t\_0
small enough to absorb the t\^{}2 term. For t\textless t\_0 a positive
kappa gap follows. For t\textgreater=t\_0 the already proved
compactified estimate is at least a fixed multiple of (1-2t)\^{}2. The
case kappa\textgreater1 is immediate. Finally
r+tD/V\textgreater=t(r+D/V). This proves only the implication
(1.4)=\textgreater(1.6). It would improve the upper-wall endpoint margin
from cubic to linear in h-1/4, since the existing optimal parameter
satisfies t\_*(h)\textasciitilde(2/3)(h-1/4). Neither improvement is
recorded as achieved.

\hypertarget{joint-field-certificates-without-fixing-a-global-sign}{%
\subsection{2. Joint-field certificates without fixing a global
sign}\label{joint-field-certificates-without-fixing-a-global-sign}}

Assume exactly the multiple actual conformal-class inputs of
\href{../mathematics/current/line1_period_inertia.md}{I12}: classes
a\_1,\ldots,a\_m are each supplied by a fixed actual structure
throughout the conformal class, p\_i\textgreater0, and \[
 \kappa=\max_i(r+D_2(g,a_i)/V)\le1/2880.
\] Define the topological pairing matrix and the sign support function
by \[
 Q_{ij}=\frac{2\pi^2}{V}a_i\cdot a_j,\qquad
 H(Y)=\max_{\sigma\in\{\pm1\}^m}\sigma^T Y\sigma,\qquad
 A(Y)=\left(\sum_i\sqrt{Y_{ii}}\right)^2,
 \quad Y\succeq0.
\] Let CUT\_m be the convex hull of the matrices sigma sigma\^{}T, and
K\_m=CUT\_m-PSD\_m. The negative Hodge part is retained in this set.

\textbf{Proposition 2 (the supplied web certificate).} Under these
hypotheses, \[
 \operatorname{tr}(YQ)\le
       \bigl(\sqrt{H(Y)}+\sqrt{27000\kappa A(Y)}\bigr)^2.   \tag{2.1}
\] No first scalar eigenvalue, common class square, or bound on the
negative Hodge Gram matrix is assumed.

\textbf{Proof.} With normalized L2 inner product, let xi\_i=-sqrt(2) pi
H\_i, where H\_i is the self-dual harmonic representative of a\_i. I12
gives bounded u\_i with \[
 \mathbb E|u_i-\xi_i|^2\le49\kappa/2,
 \quad\mathbb E(|u_i|^2-1)^2\le180\kappa,
 \quad\mathbb E|u_i-\sigma_i u_1|^2\le6480\kappa,
\] up to arbitrarily small simultaneous allowances; sigma\_i is the
measurable sign of the inner product with u\_1, with sign(0)=+1. Choose
n=u\_1/\textbar u\_1\textbar{} away from its zero set and any measurable
unit section on that set. Then \(\mathbb E|u_1-n|^2\le180\kappa\), and
two triangle inequalities give \[
 \mathbb E|\xi_i-\sigma_i n|^2
       \le (49+26640)\kappa=26689\kappa\le27000\kappa.    \tag{2.2}
\] No measurable sign or normalized section is differentiated.

Let P be the Gram matrix of the xi\_i, C that of sigma\_i n, and E\_err
that of their differences. Thus C is in CUT\_m and Q=P-N for
N\textgreater=0. For the column operators into the L2 Hilbert space, the
Hilbert--Schmidt triangle inequality after multiplication by sqrt(Y)
gives \[
 \sqrt{\operatorname{tr}(YP)}\le\sqrt{\operatorname{tr}(YC)}
                                      +\sqrt{\operatorname{tr}(YE_{\rm err})}.
\] The terms on the right are at most sqrt(H(Y)) and sqrt(27000 kappa
A(Y)), respectively; the latter follows from (2.2) and
\(|Y_{ij}|\le\sqrt{Y_{ii}Y_{jj}}\). Since tr(YN)\textgreater=0, (2.1)
follows. Letting the allowances tend to zero is done for fixed Y; no
convergence of monopole solutions or signs is required. QED.

In particular, tr(YQ)\textgreater H(Y) implies \[
 \kappa\ge
 \frac{(\sqrt{\operatorname{tr}(YQ)}-\sqrt{H(Y)})^2}{27000A(Y)}. \tag{2.3}
\] For fixed m, if kappa tends to zero and Q has a finite limit, that
limit lies in K\_m: the column error bound gives
\(\|P-C\|_{\rm op}\le m(27000\kappa+2\sqrt{27000\kappa})+o(1)\), and
CUT\_m is compact. The conclusion is not rank-one convergence.

\hypertarget{a-precise-limitation-for-standard-signed-suppliers}{%
\subsubsection{2.1 A precise limitation for standard signed
suppliers}\label{a-precise-limitation-for-standard-signed-suppliers}}

\textbf{Proposition 3 (classical scalar budgets dominate these
certificates).} Suppose the actual selected classes have pairings of the
form \[
 Q=\sum_{j=1}^{n_+} w_j s_js_j^T-B,\qquad
 w_j\ge0,\ s_j\in\{\pm1\}^m,\ B\succeq0.                \tag{2.4}
\] Write S\_w=sum w\_j. Assume the same metric also has an actual scalar
monopole supplier with projected square at least A\_+, where
S\_w=2pi\^{}2 A\_+/V. Then, with L(r)=(1+sqrt(r)/2)\^{}2, \[
 Q/L(r)\in K_m,\qquad
 \operatorname{tr}(YQ)\le L(r)H(Y)\quad(Y\succeq0).      \tag{2.5}
\] This conclusion requires no common endpoint saturation.

\textbf{Proof.} The classical signed scalar estimate yields \[
 32\pi^2 A_+\le\int s_-^2,
 \quad\|s_-\|_{L^2(d\mu/V)}\le4+\|s+4\|_2\le4+2\sqrt r.
\] Hence S\_w\textless=L(r). A convex combination of the sign matrices
with weights w\_j/L(r), completed to total weight one with any other
sign matrix, gives C in CUT\_m with C-Q/L(r)\textgreater=0. This proves
both statements. QED.

For Y nonzero, H(Y)\textgreater0. A proposed separation
t=tr(YQ)\textgreater H(Y) already forces, by this classical estimate
alone, \[
 r\ge4\frac{(\sqrt t-\sqrt{H(Y)})^2}{H(Y)}.              \tag{2.6}
\] Since H(Y)\textless=A(Y) and r\textless=kappa, this is stronger than
(2.3) on this supplier family. Thus finding such a certificate here does
not demonstrate an application beyond the scalar projection budget.

The hypotheses apply to the standard actual signed connected-sum
suppliers with positive orthogonal blocks b\_j\^{}2\textgreater0 and a
negative-definite exceptional block. Their cup matrix has (2.4), with
w\_j=2pi\^{}2 b\_j\^{}2/V. Full independent sign supply guarantees
A\_+=sum b\_j\^{}2: average the positive Hodge squares over all signs;
the result is the sum of the individual projected squares, at least sum
b\_j\^{}2. Some sign choice therefore supplies that budget.
Alternatively, use the already established metricwise selected supplier.
This includes the relevant
\href{../mathematics/current/line1_bauer_furuta_geography.md}{BF
application} and ordinary blow-ups. A twisted version requires one
common real local system, its actual pairing and its established scalar
estimate; no cross pairing between unrelated local systems is asserted.

For a full independent sign cube with a diagonal negative exceptional
block, the statement has an exact algebraic version: Q belongs to K\_m
if and only if S\_w\textless=1. Sufficiency is the preceding convex
argument. For necessity, let Y be the orthogonal projection onto the
span of the positive sign coordinate columns. These Walsh columns have
squared norm m and are orthogonal to all exceptional columns. Hence
tr(YQ)=m S\_w, while H(Y)=m (a positive coordinate column attains the
maximum). Membership in K\_m implies m S\_w\textless=m. This necessity
is not asserted for arbitrary subsets of sign rows or arbitrary negative
matrices B.

The report's higher slow-spectrum estimate is a separate
conditional-field consequence, not a new scalar-gap-free topological
obstruction. Round18 already contains a spectral-cluster mechanism; the
later assumptions allow unequal squares and an uncontrolled negative
Hodge part. No actual family with a new separation beyond established
budgets is supplied in this note.

\hypertarget{an-actual-matrix-realization-which-does-not-increase-obstruction-scope}{%
\subsubsection{2.2 An actual matrix realization which does not increase
obstruction
scope}\label{an-actual-matrix-realization-which-does-not-increase-obstruction-scope}}

The formal three-by-three example in the supplied report can in fact be
realized by actual BF classes, but not in its required small-error
regime. Let X be a smooth degree-eight hypersurface in CP\^{}3.
Adjunction gives c\_1\^{}2=128 and p\_g=35; Lefschetz gives b\_1=0,
hence b\_+=71, congruent to three modulo four. On
\(M=X\#X\#X\#128\overline{\mathbb{CP}}^2\), the established three-block
BF supply gives the independent canonical and exceptional sign classes.
Choose positive sign rows (-,+,+), (+,-,+), (+,+,-). Choose exceptional
sign rows by repeating each of the vectors (1,1,1,1), (1,1,-1,-1),
(1,-1,1,-1) thirty-two times. The three exceptional rows are orthogonal.
The resulting actual classes have diagonal square 256 and off-diagonal
pairing -128. At volume 512pi\^{}2 their Q is exactly \[
 Q_* = \begin{pmatrix}1&-1/2&-1/2\\-1/2&1&-1/2\\-1/2&-1/2&1\end{pmatrix}.
\] Thus Y=I-11\^{}T/3 has tr(YQ\_*)=3\textgreater8/3=H(Y). However
A\_+=384 and S\_w=3/2, so (2.5) already forces
r\textgreater=10-4sqrt(6)\textgreater0.202. In particular, the common
small-kappa premise cannot hold. Also h=(128+8)/384\textgreater1/3; the
classical BF Einstein obstruction already applies. The precise supplier
and obstruction are
\href{https://arxiv.org/pdf/math/0111228}{Ishida--LeBrun, Proposition
10, Corollary 11 and Theorem D}. This example separates actual
realizability of the intersection matrix from a genuinely new
application of its certificate. It supplies the former only.

\hypertarget{small-h-in-the-actual-two-chirality-window}{%
\subsection{3. Small h in the actual two-chirality
window}\label{small-h-in-the-actual-two-chirality-window}}

Define E(g)=inf over nonzero lambda of
\(\int|\operatorname{Ric}-\lambda g|^2/(\lambda^2 V)\). For
R=int\textbar Ric\textbar\^{}2\textgreater0 and I=int s, minimization in
1/lambda gives \[
 \mathcal E=4-I^2/(VR).                                  \tag{3.1}
\] This includes I=0 as an infimum, not necessarily a minimum.

\textbf{Proposition 4 (explicit upper bound from classical gluing).} If
X is a compact complex-hyperbolic surface, A=c\_1\^{}2(X), and
\(M=X\#k\overline{\mathbb{CP}}^2\#\ell(S^1\times S^3)\), then \[
 \inf_g\mathcal E(g)\le\frac{4h}{1+h},\qquad h=(k+4\ell)/A. \tag{3.2}
\] The upper bound does not assert attainment.

\textbf{Proof.} Normalize Ric=-g on X. The scalar-flat anti-self-dual
insertions in \href{https://arxiv.org/html/math/0003068}{LeBrun, Theorem
4.3 and its proof} give, on each fixed M, metrics with \[
 V_\epsilon\to2\pi^2A,\quad I_\epsilon\to-8\pi^2A,
 \quad R_\epsilon\to8\pi^2A(1+h).                        \tag{3.3}
\] For the first two limits, the exterior converges to the normalized
background, the inserted scalar-flat regions have vanishing volume, and
the transition annuli have bounded scalar curvature and vanishing
volume. These are features of that construction, not consequences of the
Ricci-energy infimum alone. Substitution in (3.1) proves (3.2). QED.

The no-handle limit 4k/(A+k) was already recorded in the project's
\href{../mathematics/current/line1_kahler_energy.md}{Kaehler controls,
section 5}. The present use includes handles and places the bound in the
same numerical window as I08; it is not a newly invented gluing
construction.

\textbf{Theorem 5 (same-scope vanishing family, conditional on I08).}
There are actual ordinary-monopole manifolds M\_j, with h\_j tending to
zero, such that \[
 0<U_j-F(h_j)=o(h_j^2),\qquad
 c_1h_j^2\le\inf_g\mathcal E(g)\le\frac{4h_j}{1+h_j}\to0.  \tag{3.4}
\] All sufficiently late M\_j admit no Einstein metric of any sign. Here
F(h)=(4-sqrt(1-3h))\^{}2/9 and U=3min(T\_+,T\_-)/A, precisely as in I08.

\textbf{Proof.} Use connected finite unramified covers X\_j of a fixed
compact complex-hyperbolic surface, taking a subsequence with
A\_j\textgreater=j\^{}4. Their existence is explained in the retained
\href{../mathematics/current/line1_two_chirality_topology.md}{two-chirality
construction, section 3}: the infinite finitely generated lattice is
linear and residually finite. No prescribed cover degree or divisibility
is required. Put eta\_j=1/j and, eventually,
0\textless eta\_j\textless=19/300. Define \[
 x(\eta)=\frac{F(\eta)-1+3\eta}{18},\qquad
 y(\eta)=\frac{15\eta-F(\eta)+1}{72},
 \quad k_j=\lfloor A_jx(\eta_j)\rfloor,\quad
 \ell_j=\lfloor A_jy(\eta_j)\rfloor-2.                   \tag{3.5}
\] They are nonnegative for large j, since x\textasciitilde2eta/9 and
y\textasciitilde7eta/36. Let M\_j be the indicated connected sum. Write
the two fractional parts as alpha\_j,beta\_j in {[}0,1). The identities
x+4y=eta and 1+15x-12y=F(eta) imply \[
 h_j=\eta_j-(\alpha_j+4\beta_j+8)/A_j,
 \quad U_j=F(\eta_j)+(-15\alpha_j+12\beta_j+24)/A_j.       \tag{3.6}
\] For the second identity, tau(M\_j)=A\_j/9-k\_j\textgreater0
eventually, so the smaller chiral budget is T\_-; the absolute value in
U has not been omitted. Since 0\textless F'\textless2 on {[}0,19/300{]},
\[
 0<U_j-F(h_j)<(36+26)/A_j=62/A_j=o(h_j^2),\quad h_j\sim1/j.
\] Ordinary polarized supply gives, for every metric, a class with
p\textgreater=A\_j and
a\^{}2=A\_j-k\_j\textgreater=T\_+=A\_j-k\_j-4ell\_j, fixed through its
conformal weights. Handle additions do not decrease this class square;
their contribution to T\_+ is a separate topological term. I08 therefore
yields E\textgreater=c h\_j\^{}2-62/A\_j\textgreater=(c/2)h\_j\^{}2 for
large j. Apply Proposition 4 to each fixed M\_j for the upper bound.
Nonzero Einstein constants have E=0. Ricci-flat metrics are
independently excluded by \(\int s_-^2\ge32\pi^2A_j>0\). QED.

An actual diagonal sequence can approximate (3.2) on each M\_j to any
chosen o(h\_j) error. Alternatively the fixed-radius surgery
construction in the supplied synthesis gives O(h\_j) directly. Neither
gives pointwise-small Ricci error. Thus (3.4) disproves an h-independent
positive integral gap in this same class, but proves neither h\^{}2
sharpness nor a Yau26 counterexample.

\hypertarget{what-is-actually-missing-for-linear-h-order}{%
\subsection{4. What is actually missing for linear
h-order}\label{what-is-actually-missing-for-linear-h-order}}

Put S=int s\^{}2, P=int\textbar Ric\^{}circ\textbar\^{}2,
Sigma=S-32pi\^{}2 p\textgreater=0 and D=D\_2. Classical Chern--Weil
gives the exact identity \[
 36\pi^2(p-T_+)=\frac92P+\frac98\Sigma-D.                \tag{4.1}
\] The proposed strengthening Sigma\textless=C(P+D) would yield a linear
h lower bound on (3.4). It is not proved by the existing parameter
inequalities.

A precise estimate that does follow is \[
 \Sigma\le\frac49D+\frac{4\sqrt2}{3}\sqrt{RD}.            \tag{4.2}
\] Indeed let z=s+2sqrt(6)\textbar W\^{}+\textbar. The completed-square
inequality gives int z\^{}2\textless=2D, and direct substitution gives
\(\Sigma=(4/9)D+(2/3)\int sz-(1/3)\int z^2\). Cauchy--Schwarz and
S\textless=4R prove (4.2). For a Kaehler metric in its complex
orientation, \textbar W\^{}+\textbar\^{}2=s\^{}2/24 gives the stronger
identity Sigma=(4/9)D, for any p used in these definitions. The general
input is not assumed Kaehler, so this special identity cannot be
transplanted.

Here is an algebraic test isolating the missing differential
information. Write C\_A=8pi\^{}2 A and, for small positive h, prescribe
formal moments \[
 T_+=A(1-h),\quad p=A,\quad
 R=C_A(1+h+2h^2),\quad P=C_Ah^2,\quad
 S=4C_A(1+h+h^2),\quad W_+:=\int|W^+|^2=C_A(1-h)^2/6.
\] Choose constant formal scalar curvature so that I\^{}2/V=S. These
data obey the plus Chern--Weil identity, the scalar bound and all
scalar--Weyl pencil nonnegativity tests: with x=2sqrt(1+h+h\^{}2) and
y=1-h, x\textgreater=2 and x+y\textgreater=3. For U=F(h)+o(h\^{}2), the
minus Weyl energy is also positive for small h, because it equals
C\_A{[}U-1-h+2h\^{}2{]}/12. Yet \[
 D=9C_Ah^2,\quad\Sigma=4C_Ah(1+h),\quad
 \mathcal E+D/R=\frac{13h^2}{1+h+2h^2},\qquad
 \frac{\Sigma}{P+D}=\frac{2(1+h)}{5h}\to\infty.           \tag{4.3}
\] These are not actual metrics or monopole solutions, and so are not a
counterexample to the proposed inequality. They prove a narrower useful
fact: the displayed numerical identities and near-Einstein moment
conditions alone cannot prove that inequality. A proof would need
additional actual-field or differential compatibility information. The
topology in (3.4) is real; the moments in (4.3) are only a feasibility
test and are not identified with the metrics used in (3.3).

For precision, the implication of the proposed linear inequality has no
additional hidden small-deficit branch. Suppose Sigma\textless=C(P+D)
holds for a selected Spin\^{}c structure whose unperturbed equations are
solvable for every metric conformal to g, fixed through all weights,
with A\textgreater0, p\textgreater=A, a\^{}2\textgreater=T\_+=A(1-h),
and E+D/R\textless eta\_0, with uniform C and eta\_0. On the family
(3.4), eventually U\textless=1+h+3h\^{}2. If E\textless1/2, the
classical lower Ricci budget and the chiral upper budget give
1+h\textless=e\textless=U/(1-E)\textless=3. Since \[
 D/C_A=\frac92\{e-(1+h)-2(p/A-1)\},\qquad P/C_A\le eE/4,
\] one may use the loose bounds D/C\_A\textless=27h\^{}2+18E and
P/C\_A\textless=3E/4. It follows that E+D/R\textless=19E+27h\^{}2. If
the small-deficit hypothesis fails, small enough h forces
E\textgreater=eta\_0/38. If it holds, (4.1) and p-T\_+\textgreater=Ah
give h\textless=K(E+h\^{}2), hence E\textgreater=h/(2K) for sufficiently
small h. The case E\textgreater=1/2 is immediate. Together with (3.2),
the proposed inequality would therefore imply inf E asymptotic to h up
to two fixed positive multiplicative constants. Its unproved ingredient
is exactly the scalar-slack inequality, not the conversion from that
inequality to an all-metric conclusion.

\hypertarget{realization-and-the-fixed-manifold-question-remain-separate}{%
\subsection{5. Realization and the fixed-manifold question remain
separate}\label{realization-and-the-fixed-manifold-question-remain-separate}}

The supplied off-boundary identity is correct classical algebra: \[
 D_2+72\pi^2(p-T_+)+54\int|W^-|^2
      =18P+144\pi^2(\chi-3\tau).                         \tag{5.1}
\] At p\textgreater=T\_+ and chi=3tau it supplies the saturation used in
I26. It does not remove diameter or identify the filling of a long thin
end. The suggested diameter-free numerical-boundary realization and the
quantitative W\^{}\{2,p\} repair rate remain targets, not consequences
established in this continuation.

The fixed Dehn-twist bundle E\_D is already registered as C02/O02. Its
reported pointwise approximations, virtual non-Kaehler restrictions and
candidate Einstein volume window belong to the retained cross-line
proofs. This continuation neither reconstructs its full pointwise
sequence nor excludes all its negative Einstein outputs. Its global
existence question is not replaced by the small-h integral family above,
where topology varies. The existing elimination of nonnegative Einstein
outputs is a separate topological/scalar argument; no failure of one
repair algorithm supplies the remaining negative-output exclusion.

The current common problem is therefore stable: use actual fields to
decide which numerical boundaries are geometrically realizable. Sections
1--4 make specific parts of that question more precise. They do not
combine the obstruction, compact realization and unrestricted Yau26
statements into an unproved exhaustive dichotomy.

END COMPLETE SOURCE

\begin{center}\rule{0.5\linewidth}{0.5pt}\end{center}

\hypertarget{source-11-20260926_actual_topological_cost.md}{%
\subsection{Source 11:
\nolinkurl{20260926_actual_topological_cost.md}}\label{source-11-20260926_actual_topological_cost.md}}

Original path:
\nolinkurl{continuation/20260926_actual_topological_cost.md}

SHA-256:
\nolinkurl{5b189c3143977efe8ab65b4f22c06ef2e8d354d30faa69180a0f07eadf95772b}

BEGIN COMPLETE SOURCE

\hypertarget{actual-spinor-energy-and-the-location-of-a-possible-topological-cost}{%
\section{Actual spinor energy and the location of a possible topological
cost}\label{actual-spinor-energy-and-the-location-of-a-possible-topological-cost}}

26--27 September 2026, Beijing time. Continuation of
\href{20260926_boundary_compatibility.md}{the boundary proof}, based at
local commit \texttt{f48f385}. The owner's
\href{../guidance/source/20260926_web_parallel_execution.zh.md}{parallel
execution brief} proposes zero-Weyl rigidity and a linear small-h lower
bound. Neither is assumed proved here. The inherited actual
weighted-field construction is an explicit premise, with its original
fixed-structure/conformal-solvability quantifiers. No inherited proof or
eight-volume source is replaced.

Two connections are proved below using the \textbf{same bounded weighted
solution}. First, an exact remainder identity, the localized spinor
equation and signed Chern transgression force any sublinear small-h
sequence to carry both projective first-order energy and critical
transition-layer concentration. This avoids comparing unidentified
ordinary and weighted solutions. Second, a smooth cubic torsion field
satisfies an additional Dirac identity, giving a vanishing weak residual
at the hypothetical zero-Weyl endpoint. These are necessary-structure
and weak-equation results; neither the linear small-h gap nor the
zero-end exclusion is assumed or claimed.

\hypertarget{an-exact-identity-without-differentiating-the-weight}{%
\subsection{1. An exact identity without differentiating the
weight}\label{an-exact-identity-without-differentiating-the-weight}}

Let M be a closed oriented Riemannian four-manifold. Use the spinor and
curvature conventions of
\href{../mathematics/current/line1_functional_pencil.md}{the functional
pencil}. Fix an actual Spin\^{}c structure and a smooth positive weight
f, and a solution of the corresponding rescaled Seiberg--Witten
equations. Write

\[
 \rho=|\Phi|^2,\quad K_A=\int|\nabla_A\Phi|^2,\quad
 p=(a_g^+)^2,\quad \Sigma=\int s^2-32\pi^2p,
 \quad J_+=\|iF_A^+-2\pi H\|_2^2,
\]

where H is the self-dual harmonic representative of a\_g\^{}+. The sign
of the quadratic map is chosen so that the determinant curvature has
harmonic projection 2pi H. The two identities needed are

\[
 4K_A+\int s\rho+\int f\rho^2=0,\qquad
 \int f^2\rho^2=32\pi^2p+8J_+.                 \tag{1.1}
\]

\textbf{Proposition 1.} For every constant c\textgreater0, the same
solution satisfies

\[
 \boxed{\Sigma+\int(f-c)^2\rho^2
 =\|s+c\rho\|_2^2+8cK_A+8J_+.}              \tag{1.2}
\]

\textbf{Proof.} Expand the square on the right and substitute (1.1). The
coefficient of rho\^{}2 after cancelling the scalar-spinor term is
c\^{}2-2cf+f\^{}2=(f-c)\^{}2. No derivative of f is taken. QED.

For f=c=1 this recovers the ordinary-solution remainder in the supplied
brief. For a variable f, it is not legitimate to omit the additional
nonnegative term on the left. In particular, replacing Phi by sqrt(f)Phi
would introduce derivatives of f; (1.2) does not make that replacement.
The equations and Weitzenboeck input are classical; retaining and using
the same weighted solution is the point of the present composition.

\hypertarget{uniform-amplitude-and-potential-estimates-for-the-supplied-solution}{%
\subsection{2. Uniform amplitude and potential estimates for the
supplied
solution}\label{uniform-amplitude-and-potential-estimates-for-the-supplied-solution}}

Retain the actual conformal-solvability hypotheses of the inherited
pencil, with the same structure throughout all auxiliary weights. Put

\[
 r=\fint|\operatorname{Ric}+g|^2,\quad
 \kappa=r+D_\beta/V,\quad 2\le\beta<\infty,
 \quad q_\beta=4(1+1/\beta),
 \quad c_\beta=\sqrt{q_\beta},\quad \rho_\beta=2c_\beta/3.
\]

Here c\_beta denotes the reference \textbf{weight}, not the form
rescaling constant denoted similarly in some earlier manuscripts. Set
U=s+f rho and U\_beta=-4+c\_beta rho\_beta=(4/3)(2/beta-1). The selected
smooth majorants may have an arbitrarily small total allowance epsilon V
in the following estimates. Every estimate concerns the same actual
solution; one can choose epsilon after fixing each metric and class.

\textbf{Lemma 2.} The inherited construction gives, uniformly in beta,

\[
 \rho\le4,\qquad
 \fint(\rho-\rho_\beta)^2\le40\kappa+o(1),\qquad
 \fint(f-c_\beta)^2\le\tfrac{33}{2}\kappa+o(1),
                                                        \tag{2.1}
\] \[
 \fint|U-U_\beta|^2\le1504\kappa+o(1),\quad
 \fint|d\rho|^2\le\tfrac{26}{3}\kappa+o(1),\quad
 8J_+/V\le52\kappa+o(1).                              \tag{2.2}
\]

\textbf{Proof.} With the original notation w=sqrt(6)\textbar W+\textbar,
q=-s+w, Q=max(q,3), e=s+4, u=3rho/2 and f\_0=max(sqrt(Q),(-s)\_+/4), the
exact square completion gives

\[
 \int(q-q_\beta)^2\le18\kappa V,\quad
 \int e^2\le4\kappa V,\quad
 \int Q(u-\sqrt Q)^2\le117\kappa V+o(V).
\]

Clipping at three does not increase distance to q\_beta in {[}4,6{]}.
Thus int(sqrt(Q)-c\_beta)\^{}2\textless=6kappa V and
int(u-sqrt(Q))\^{}2\textless=39kappa V. Combining gives the first
amplitude estimate. Also
(f\_0-sqrt(Q))\^{}2\textless=f\_0\^{}2-Q\textless=9e\^{}2/16, which
gives the weight estimate. The potential difference is
e+c\_beta(rho-rho\_beta)+(f-c\_beta)rho. First bound the first two terms
together by 488kappa V and the last by 264kappa V, and then their sum by
1504kappa V. Finally D\_0\textless=52kappa V/3 controls
2int\textbar drho\textbar\^{}2. The actual curvature moment C\_act
satisfies C\_act-72pi\^{}2p\textless=117kappa V+o(V); since
8J\_+=(4/9)(C\_act-72pi\^{}2p), this proves the final estimate. QED.

\textbf{Corollary 3 (same-solution scalar-slack comparison at beta=2).}
Take c=sqrt(6). Then

\[
 -540\kappa V-o(V)\le8\sqrt6 K_A-\Sigma
 \le264\kappa V+o(V).                                 \tag{2.3}
\]

Indeed rho\_2=4/sqrt(6), so int(s+c rho)\^{}2\textless=488kappa V+o(V),
int(f-c)\^{}2rho\^{}2\textless=264kappa V+o(V), and
8J\_+\textless=52kappa V+o(V). Substitute in (1.2). The little-o
allowance here is freely selectable smoothing error, not an assumed
limiting solution or a rate in kappa.

\hypertarget{actual-first-order-energy-cannot-be-hidden-at-spinor-zeros}{%
\subsection{3. Actual first-order energy cannot be hidden at spinor
zeros}\label{actual-first-order-energy-cannot-be-hidden-at-spinor-zeros}}

Use a fixed pair 0\textless a\textless b\textless4/3. A Lipschitz
nonincreasing cutoff chi(rho) equals one below a and zero above b, with
\textbar chi'\textbar\textless=1/(b-a). Smooth approximations give the
same inequalities. Testing

\[
             4\nabla_A^*\nabla_A\Phi+U\Phi=0
\]

against chi(rho)Phi gives the exact real identity

\[
 4\int\chi|\nabla_A\Phi|^2+
 2\int\chi'|d\rho|^2+\int\chi U\rho=0.
\]

If mu\_b=Vol\{rho\textless b\}/V, Cauchy--Schwarz gives

\[
 \fint_{\{\rho\le a\}}|\nabla_A\Phi|^2
 \le\frac{\fint|d\rho|^2}{2(b-a)}
 +\frac b4\left(|U_\beta|\mu_b+
      \|U-U_\beta\|_{L^2_{\rm avg}}\sqrt{\mu_b}\right). \tag{3.1}
\]

The symbol fint over a subset here and below means integral over the
subset divided by \textbf{total} V. By (2.1), mu\_b\textless=
(40kappa+o(1))/(rho\_beta-b)\^{}2; rho\_beta\textgreater=4/3 and
\textbar U\_beta\textbar\textless=4/3. Consequently

\[
 \boxed{\int_{\{\rho\le a\}}|\nabla_A\Phi|^2
             \le C_{a,b}\kappa V+o(V).}                \tag{3.2}
\]

This is uniform over all finite beta\textgreater=2, and uses the actual
spinor equation and moment remainders. It is stronger than applying a
generic Hessian/L4 estimate, which would give only an O(sqrt(kappa))
bound. It does not estimate the fourth-power gradient or squared Hessian
energy on the same set.

There is a further phase/radial decomposition. On \{rho\textgreater0\},
let Pi\_Phi be Hermitian projection onto the complex line spanned by
Phi, and write gamma=Im\textless{} Phi,nabla\_A Phi \textgreater, with
the pairing linear in its second entry. For the normalized quadratic map
psi used in the original proof, direct differentiation gives

\[
 4\rho|\nabla_A\Phi|^2-|\nabla\psi|^2=4|\gamma|^2,
 \qquad
 |\Pi_\Phi\nabla_A\Phi|^2
       =\frac{|d\rho|^2}{4\rho}+\frac{|\gamma|^2}{\rho}.
                                                        \tag{3.3}
\]

These identities can be checked in a unitary frame Phi=(sqrt(rho),0),
with an arbitrary complex two-component derivative; they are gauge
invariant. Since both remainders are summands of D\_0,

\[
 \int_{\{\rho>a\}}|\Pi_\Phi\nabla_A\Phi|^2
       \le\frac{13}{2a}\kappa V+o(V).                 \tag{3.4}
\]

Thus all but O(kappa V) of K\_A lies in
int\_\{rho\textgreater a\}\textbar(1-Pi\_Phi)nabla\_A Phi\textbar\^{}2.
This is a change of direction of the spinor line, not a gauge phase
change or a change of amplitude.

For completeness, testing the spinor equation against Phi and using
(2.1) --(2.2) gives

\[
 K_A/V=-\tfrac14 U_\beta\rho_\beta+O(\sqrt\kappa)+o(1)
       =\tfrac29\sqrt{q_\beta}(1-2/\beta)
          +O(\sqrt\kappa)+o(1).                       \tag{3.5}
\]

In a putative actual zero-end sequence beta-\textgreater infinity,
kappa-\textgreater0, the projective-direction energy outside
\{rho\textless=a\} therefore tends to 4/9. The previously proved
two-form energy limit 8/3 is consistent with this. Neither limit is
zero-Weyl rigidity: it identifies which nonparallel energy that rigidity
would have to detect.

\hypertarget{a-necessary-structure-theorem-for-any-sublinear-small-h-sequence}{%
\subsection{4. A necessary structure theorem for any sublinear small-h
sequence}\label{a-necessary-structure-theorem-for-any-sublinear-small-h-sequence}}

Fix precisely the varying-topology family in Theorem 5 of
\href{20260926_boundary_compatibility.md}{the preceding continuation},
with T\_+=A(1-h), p\textgreater=A and U\_top=F(h)+o(h\^{}2). Its
selected actual class has a\^{}2=A-k and c\_2(S\^{}+){[}M{]}=ell, with
ell/(Ah)-\textgreater7/36. The positive budget, actual
conformal-solvability and chiral identities are unchanged.

\textbf{Theorem 4 (where a hypothetical sublinear cost must go).}
Suppose metrics g\_j on this same family satisfy E(g\_j)=o(h\_j). After
the optimal negative rescaling, there exist the bounded beta=2 actual
weighted solutions above, with freely chosen allowances o(h\_j), such
that

\[
 \frac{V_j}{A_j}\to2\pi^2,\qquad
 \kappa_j=o(h_j),\qquad
 \frac{\Sigma_j}{A_jh_j}\to32\pi^2,
                                                        \tag{4.1}
\] \[
 \boxed{\frac1{A_jh_j}
   \int_{\{\rho_j>a\}}|(1-\Pi_{\Phi_j})\nabla_{A_j}\Phi_j|^2
       \longrightarrow\frac{4\pi^2}{\sqrt6}.}           \tag{4.2}
\]

The integral on \{rho\_j\textless=a\} and the line/radial/phase
derivative on \{rho\_j\textgreater a\}, divided by A\_j h\_j, both tend
to zero.

\textbf{Proof, including the normalization and sign.} Put
R=int\textbar Ric\textbar\^{}2, I=int s, L\_s=int(s-I/V)\^{}2 and
e=R/(8pi\^{}2 A). For R\textgreater0,

\[
 E=4-I^2/(VR)=(4P+L_s)/R.
\]

If I\textgreater=0, the actual scalar estimate gives 32pi\^{}2
A\textless=int s\_-\^{}2\textless=L\_s\textless=RE, hence 4\textless=eE.
On the other hand, e(1-E)\textless=U\_top and U\_top\textless=2
eventually, so for E\textless1/2 we have e\textless=4 and eE\textless2,
a contradiction. Thus I\textless0 for all late j. The minimizing
constant is lambda\_\emph{=R/I\textless0. Set g'=(-lambda\_})g. Then
r(g')=E(g), V(g')=R/(4-E), and R,P,D\_2,p are invariant under this
constant rescaling. Relabel g' as g.

The exact identities are

\[
 D_2=36\pi^2A\{e-(1+h)-2(p/A-1)\},\qquad
 36\pi^2(p-T_+)=\tfrac92P+\tfrac98\Sigma-D_2.
\]

Together with p\textgreater=A, D\_2\textgreater=0 and
e(1-E)\textless=U\_top=1+h+O(h\^{}2), these give e=1+h+o(h), p-A=o(Ah),
D\_2=o(Ah), P\textless=RE/4=o(Ah), and Sigma/(Ah)-\textgreater32pi\^{}2.
Also V/A=8pi\^{}2e/(4-E)-\textgreater2pi\^{}2, proving (4.1). Equation
(2.3) then gives K\_A/(Ah)-\textgreater4pi\^{}2/sqrt(6). Choose all
smoothing allowances o(h). Equations (3.2) and (3.4), with kappa=o(h),
leave exactly (4.2). QED.

This theorem goes beyond a formal moment model: it concerns the same
actual bounded solution used in the curvature-deficit proof. It does not
assert that (4.2) is impossible. The still-open estimate
P+D\_2\textgreater=c ell would exclude the entire sequence and imply a
linear lower bound; no such estimate is assumed in Theorem 4.

\hypertarget{positive-chern-charge-forces-critical-transition-concentration}{%
\subsubsection{4.1 Positive Chern charge forces critical transition
concentration}\label{positive-chern-charge-forces-critical-transition-concentration}}

The necessary condition can be strengthened using the \textbf{sign} of
the determinant curvature, not an absolute bound on all its components.
Write Theta=(kappa+epsilon)V, where epsilon is the selectable allowance.
For bundle curvature in this subsection use the two-form
Hilbert--Schmidt norm sum\_\{i\textless j\}\textbar F\_ij\textbar\^{}2,
half the double-sum convention of the inherited Hessian lemma. Let F be
the curvature of the positive spinor connection, F=Omega+(F\_A/2)I,
where Omega is traceless.

\textbf{Theorem 5 (localized positive-charge bound).} For the same
beta=2 supplied solutions there is a universal C such that

\[
 \boxed{c_2(S^+)[M]\le C\Theta+
               C\int_{\{1/2<\rho<1\}}|\nabla_A\Phi|^4.} \tag{4.3}
\]

No bound for W-, F\_A- in L-infinity, derivatives of f, a uniform
Sobolev constant or transverse zeros is required. The inequality holds
also when c\_2 is negative; its useful application here is
c\_2=ell\textgreater0.

\textbf{Proof.} For a unitary rank-two connection, c\_2(F)=(tr(F wedge
F)-trF wedge trF)/(8pi\^{}2). In our conventions,

\[
 16\pi^2 c_2(F)=\left(|F_A^+|^2-|F_A^-|^2
                 -2|\Omega^+|^2+2|\Omega^-|^2\right)d\mu,
 \quad |\Omega^+|^2=|W^+|^2+s^2/48,\quad
 |\Omega^-|^2=|\operatorname{Ric}^\circ|^2/4.            \tag{4.4}
\]

In particular its positive part is bounded by
(\textbar F\_A+\textbar\^{}2+\textbar Ric\^{}circ\textbar\^{}2/2)/(16pi\^{}2),
while the term -\textbar F\_A-\textbar\^{}2/(16pi\^{}2) must be retained
for the argument below.

Let B=\{rho\textless=1\} and T=\{1/2\textless rho\textless1\}. Since
(4/sqrt(6)-1)\^{}2\textgreater2/5, (2.1) gives Vol(B)\textless=100Theta.
The actual equation and weight estimate give

\[
 \int_B|F_A^+|^2=\tfrac18\int_B f^2\rho^2
       \le\tfrac18(12\operatorname{Vol}(B)+33\Theta)
       \le155\Theta.                                  \tag{4.5}
\]

On B the bounded model A\_\{omega,1/2\}, with omega=sqrt(3)psi/2, has
\textbar A\_\{omega,1/2\}\textbar\^{}2=3rho\^{}4/32\textless=3/32. The
inherited squared model error is at most CTheta. Thus
int\_B\textbar W+\textbar\^{}2\textless=CTheta; also int\_B
s\^{}2\textless=CTheta and
int\textbar Ric\^{}circ\textbar\^{}2\textless=Theta. Consequently

\[
 \int_B G^2\le C\Theta,\qquad G^2=|\Omega|^2+|F_A^+|^2. \tag{4.6}
\]

On \{rho\textgreater0\}, set u=Phi/sqrt(rho), Pi=uu*, Q=1-Pi. Define a
unitary connection nabla\^{}0 by nabla\^{}0(zu+w)=dz tensor u+Q nabla\_A
w for w perpendicular to u. It fixes u, so c\_2(F\^{}0)=0. In a local
unitary frame (u,v), write nabla\_A u=i alpha u+b v; the difference
B\_c=nabla\_A-nabla\^{}0 is the matrix {[}{[}i
alpha,-b*{]},{[}b,0{]}{]}, with
\textbar B\_c\textbar\textless=sqrt(2)\textbar nabla\_A u\textbar. The
projected-connection curvature formula gives
\textbar F\^{}0\textbar\textless=C(\textbar F\textbar+\textbar nabla\_A
u\textbar\^{}2). For nabla\^{}t=nabla\^{}0+tB\_c,
F\^{}t=(1-t)F\^{}0+tF-t(1-t)B\_c wedge B\_c, hence the same bound holds.
The Chern--Weil transgression form

\[
 T_u=\frac1{4\pi^2}\int_0^1
       \{\operatorname{tr}(B_c\wedge F^t)
             -\operatorname{tr}B_c\wedge\operatorname{tr}F^t\}\,dt
\]

satisfies dT\_u=c\_2(F), and

\[
 |T_u|\le C\left(|F||\nabla_Au|+|\nabla_Au|^3\right).  \tag{4.7}
\]

Choose a smooth eta equal to zero below 1/2 and one above 1, with
\textbar eta'\textbar\textless=4, and put chi=(1-eta)\^{}2. Then
\textbar chi'\textbar\textless=8 and (chi')\^{}2/chi\textless=64 on
\{chi\textgreater0\}. The form (1-chi(rho))T\_u extends by zero smoothly
across Phi=0. Stokes' theorem therefore yields exactly

\[
 c_2(S^+)[M]=\int\chi(\rho)c_2(F)+\int d\chi(\rho)\wedge T_u.
                                                        \tag{4.8}
\]

On T, \textbar nabla\_Au\textbar\textless=\textbar nabla\_A
Phi\textbar/sqrt(rho) and
\textbar drho\textbar\textless=2sqrt(rho)\textbar nabla\_A Phi\textbar.
The second integral is bounded by
C\textbar chi'\textbar{[}(G+\textbar F\_A-\textbar)\textbar nabla\_A
Phi\textbar\^{}2+\textbar nabla\_A Phi\textbar\^{}4{]} integrated over
T. Young's inequality bounds its F\_A- part by

\[
 \frac{\chi|F_A^-|^2}{32\pi^2}
       +C\frac{(\chi')^2}{\chi}|\nabla_A\Phi|^4.
\]

The first term is absorbed by half the \textbf{negative} central
contribution in (4.4). The G term costs CG\^{}2+C\textbar nabla\_A
Phi\textbar\^{}4. Apply (4.5)--(4.6) to all remaining curvature terms.
This proves (4.3). QED.

\textbf{Corollary 6 (two distinct energy locations are compulsory).}
Under the hypothesis of Theorem 4, take epsilon\_j=o(h\_j). Then

\[
 \operatorname{Vol}(T_j)=o(A_jh_j),\qquad
 \int_{T_j}|\nabla_{A_j}\Phi_j|^2=o(A_jh_j),\qquad
 \int_{T_j}|\nabla_{A_j}\Phi_j|^4\ge cA_jh_j.            \tag{4.9}
\]

Indeed the first two statements follow from (2.1), (3.2), choosing
a=1\textless b\textless4/3. Since ell\_j/(\nolinkurl{A_jh_j})-\textgreater7/36,
(4.3) gives the third. In particular, with kappa including an
arbitrarily small allowance,

\[
 \sup_{T_j}|\nabla_{A_j}\Phi_j|^2
       \ge\frac{c h_j}{\kappa_j+\epsilon_j}\longrightarrow\infty.
                                                        \tag{4.10}
\]

Thus a sublinear defect requires both the diffuse leading L2 projective
energy from (4.2) on normal-amplitude points and critical L4
concentration in a small-amplitude transition layer. Neither a purely
slow diffuse field nor the negative determinant-curvature bulk term
alone can account for it. Here ``diffuse leading L2'' means the leading
term is outside the declared small-amplitude set; it does not assert
uniform spatial equiintegrability over every other set.

This is not yet a contradiction: the inherited global L4 budget permits
(4.9). Even ordinary equiintegrability with respect to dmu/V would give
only o(V), insufficient when the required charge is hV-\textgreater0
relative to V. Closing the linear bound needs control at the \textbf{hV
scale}, or a direct actual-equation exclusion of (4.9). No such control
is silently added.

\hypertarget{the-topological-quantities-and-the-critical-energy-not-yet-controlled}{%
\subsection{5. The topological quantities and the critical energy not
yet
controlled}\label{the-topological-quantities-and-the-critical-energy-not-yet-controlled}}

For a Spin\^{}c structure with determinant L and a=c\_1(L), the
splitting principle gives

\[
 c_2(S^+)[M]=(a^2-T_+)/4,\quad
 c_2(S^-)[M]=(a^2+2\chi-3\tau)/4,
\] \[
 c_2(S^-\otimes L)[M]
   =5\chi+6\tau+\tfrac94(a^2-T_+).                    \tag{5.1}
\]

Indeed twisting a rank-two bundle of first Chern class a by L increases
its second Chern class by 2a\^{}2. These are identities, not energy-gap
theorems. With the manuscript curvature norms, Chern--Weil also gives

\[
 \frac{8\pi^2(5\chi+6\tau)}V
 =\fint(9|W^+|^2+|W^-|^2+5s^2/24-5|\operatorname{Ric}^\circ|^2/2)
 \ge10/3-(10/3)\sqrt r-(5/2)r.                         \tag{5.2}
\]

Armstrong's compact exact argument uses a nowhere-zero torsion section
to force the corresponding characteristic number to vanish. Equations
(5.1)--(5.2) provide the desired opposite inequality for actual inputs,
but the existence of an approximate torsion section alone does not force
vanishing. Its zeros and concentrated critical energy must be retained.
See
\href{https://johnarmstrongmaths.com/publications/resources/thesis.pdf}{Armstrong,
section 2.2.2}.

The actual estimates (3.2)--(3.4) remove concentration of the
\textbf{first} spinor derivative energy as an explanation for Theorem
4's leading term. They do not prove uniform integrability of
\textbar nabla Phi\textbar\^{}4 or \textbar nabla\^{}2 Phi\textbar\^{}2,
nor any L4 bound for the derivative of a torsion section containing
nabla Phi. Those are different orders of information. The zero-end weak
gap and the linear small-h lower bound therefore remain open.

\hypertarget{a-polynomial-opposite-spinor-field-and-its-weak-dirac-equation}{%
\subsubsection{5.1 A polynomial opposite-spinor field and its weak Dirac
equation}\label{a-polynomial-opposite-spinor-field-and-its-weak-dirac-equation}}

Let B=S\^{}- tensor L, with the tensor-product connection; its
determinant connection is 3A. Define, without division by the spinor
length,

\[
 t_\Phi=\sum_i c(e^i)\Phi\otimes(\Phi\wedge\nabla^A_i\Phi)
              \in\Gamma(B).                           \tag{5.3}
\]

The contraction is invariant under orthonormal-frame and unitary
Spin\^{}c changes. It is smooth across the zero set of Phi. It is the
candidate suggested in the owner's execution brief; the following
equation is a further calculation, not an assertion supplied by the
exact rigidity theorem.

\textbf{Proposition 7 (actual weak Dirac bridge at the zero endpoint).}
For the same weighted solution, put P\_i=nabla\_i Phi, w\_i=Phi wedge
P\_i. In the Clifford convention described below,

\[
 D^-_{3A}t_\Phi=-4\sum_i P_i\otimes w_i+\mathcal K,\qquad
 |\mathcal K|\le2|W^+||\Phi|^3.                       \tag{5.4}
\]

In particular, if beta\_j-\textgreater infinity, kappa\_j-\textgreater0
and a\_j\^{}2\textgreater=T\_\{+,j\}, the actual supplied solutions,
with sufficiently small allowances, satisfy

\[
 \|t_{\Phi_j}\|_{L^4_{\rm avg}}+
 \|\nabla t_{\Phi_j}\|_{L^2_{\rm avg}}\le C,
 \quad \fint|t_{\Phi_j}|^2\longrightarrow128/81,
                                                        \tag{5.5}
\] \[
 \boxed{\|D^-_{3A_j}t_{\Phi_j}\|_{L^q_{\rm avg}}
          \longrightarrow0\quad(1\le q<2).}            \tag{5.6}
\]

This produces an actual opposite-chirality twisted spinor with
nonvanishing mean mass and a vanishing weak Dirac residual. It is not a
new monopole solution: no curvature equation for its quadratic map has
been established.

\textbf{Proof of the differential identity.} At a normal orthonormal
frame, the product rule for (5.3) splits into

\[
 Q=\sum_{ji}c_jc_i\{P_j\otimes w_i+
                       \Phi\otimes(P_j\wedge P_i)\},
 \quad H=\sum_{ji}c_jc_i\Phi\otimes(\Phi\wedge\nabla_jP_i).
\]

The two-dimensional determinant identity Phi(P\_j wedge P\_i)=P\_j
w\_i-P\_i w\_j and D\_A Phi=0 give Q=-4 sum\_i P\_i tensor w\_i, using
\nolinkurl{c_jc_i}+\nolinkurl{c_ic_j}=-2delta\_ij. The symmetric part of H is Phi tensor(Phi
wedge nabla\^{}\emph{nabla Phi), which vanishes because the same
weighted spinor equation makes nabla\^{}}nabla Phi a scalar multiple of
Phi. The remaining part is exactly

\[
 \mathcal K=\tfrac12\sum_{ji}c_jc_i\Phi\otimes
                   (\Phi\wedge R^{S^+}_{ji}\Phi).       \tag{5.7}
\]

The central determinant-curvature term vanishes in the wedge. Clifford
action on S\^{}+ annihilates the anti-self-dual first-pair component, so
the Ricci-off-diagonal block also disappears. For a real orthonormal
basis sigma\_alpha of Lambda+, set J\_alpha=c(sigma\_alpha)/sqrt(2). The
scalar curvature term disappears by sum\_alpha J\_alpha Phi tensor(Phi
wedge J\_alpha Phi)=0. With the convention
D\_A\^{}2=nabla\^{}*nabla+s/4+c(F\_A)/2, the remaining expression is

\[
 \mathcal K=-\sum_{\alpha\beta}W^+_{\alpha\beta}
                  J_\alpha\Phi\otimes(\Phi\wedge J_\beta\Phi).
\]

The matrix with columns J\_alpha Phi has operator norm
sqrt(2)\textbar Phi\textbar, and the vector of wedges has norm
sqrt(2)\textbar Phi\textbar\^{}2. This proves (5.4). Equivalently (5.7),
whose sign does not depend on identifying a curvature operator
convention, can be used throughout.

Here are explicit pointwise checks of the remaining representation
facts. On S+-\textgreater S-, take Clifford matrices

\[
 C_1=I,\quad C_2=\operatorname{diag}(i,-i),\quad
 C_3=\begin{pmatrix}0&1\\-1&0\end{pmatrix},\quad
 C_4=\begin{pmatrix}0&-i\\-i&0\end{pmatrix}.
\]

On S--\textgreater S+ use -C\_i\^{}*. Thus c(e12+e34)=-2i diag(1,-1) on
S+. At a nonzero spinor use a unitary frame Phi=sqrt(rho)u and write
nabla Phi=alpha u+beta v, nabla u=i gamma u+b v. In the notation
Omega=sqrt(2)psi/rho, differentiation gives

\[
 |\nabla\Omega|^2=8|b|^2,\quad |d\Omega|^2=8|b_L|^2,
 \quad |t_\Phi|^2=2\rho^3|b_N|^2
    =\frac{\rho^3}{4}(|\nabla\Omega|^2-|d\Omega|^2).
                                                        \tag{5.8}
\]

Here b\_N is the component in ker{[}c(.)v{]} and b\_L its
Hermitian-orthogonal complement. For closed Omega this is precisely the
non-Kaehler intrinsic torsion component, with the appropriate
conjugate-bundle convention. Thus (5.3) is not merely an arbitrary cubic
field with the same Chern class.

Let zeta=\textless Phi,nabla Phi\textgreater, with the Hermitian pairing
linear in its second entry. The Dirac equation implies
\textbar beta\_L\textbar\textless=\textbar alpha\textbar. Each complex
isotropic subspace has beta\_N dot beta\_N=beta\_L dot beta\_L=0.
Therefore

\[
 \left|\sum_iP_i\otimes w_i\right|^2
   =\rho(|\alpha\mathbin\cdot\beta|^2+|\beta\mathbin\cdot\beta|^2)
   \le5|\zeta|^2|\nabla\Phi|^2,\qquad
 \fint|\zeta|^2\le D_0/(4V).                          \tag{5.9}
\]

The last inequality follows from (3.3) and
\textbar zeta\textbar\^{}2=\textbar drho\textbar\^{}2/4+\textbar gamma\_original\textbar\^{}2.
This notation distinguishes gamma\_original=Im\textless Phi,nabla
Phi\textgreater{} from the normalized-frame gamma above. The inherited
bounded-spinor Hessian and fourth-power gradient estimates are uniform
when a\^{}2\textgreater=T+ and kappa\textless=1. Directly
differentiating (5.3) then gives the bounds in (5.5). Holder and (5.4),
(5.9) yield

\[
 \|D^-_{3A}t_\Phi\|_{L^{4/3}_{\rm avg}}
 \le2\|\Phi\|_\infty^3\|W^+\|_{L^2_{\rm avg}}
   +4\sqrt5\|\zeta\|_{L^2_{\rm avg}}
                     \|\nabla\Phi\|_{L^4_{\rm avg}}
 \le C(\sqrt\kappa+\beta^{-1})+o(1).                  \tag{5.10}
\]

The same expression is bounded in L2, since
\textbar zeta\textbar\textless=\textbar Phi\textbar\textbar nabla
Phi\textbar. Interpolation proves (5.6). For the mass limit, (5.8) also
gives

\[
 |t_\Phi|^2=\frac\rho2\left(|\nabla\psi|^2-|d\rho|^2
                -|d\psi-\rho^{-1}d\rho\wedge\psi|^2\right).
\]

Its difference from (rho/2)\textbar nabla psi\textbar\^{}2 has absolute
value bounded by
C(\textbar drho\textbar\^{}2+\textbar dpsi\textbar\^{}2), also through
spinor zeros by the polynomial identity and continuity. Its average is
O(kappa). Now rho-\textgreater4/3 in L2, nabla psi is bounded in L4, and
the previous Hodge computation gives avg\textbar nabla
psi\textbar\^{}2-\textgreater64/27. The mass limit is consequently
128/81. QED.

The remaining critical quantity is particularly explicit:

\[
 \|D^-_{3A}t_\Phi\|_{L^2_{\rm avg}}^2
 \le C\left(\fint|W^+|^2+
               \fint|\zeta|^2|\nabla_A\Phi|^2\right).   \tag{5.11}
\]

Existing L2 smallness of zeta and L4 boundedness of nabla Phi do not
make the product in (5.11) small. Nor does the mean mass limit make
t\_Phi nowhere zero. Thus (5.6) is a weak Dirac bridge, not the claimed
zero-end exclusion or a characteristic-number vanishing theorem.

\hypertarget{what-the-opposite-spinor-topological-argument-must-additionally-pay}{%
\subsubsection{5.2 What the opposite-spinor topological argument must
additionally
pay}\label{what-the-opposite-spinor-topological-argument-must-additionally-pay}}

Writing the curvature of B=S\^{}- tensor L as
F\_B=Omega\_-spin+(3/2)F\_A I, the same conventions give

\[
 16\pi^2c_2(F_B)=\bigl(2|W^-|^2+s^2/24-|\operatorname{Ric}^\circ|^2/2
                     +9|F_A^+|^2-9|F_A^-|^2\bigr)d\mu,
\] \[
 \|F_B\|_2^2=8\pi^2c_2(B)[M]+\tfrac12P+9\|F_A^-\|_2^2. \tag{5.12}
\]

At the zero endpoint, a\^{}2\textgreater=T+ bounds
avg\textbar F\_A-\textbar\^{}2 above by 2/9+o(1). It does \textbf{not}
bound avg\textbar F\_B\textbar\^{}2 universally, because c\_2(B)/V,
equivalently the opposite Weyl energy, need not be bounded above. The
uniform bound in (5.5) concerns the derivative of a particular section,
not the curvature of its whole bundle.

For any smooth section t of a Hermitian rank-two bundle, the
construction of (4.7) applied to t/\textbar t\textbar{} and a cutoff in
\textbar t\textbar/delta gives

\[
 |c_2(B)[M]|\le C\left(
    \int_{\{|t|<2\delta\}}|F_B|^2+
    \delta^{-4}\int_{\{\delta<|t|<2\delta\}}|\nabla t|^4\right).
                                                        \tag{5.13}
\]

For clarity, \textbar dchi\textbar\textless=C\textbar nabla
t\textbar/delta and
\textbar nabla(t/\textbar t\textbar)\textbar\textless=C\textbar nabla
t\textbar/delta on the transition layer; (4.7) and Young's inequality
give exactly (5.13). The cutoff is constant near zeros, so no zero-set
transversality is assumed. Formula (5.13) isolates the missing terms
rather than discarding boundary flux.

Under a\^{}2\textgreater=T+ and r-\textgreater0, (5.1)--(5.2) give
c\_2(B)/V\textgreater=5/(12pi\^{}2)+o(1). Hence decay of both terms in
(5.13), divided by V, would close the topological contradiction.
Proposition 7 by itself supplies neither: nabla t is bounded in L2, not
L4, and W- may concentrate where t is small. These facts distinguish the
actual weak Dirac bridge from an already completed weak-error rigidity
theorem.

\hypertarget{scope-correction-for-the-inherited-harmonic-zero-example}{%
\subsection{6. Scope correction for the inherited harmonic-zero
example}\label{scope-correction-for-the-inherited-harmonic-zero-example}}

The owner's new report correctly narrows its previous interpretation of
I37. In \href{20260924_harmonic_zero_scale.md}{the original proof},
lines 422--432, the scaled zero-circle construction explicitly has
divergent pointwise Ricci error and does not establish
r+q-\textgreater0. It cannot refute automatic whole-core matching from
the full near-Einstein PDE input.

In fact the particular nonzero-Ricci model fixed in that construction
gives a little more. On a fixed interior open set U of the model with
int\_U\textbar Ric(g\_tilde)\textbar\^{}2\textgreater0, its exact
rescaling g\_rho=rho\^{}2 g\_tilde satisfies

\[
 \int_{\rho U}|\operatorname{Ric}(g_\rho)+g_\rho|^2
 =\int_U|\operatorname{Ric}(\widetilde g)+\rho^2\widetilde g|^2
 \longrightarrow\int_U|\operatorname{Ric}(\widetilde g)|^2>0.
\]

Since total volume tends to the fixed exterior volume, liminf
r\_rho\textgreater0. This calculation concerns that chosen model, not
every possible harmonic zero construction. The original sufficient scale
criterion remains valid; its automatic derivation under complete PDE
input remains an open problem.

END COMPLETE SOURCE

\end{document}
